\documentclass[10pt,twoside]{book}
\usepackage[margin=0.9in,inner=1.05in,top=0.85in,bottom=0.9in]{geometry}
\usepackage{amsmath,amssymb}
\usepackage{xcolor}
\usepackage{enumitem}
\usepackage{titlesec}
\usepackage{fancyhdr}
\usepackage{microtype}
\usepackage{needspace}
\usepackage[hidelinks]{hyperref}
\definecolor{ink}{RGB}{25,32,44}
\definecolor{accent}{RGB}{0,74,132}
\definecolor{muted}{RGB}{110,120,135}
\definecolor{flag}{RGB}{150,45,40}
\setlist[itemize]{leftmargin=1.1em, itemsep=1pt, topsep=2pt, parsep=0pt}
\titleformat{\chapter}[display]{\sffamily\huge\bfseries\color{accent}}{}{0pt}{}
\titlespacing*{\chapter}{0pt}{-30pt}{18pt}
\pagestyle{fancy}\fancyhf{}
\fancyhead[LE,RO]{\sffamily\small\color{muted}\thepage}
\fancyhead[LO,RE]{\sffamily\small\color{muted}\leftmark}
\renewcommand{\headrulewidth}{0pt}
\renewcommand{\chaptermark}[1]{\markboth{#1}{}}
\newcommand{\qhead}[3]{%
  \par\vspace{14pt}\noindent
  {\sffamily\small\color{muted}#1 \textperiodcentered\ #2}\par\vspace{1pt}
  \noindent{\sffamily\large\bfseries\color{ink}\raggedright #3\par}\vspace{4pt}}
\newcommand{\lab}[1]{\par\vspace{5pt}\noindent{\sffamily\footnotesize\bfseries\color{accent}#1}\par\vspace{1pt}}
\newcommand{\flab}[1]{\par\vspace{5pt}\noindent{\sffamily\footnotesize\bfseries\color{flag}#1}\par\vspace{1pt}}
\color{ink}
\begin{document}
\frontmatter
\thispagestyle{empty}
\vspace*{2.2in}
{\sffamily\Huge\bfseries\color{accent} The Drill Deck\par}
\vspace{10pt}
{\sffamily\Large\color{ink} ML breadth and depth rounds, answered out loud\par}
\vspace{26pt}
{\normalsize\color{muted}%
Every question here is one you get asked. Every answer is what you say---%
two to three minutes of speech, not a page of derivation. The equations that
survived are the ones an interviewer actually hands you a marker for.\par}
\vspace{12pt}
{\normalsize\color{muted}Work the \textit{Say this} block until you can deliver it
without reading. The \textit{Numbers} are what you quote; the \textit{Loses the
signal} lines are what you must not do; the \textit{Then they ask} lines are where
the conversation goes next.\par}
\vfill
{\small\color{muted}Generated from the four-volume reference set. Regenerate with \texttt{make drill}.\par}
\cleardoublepage
\tableofcontents
\mainmatter
\chapter{Foundations}
\addtocontents{toc}{\protect\thispagestyle{fancy}}
\needspace{6\baselineskip}
\qhead{L5}{Conceptual}{Explain SVD and give three applications in machine learning.}
\lab{Say this}
SVD says any matrix at all --- rectangular, singular, doesn't matter --- factors into a rotation, a scaling, and another rotation. The scale factors are the singular values, sorted largest first, and truncating after the top k gives you provably the best rank-k approximation in Frobenius norm. That's Eckart--Young, and it's why SVD shows up anywhere you want to say ``this matrix is nearly low-rank.''\par\vspace{3pt}
Three places it earns its keep. LoRA: fine-tuning updates empirically live in a low-dimensional subspace, so instead of updating a d-by-d weight matrix you learn two thin factors of rank r and train 2dr parameters instead of d-squared. Be precise here, because interviewers probe it --- LoRA does not compute an SVD, it learns those factors by gradient descent. SVD is the justification, not the algorithm. Second, recommenders: the user-item matrix is approximately low-rank because users cluster into types and items into genres, and matrix factorization is learning truncated-SVD-like factors over the observed entries. Third, model compression: swap a weight matrix for its rank-k approximation and you go from m-times-n parameters to k times m-plus-n, with FLOPs falling proportionally.\par\vspace{3pt}
For choosing k, I plot the singular value spectrum and look for a gap, or take enough components to hit an explained-variance target. And PCA on centered data is just the SVD of the data matrix, which is usually the first follow-up.\par\vspace{3pt}
\lab{Numbers}
\begin{itemize}
\item LoRA: d$^2$ trainable params $\to$ 2dr for rank r
\item Rank-k compression: mn params $\to$ k(m+n)
\end{itemize}
\flab{Loses the signal}
\begin{itemize}
\item Confusing SVD with eigendecomposition --- SVD works on any matrix, eigendecomposition needs a square one.
\item Not knowing truncated SVD is the optimal rank-k approximation (Eckart--Young).
\item Saying LoRA computes an SVD of the update --- it learns the factors by gradient descent.
\end{itemize}
\lab{Then they ask}
\begin{itemize}
\item How does SVD relate to PCA?
\item How do you pick the rank k?
\item What does a full SVD cost, and when do you use randomized SVD?
\end{itemize}
\vspace{2pt}\noindent{\color{muted}\rule{\linewidth}{0.3pt}}
\needspace{6\baselineskip}
\qhead{L5}{First Principles}{Why does the chain rule matter for backpropagation? Walk through a 3-layer example.}
\lab{Say this}
Backpropagation isn't like the chain rule, it is the chain rule, applied systematically over a computational graph. Take a three-layer net: h1 is a nonlinearity of W1 x, h2 the same on W2 h1, and the output is W3 h2. The gradient of the loss with respect to W1 is a product of four factors - loss to output, output to h2, h2 to h1, h1 to W1. Each one is a Jacobian, and the backward pass evaluates that product from the loss end, reusing intermediates as it goes.\par\vspace{3pt}
Two things fall out of that. First, the gradient at any layer is a product of all the Jacobians above it, so if those factors run consistently under one you get vanishing gradients and over one you get exploding gradients. That product is the entire reason residual connections and normalization exist.\par\vspace{3pt}
Second, why reverse mode. The loss is a scalar, so propagating backward from it gets gradients for every parameter in one pass, at roughly twice the cost of a forward pass. Forward mode propagates a perturbation from the inputs, so you'd need one pass per parameter - hopeless for billions of weights. The price is memory: you keep the forward activations around to evaluate those Jacobians, which is exactly what gradient checkpointing trades back for compute.\par\vspace{3pt}
\lab{Numbers}
\begin{itemize}
\item Backward pass costs roughly 2x the forward pass
\item Gradient checkpointing: activation memory O(L) -> O(sqrt(L))
\end{itemize}
\lab{If they hand you the marker}
$\frac{\partial \mathcal{L}}{\partial W_1} = \frac{\partial \mathcal{L}}{\partial \hat{y}}\cdot\frac{\partial \hat{y}}{\partial h_2}\cdot\frac{\partial h_2}{\partial h_1}\cdot\frac{\partial h_1}{\partial W_1}$ Write the four factors out, then say the backward pass evaluates them starting from the loss factor on the left, accumulating a vector-Jacobian product at each step, which is why one pass gets every parameter. Point at the two middle Jacobians as the ones whose product vanishes or explodes.\par
\flab{Loses the signal}
\begin{itemize}
\item Describes backprop as a separate algorithm from the chain rule rather than an application of it.
\item Can't say why reverse mode beats forward mode when there are millions of parameters.
\item Forgets that the backward pass needs the stored forward activations, so misses the memory cost.
\end{itemize}
\lab{Then they ask}
\begin{itemize}
\item What does gradient checkpointing buy you, and what does it cost?
\item When would forward-mode autodiff actually be the right choice?
\item Where in that Jacobian product do vanishing gradients come from?
\end{itemize}
\vspace{2pt}\noindent{\color{muted}\rule{\linewidth}{0.3pt}}
\needspace{6\baselineskip}
\qhead{L5}{Conceptual}{What is KL divergence? Why is it not a true distance metric? When do you use it in practice?}
\lab{Say this}
KL divergence is the expected number of extra bits you pay for encoding data that really comes from p using a code built for q. It is never negative and it is zero only when the two distributions match, but it is not a distance: it is asymmetric and it does not obey the triangle inequality.\par\vspace{3pt}
The asymmetry is the part that actually matters, and I would argue it is a feature. Forward KL, with p in front, is mode-covering: it blows up wherever q puts near-zero mass on something p thinks is likely, so q is forced to spread out and cover everything. That is what you get from maximum likelihood and cross-entropy training, and from knowledge distillation, where you want the student to cover the teacher.\par\vspace{3pt}
Reverse KL, with q in front, is mode-seeking: it punishes q for putting mass where p has none, but it lets q quietly ignore whole modes. That is variational inference. The VAE's ELBO term is reverse KL against the prior, which is why a unimodal Gaussian posterior is tolerable there. It is also the shape of the KL penalty in RLHF, where you are keeping the policy inside the region the reward model was trained on.\par\vspace{3pt}
The last thing I would say is that cross-entropy equals the entropy of p plus the KL. The entropy does not depend on your parameters, so minimising cross-entropy is minimising KL. They are not different objectives.\par\vspace{3pt}
\flab{Loses the signal}
\begin{itemize}
\item Calls KL a distance, or treats KL(p||q) and KL(q||p) as interchangeable.
\item Cannot say which direction is mode-covering versus mode-seeking, or name where each one shows up.
\item Treats KL and cross-entropy as the same quantity - they differ by the entropy of the data distribution.
\end{itemize}
\lab{Then they ask}
\begin{itemize}
\item Which KL direction does the VAE's ELBO term use, and why?
\item How does mutual information relate to KL?
\item Why does mode-covering matter for distillation?
\end{itemize}
\vspace{2pt}\noindent{\color{muted}\rule{\linewidth}{0.3pt}}
\needspace{6\baselineskip}
\qhead{L6}{First Principles}{Derive the gradient of softmax cross-entropy loss with respect to the logits.}
\lab{Say this}
The gradient is predicted minus target --- the softmax probabilities minus the one-hot label. That's the whole result, and it's worth knowing cold precisely because it's so clean.\par\vspace{3pt}
Getting there is two steps. The softmax Jacobian isn't diagonal, and that's the piece people skip --- nudging one logit moves every probability, because they all share the same normalizer. That coupling is the actual content of the derivation. Then you chain it through cross-entropy, and the probability factor coming out of the Jacobian cancels the reciprocal coming out of the log derivative. What's left collapses using the fact that the one-hot target sums to one, and you land on probability minus target.\par\vspace{3pt}
Why it matters in practice: this is exactly why every framework fuses softmax and cross-entropy into a single kernel. Computed separately, you take the log of a possibly tiny softmax output and then divide by that tiny number in the backward pass, and it overflows. Fused, the backward pass is a subtraction --- no logs, no division, nothing to blow up. So call the framework's cross-entropy on raw logits rather than hand-rolling log-softmax and negative log-likelihood.\par\vspace{3pt}
And it generalizes cleanly. With label smoothing you just replace the one-hot target with the smoothed one, and the gradient is still predicted minus target --- which is a nice way to say why smoothing damps overconfidence: it never lets the gradient go all the way to zero.\par\vspace{3pt}
\lab{If they hand you the marker}
$\frac{\partial p_i}{\partial z_j} = p_i(\delta_{ij} - p_j) \quad\Longrightarrow\quad \nabla_{z}\mathcal{L} = p - y$ --- write the Jacobian first, then say out loud that the $p_i$ cancels the $1/p_i$ from the log, and the sum collapses because the target sums to one.\par
\flab{Loses the signal}
\begin{itemize}
\item Can't state the clean final form, predicted minus target.
\item Skipping the softmax Jacobian --- it is the derivation, not a detail.
\item Confusing the gradient with respect to logits with the gradient with respect to probabilities.
\end{itemize}
\lab{Then they ask}
\begin{itemize}
\item What changes with label smoothing?
\item Why is this numerically unstable if you compute softmax and cross-entropy separately?
\item What does the gradient look like for a soft (non-one-hot) target, as in distillation?
\end{itemize}
\vspace{2pt}\noindent{\color{muted}\rule{\linewidth}{0.3pt}}
\needspace{6\baselineskip}
\qhead{L5}{Conceptual}{Explain the bias-variance tradeoff. Does it apply to deep learning?}
\lab{Say this}
It applies, but the famous U-shaped curve does not. The decomposition itself is a mathematical identity --- expected squared error splits into bias squared, variance, and irreducible noise --- so it holds for a 70B transformer exactly as much as for a linear model. Bias is systematic error, how far the average prediction sits from the truth. Variance is sensitivity to which training set you happened to draw.\par\vspace{3pt}
What breaks in deep learning is the classical claim that there's a sweet spot at intermediate complexity. We train deep in the over-parameterized regime, where classical theory says variance should dominate, and instead test error comes back down past the interpolation point --- double descent. The resolution is that parameter count is a bad proxy for effective capacity: SGD has an implicit bias toward flat, simple solutions, so variance stays controlled even when the model is enormous.\par\vspace{3pt}
Practically I reduce bias by making the model bigger and control variance with more data, augmentation, weight decay, dropout and early stopping. Diagnosing which one you have is cheap --- high training error is bias, low training error with high validation error is variance.\par\vspace{3pt}
The line I'd land on is that the tradeoff is real but the optimal complexity is far larger than classical theory predicts, and in modern deep learning the binding constraint is almost always data quality and compute, not model capacity.\par\vspace{3pt}
\lab{If they hand you the marker}
$$\mathbb{E}\big[(y-\hat f(x))^2\big] = \underbrace{\mathrm{Bias}[\hat f(x)]^2}_{\text{underfit}} + \underbrace{\mathrm{Var}[\hat f(x)]}_{\text{overfit}} + \underbrace{\sigma^2}_{\text{noise}}$$ Write the three terms and name them as you go: systematic error, sensitivity to the training draw, and noise you cannot touch. Then point at the third term and say that's the floor --- no model gets below it.\par
\flab{Loses the signal}
\begin{itemize}
\item ``The bias-variance tradeoff doesn't apply to deep learning'' --- the decomposition always holds, it's the U-curve that breaks.
\item Cannot state the three terms of the decomposition when asked.
\item Never mentions double descent or implicit regularization from SGD.
\end{itemize}
\lab{Then they ask}
\begin{itemize}
\item How do you tell high bias from high variance in practice?
\item Why does weight decay reduce variance?
\item What is double descent actually measuring on the x-axis?
\end{itemize}
\vspace{2pt}\noindent{\color{muted}\rule{\linewidth}{0.3pt}}
\needspace{6\baselineskip}
\qhead{L5}{Conceptual}{Explain double descent. Why does test error decrease again in the over-parameterized regime?}
\lab{Say this}
Double descent is the observation that test error as a function of model capacity isn't the classical U. It goes down, up, and then down again, and that second descent is the regime every modern model lives in.\par\vspace{3pt}
The first part is familiar: adding parameters cuts bias faster than it raises variance, until variance takes over and error climbs. The peak sits right at the interpolation threshold, where the model has just barely enough capacity to fit the training set exactly. That's the worst place to be, because there's essentially one interpolating solution and you're forced into it however jagged it is.\par\vspace{3pt}
Past that point you keep adding capacity and error falls again. Now many parameter settings achieve zero training loss, so the question becomes which one the optimizer picks - and from small random initialization, gradient descent converges to something close to the minimum-norm interpolator, which is smoother and generalizes better. That's implicit regularization: the optimizer, not a penalty term, selects the good solution.\par\vspace{3pt}
It also shows up along the epoch axis - test error can peak partway through training and come down again well after the model has memorized the data. So ``more parameters means more overfitting'' just isn't a reliable rule.\par\vspace{3pt}
\flab{Loses the signal}
\begin{itemize}
\item ``More parameters always means more overfitting'' - that is exactly the claim double descent refutes.
\item Describes the second descent without explaining it: no implicit regularization, no minimum-norm story.
\item Treats the interpolation threshold as ordinary overfitting rather than a specific capacity point.
\end{itemize}
\lab{Then they ask}
\begin{itemize}
\item Does explicit regularization remove the peak?
\item Does this happen along the epoch axis too?
\item How does this square with scaling laws?
\end{itemize}
\vspace{2pt}\noindent{\color{muted}\rule{\linewidth}{0.3pt}}
\needspace{6\baselineskip}
\qhead{L6}{Trade-off}{Your 10B parameter model is clearly overparameterized for your task. Should you reduce it?}
\lab{Say this}
No, not during training. Over-parameterisation is doing useful work for you, and cutting the model down is the wrong lever.\par\vspace{3pt}
The reason is that the classical ``bigger model, more overfitting'' intuition does not survive in this regime. Over-parameterised models have a friendlier loss landscape with fewer bad basins for SGD to fall into, and SGD's implicit bias pushes it toward minimum-norm solutions that generalise. Double descent says the same thing from the other direction: shrinking a 10B model back toward the interpolation threshold can make test error worse, not better. And the lottery-ticket picture is that training the big model is what lets you find the good subnetwork at all.\par\vspace{3pt}
There are legitimate reasons to want a smaller model and I would name them. If inference cost is the bottleneck, that is real, but the answer is train large and deploy small: prune eighty to ninety percent of the weights, distil into a smaller student, or quantise to INT8 or INT4. If training compute is the constraint, that is also real, but the Chinchilla answer is a smaller model on more tokens, not an arbitrarily trimmed one. And if you genuinely have under a thousand examples, a 10B model will memorise them, but the fix there is more data or few-shot prompting rather than a smaller architecture.\par\vspace{3pt}
So: keep the capacity for training, compress for serving.\par\vspace{3pt}
\lab{Numbers}
\begin{itemize}
\item Post-training pruning: 80-90\% of weights removable
\end{itemize}
\flab{Loses the signal}
\begin{itemize}
\item ``Yes, shrink it so it stops overfitting'' - classical intuition that does not hold in the over-parameterised regime.
\item Ignoring inference cost, which is the one genuinely good reason to want a smaller model.
\item No awareness of post-training compression: pruning, distillation, quantisation.
\end{itemize}
\lab{Then they ask}
\begin{itemize}
\item When does over-parameterisation actually hurt?
\item Pruning or distillation - how do you choose?
\item What changes if training compute, not serving, is the binding constraint?
\end{itemize}
\vspace{2pt}\noindent{\color{muted}\rule{\linewidth}{0.3pt}}
\needspace{6\baselineskip}
\qhead{L5}{Mathematical}{Why do we scale by $\sqrt{d_k}$? Derive the variance argument.}
\lab{Say this}
Because without it the dot products grow with dimension and drive softmax into saturation.\par\vspace{3pt}
Take queries and keys with independent, unit-variance entries. The dot product is a sum of $d_k$ products, each with mean zero and variance one, so the sum has variance $d_k$ and standard deviation root $d_k$. At head dimension 64, that's a standard deviation of 8, so typical logits span roughly plus or minus 24. Softmax over a range that wide is effectively one-hot: one entry takes everything, the rest get numerically nothing.\par\vspace{3pt}
That's fatal for training, not just for the forward pass, because the gradient of softmax goes like p times one minus p. Saturate it and the gradient goes to zero --- the attention weights stop moving and the layer stops learning where to look. Dividing by root $d_k$ puts the variance back at one regardless of head dimension, which keeps softmax in the region where it has usable gradient.\par\vspace{3pt}
The cleanest way to think about it is temperature. Root $d_k$ is a temperature chosen so that a 64-dimensional head and a 128-dimensional head behave identically at initialization, which is what lets you change head dimension without retuning anything.\par\vspace{3pt}
And I'd add that it isn't a permanent fix: the argument only holds at initialization, and logits genuinely do drift upward over a long run, which is exactly why modern large models add QK-norm or soft-capping on top of the scaling.\par\vspace{3pt}
\lab{Numbers}
\begin{itemize}
\item $d_k$ = 64: raw dot-product std \textasciitilde{}= 8, so logits span roughly +/-24 before scaling
\end{itemize}
\flab{Loses the signal}
\begin{itemize}
\item ``It normalizes the values'' with no quantitative argument.
\item Can't show that the dot-product variance is $d_k$.
\item Confusing this with layer normalization.
\end{itemize}
\lab{Then they ask}
\begin{itemize}
\item What happens if you use a different scaling factor?
\item What's the relationship between this and softmax temperature?
\item Does the argument still hold after training, when entries aren't unit-variance?
\end{itemize}
\vspace{2pt}\noindent{\color{muted}\rule{\linewidth}{0.3pt}}
\chapter{Architectures}
\addtocontents{toc}{\protect\thispagestyle{fancy}}
\needspace{6\baselineskip}
\qhead{L5}{Conceptual}{Explain the dying ReLU problem. How do modern activations solve it?}
\lab{Say this}
A ReLU unit is dead when its pre-activation is negative for every example in the dataset. The output being zero isn't the problem; the problem is that the gradient is also exactly zero there, so the weights never update and nothing revives it. It's a one-way door.\par\vspace{3pt}
Neurons get there from a large update driving the bias very negative, from a bad initialization, or from a learning rate high enough to do that in one step. You diagnose it by hooking activations and logging the fraction of zeros per layer - if thirty to forty percent of units are zero across the whole dataset, not just per batch, you have it.\par\vspace{3pt}
Every modern activation fixes it the same way: guarantee a nonzero gradient on the negative side. Leaky ReLU and PReLU do it bluntly with a small negative slope. ELU does it smoothly and pulls mean activations toward zero, but you pay for the exponential. GELU multiplies the input by its Gaussian CDF, so instead of a hard on-off decision you get a soft gate, with small but nonzero gradient for negative inputs. SiLU is the same idea self-gated by a sigmoid, dipping slightly negative around minus one point three, and it's the one inside SwiGLU, which is what most current LLMs ship.\par\vspace{3pt}
\lab{Numbers}
\begin{itemize}
\item Alarm threshold: >30-40\% zero activations per layer across the dataset
\item SiLU minimum \textasciitilde{} -0.28 at x \textasciitilde{} -1.28
\end{itemize}
\flab{Loses the signal}
\begin{itemize}
\item Says the problem is that the output is zero; the problem is the zero gradient, which removes any path to recovery.
\item ``Just lower the learning rate'' - reduces the odds, doesn't remove the mechanism.
\item Thinks GELU works by never being near zero; the point is the gradient is nonzero, not the activation.
\end{itemize}
\lab{Then they ask}
\begin{itemize}
\item Can a dead ReLU unit ever come back?
\item How does He initialization help here?
\item Is activation sparsity always bad?
\end{itemize}
\vspace{2pt}\noindent{\color{muted}\rule{\linewidth}{0.3pt}}
\needspace{6\baselineskip}
\qhead{L5}{Trade-off}{Why GELU over ReLU? Why SwiGLU over GELU?}
\lab{Say this}
GELU buys you gradient flow where ReLU gives you none, and SwiGLU buys you multiplicative gating inside the FFN. Both are small wins that compound over a lot of layers.\par\vspace{3pt}
ReLU's problem in a deep Transformer is that a unit sitting on the negative side has exactly zero gradient and no path back, which is the dying ReLU problem; across hundreds of sublayers whole channels go permanently dead. There is also a hard kink at zero where the gradient jumps from zero to one. GELU weights the input by its Gaussian CDF instead, so small negative inputs get small but nonzero outputs and gradients, and the curve is smooth, which adaptive optimisers like Adam prefer. Empirically it is worth about half a point to a point on Transformer benchmarks at negligible compute cost.\par\vspace{3pt}
SwiGLU replaces the two-matrix FFN with a three-matrix gated one. You run the input through two separate projections, pass one through Swish, multiply them elementwise, then project down. The gate is the point: that elementwise product creates feature interactions a plain pointwise nonlinearity cannot. The cost is a third weight matrix, so to hold the parameter count flat you shrink the hidden dimension from four-d to eight-thirds-d. Even with fewer hidden units it wins by roughly half a point to two points across scales, which is the PaLM and LLaMA ablations, and it is why essentially every modern LLM uses it.\par\vspace{3pt}
\lab{Numbers}
\begin{itemize}
\item GELU over ReLU: \textasciitilde{}0.5-1\% on Transformer benchmarks (BERT, GPT-2, ViT)
\item SwiGLU over GELU: \textasciitilde{}0.5-2\% across scales (PaLM, LLaMA ablations)
\item SwiGLU FFN hidden dim: 8/3 d instead of 4d, to hold parameters constant
\end{itemize}
\flab{Loses the signal}
\begin{itemize}
\item ``GELU is just a smoother ReLU'' - misses the dead-neuron and gradient-flow argument.
\item ``SwiGLU is better because it is newer'' - you have to explain the gating.
\item Not mentioning the third weight matrix and the shrunken hidden dimension.
\end{itemize}
\lab{Then they ask}
\begin{itemize}
\item What else is in the GLU family, and how do the variants compare?
\item How does activation choice interact with initialisation?
\item Would you use SwiGLU in a CNN?
\end{itemize}
\vspace{2pt}\noindent{\color{muted}\rule{\linewidth}{0.3pt}}
\needspace{6\baselineskip}
\qhead{L5}{Conceptual}{Explain Flash Attention. Why is it faster even though it does the same computation?}
\lab{Say this}
Flash Attention is exact - same result as standard attention, same FLOPs. It just stops moving the attention matrix through slow memory. Attention at typical sequence lengths is memory-bound, not compute-bound, so what you optimize is I/O, not arithmetic.\par\vspace{3pt}
Standard attention materializes the full n-by-n score matrix in HBM: write it, read it back for the softmax, write again, read again for the value multiply. On an H100 that's eighty gigabytes of HBM at roughly 3.35 terabytes a second against about twenty megabytes of on-chip SRAM an order of magnitude faster. The kernel spends its life on round trips while the tensor cores sit idle.\par\vspace{3pt}
Flash Attention tiles the computation so a block of queries and keys fits in SRAM, and uses the online softmax trick - carry a running max and running sum, rescale as new tiles arrive - to get the exact softmax without ever materializing the full matrix. Memory drops from quadratic in sequence length to linear, and attention runs two to four times faster because it becomes compute-bound.\par\vspace{3pt}
If they push on versions: FA2 improved work partitioning across thread blocks and cut non-matmul FLOPs; FA3 is Hopper-native, with warp specialization, async copies and FP8.\par\vspace{3pt}
\lab{Numbers}
\begin{itemize}
\item H100: 80 GB HBM at \textasciitilde{}3.35 TB/s vs \textasciitilde{}20 MB SRAM
\item Memory O(n\textasciicircum{}2) -> O(n); attention \textasciitilde{}2-4x faster
\end{itemize}
\flab{Loses the signal}
\begin{itemize}
\item Calls it an approximation - it is bit-for-bit the same computation.
\item Says it reduces FLOPs; the FLOPs are identical, the memory traffic is what changes.
\item No mental model of SRAM versus HBM, so can't explain why identical math runs faster.
\end{itemize}
\lab{Then they ask}
\begin{itemize}
\item What is the online softmax trick doing?
\item When does Flash Attention not help?
\item What changed in FA2, and in FA3?
\end{itemize}
\vspace{2pt}\noindent{\color{muted}\rule{\linewidth}{0.3pt}}
\needspace{6\baselineskip}
\qhead{L5}{Conceptual}{How does RoPE encode position information? Why is it preferred over learned positional embeddings?}
\lab{Say this}
RoPE encodes position by rotating the query and key vectors, not by adding anything to the input. You treat each adjacent pair of dimensions as a 2D plane and rotate it by an angle proportional to the token's position, with a different frequency per pair --- low dimensions rotate fast, high dimensions slowly. The payoff is that when you dot a rotated query against a rotated key, the rotations compose, and what survives depends only on the difference between the two positions. Relative position falls out of the mechanism instead of being bolted on.\par\vspace{3pt}
Three reasons it beat learned positional embeddings. You get relative position for free, with no separate relative-attention machinery. There are no position parameters at all, so there's no maximum length baked into a table --- a learned table simply has no row for position 5000 if you trained to 4096. And the structure degrades gracefully enough that you can extend context after training.\par\vspace{3pt}
I'd be careful with that last claim, though. Naive extrapolation past the training length does fail; attention gets noisy fast. What actually works is interpolating the rotation frequencies --- position interpolation, NTK-aware scaling, YaRN --- usually with a short fine-tune. That's how a 4K model becomes a 32K or 128K model, and it's a real procedure, not something you get for nothing.\par\vspace{3pt}
\flab{Loses the signal}
\begin{itemize}
\item Saying RoPE adds a position vector to the input --- it rotates Q and K
\item Claiming it works at any length with no modification
\item Confusing it with sinusoidal absolute encodings
\end{itemize}
\lab{Then they ask}
\begin{itemize}
\item How would you take a 4K RoPE model to 128K?
\item RoPE versus ALiBi --- when does each win?
\item Why does naive extrapolation break?
\end{itemize}
\vspace{2pt}\noindent{\color{muted}\rule{\linewidth}{0.3pt}}
\needspace{6\baselineskip}
\qhead{L5}{First Principles}{Why does attention scale by $1/\sqrt{d_k}$? What happens if you don't?}
\lab{Say this}
Because the dot product grows with dimension and softmax saturates. If query and key components are roughly independent with mean zero and unit variance, the dot product has variance equal to $d_k$, so raw scores get larger as you widen the head. Dividing by the square root of $d_k$ puts that variance back at one, independent of head width.\par\vspace{3pt}
Without it you get softmax saturation. Large-magnitude scores make the distribution nearly one-hot --- one token takes essentially all the weight --- and the gradient of a saturated softmax is close to zero, so the layer stops learning. You've effectively turned soft attention into hard attention before training has had a chance to decide where to look, and the model loses the ability to blend information from several positions. What you see in practice is unstable training and degenerate attention maps.\par\vspace{3pt}
Two things worth saying to show you actually understand it rather than reciting it. First, the variance is $d_k$, not $d_k$ squared --- that's the detail people get wrong under pressure. Second, the right frame is temperature: the scale factor is a softmax temperature chosen so that score variance is one at initialization no matter how wide the head is. Once you see it as temperature, it's obvious why any other constant would either flatten attention into uniform or sharpen it into argmax.\par\vspace{3pt}
\lab{Numbers}
\begin{itemize}
\item $\mathrm{Var}(q^\top k) = d_k$ for unit-variance independent components
\end{itemize}
\lab{If they hand you the marker}
$$\mathrm{Var}(q^\top k)=\sum_{i=1}^{d_k}\mathrm{Var}(q_ik_i)=d_k \;\Rightarrow\; \mathrm{Var}\!\left(\frac{q^\top k}{\sqrt{d_k}}\right)=1$$ Write the sum, say each term is an independent product with unit variance, so the variance is the number of terms. Then divide and say that's the whole derivation --- it's a temperature that keeps score variance at one regardless of head width.\par
\flab{Loses the signal}
\begin{itemize}
\item ``It's just a normalization trick'' with no account of why it's needed.
\item Not connecting it to softmax saturation and the vanishing gradient through softmax.
\item Saying the variance is $d_k$ squared rather than $d_k$.
\end{itemize}
\lab{Then they ask}
\begin{itemize}
\item Would additive attention need the same scaling?
\item What's the connection to softmax temperature?
\item How does this interact with mixed-precision training?
\end{itemize}
\vspace{2pt}\noindent{\color{muted}\rule{\linewidth}{0.3pt}}
\needspace{6\baselineskip}
\qhead{L5}{Conceptual}{What changed between GPT-2 (2019) and LLaMA (2023) architecturally? Why each change?}
\lab{Say this}
It's the same Transformer --- every change is component-level, and each one is motivated by training stability, inference cost, or quality per FLOP.\par\vspace{3pt}
On stability, one thing people get wrong: GPT-2 was already pre-norm, that came in with GPT-2 itself. What Llama changed is LayerNorm to RMSNorm, which drops the mean subtraction, runs ten to fifteen percent faster, and costs nothing in quality. On position, learned absolute embeddings gave way to RoPE, which encodes relative position, extrapolates better, and removes the hard length ceiling. In the feed-forward block, GELU became SwiGLU --- the gate buys you a multiplicative interaction --- and because that needs three matrices instead of two, the expansion drops from four-x to about 2.67-x to hold the parameter count. Llama 2 added GQA for KV cache. Biases came out everywhere: fewer parameters, marginally faster, no quality loss.\par\vspace{3pt}
The one that has since reversed is vocabulary. Llama went from GPT-2's 50K down to 32K on the logic that a smaller embedding table is cheaper --- but Llama 3 went to 128K and Gemma to 256K, because a bigger vocab shortens sequences, which pays directly against quadratic attention and helps multilingual fertility a lot.\par\vspace{3pt}
And it kept moving after 2023, along the same two axes: QK-norm for stability in Qwen3 and Gemma 3, interleaved sliding-window and full-attention layers, MLA-style KV compression, and RoPE bases pushed to 500K for 128K-context Llama 3.1.\par\vspace{3pt}
\lab{Numbers}
\begin{itemize}
\item SwiGLU: 3 matrices instead of 2, so FFN expansion drops 4x -> \textasciitilde{}2.67x
\item Vocab: GPT-2 50K -> Llama 32K -> Llama 3 128K -> Gemma 256K
\item RoPE base scaled to 500K for 128K context (Llama 3.1)
\end{itemize}
\flab{Loses the signal}
\begin{itemize}
\item ``It got bigger'' as the whole answer.
\item ``Llama uses a different architecture'' --- it's the same Transformer with component swaps.
\item Saying GPT-2 was post-norm; GPT-2 already moved the norm to the sub-block input.
\end{itemize}
\lab{Then they ask}
\begin{itemize}
\item Why RMSNorm specifically rather than LayerNorm?
\item What's the SwiGLU expansion factor, and why isn't it 4?
\item What changed after Llama --- walk me through 2024 onward.
\end{itemize}
\vspace{2pt}\noindent{\color{muted}\rule{\linewidth}{0.3pt}}
\needspace{6\baselineskip}
\qhead{L5}{Trade-off}{CNN vs ViT for a dataset with only 10K labeled images---how do you decide, and what changes with 10M images?}
\lab{Say this}
At 10K images the real question is not CNN versus ViT, it is whether you have a good pretrained model - and you almost always do, so fine-tune a pretrained ViT.\par\vspace{3pt}
If you genuinely had to train from scratch on 10K images, the CNN wins clearly. Convolutions hand you locality, translation equivariance and weight sharing for free; a ViT has to learn all of that from data it does not have, so a ResNet-50 beats a from-scratch ViT-B comfortably. But a ViT pretrained on CLIP, DINOv2 or ImageNet-21K and then fine-tuned on your 10K will beat that ResNet, because at this data scale pretrained features dominate inductive bias.\par\vspace{3pt}
At 10M the ViT catches up and passes, because its flexibility becomes an asset - it learns spatial structure suited to your task rather than the one convolutions assume. I would be careful about quoting a crossover number, though. The original ViT paper only brackets it somewhere between ImageNet-1K, about 1.3 million images, and ImageNet-21K, about 14 million, and it is heavily recipe-dependent: DeiT-style augmentation and distillation push it down toward a million or below. ConvNeXt also shows a modernised CNN largely matching ViT at scale, so the gap is not as clean as it is usually told.\par\vspace{3pt}
Practically: CLIP ViT-B/16 fine-tuned is a strong default, EfficientNet or MobileNet if you are latency-bound on device, Swin or ConvNeXt if you need multi-scale features for detection or segmentation.\par\vspace{3pt}
\lab{Numbers}
\begin{itemize}
\item ImageNet-1K \textasciitilde{}1.3M images, ImageNet-21K \textasciitilde{}14M - the from-scratch crossover sits somewhere in between
\end{itemize}
\flab{Loses the signal}
\begin{itemize}
\item ``ViT is better, it is newer'' - data efficiency is the entire question at 10K.
\item Never considering a pretrained model, which is the practical answer.
\item Quoting a precise crossover threshold as though it were a known constant.
\end{itemize}
\lab{Then they ask}
\begin{itemize}
\item What is ConvNeXt and why does it complicate this debate?
\item How do you handle ViT position embeddings at a new resolution?
\item Would you ever train a ViT from scratch on a small dataset?
\end{itemize}
\vspace{2pt}\noindent{\color{muted}\rule{\linewidth}{0.3pt}}
\needspace{6\baselineskip}
\qhead{L5}{Conceptual}{Explain how DETR works and why it represents a paradigm shift in object detection.}
\lab{Say this}
DETR turns detection into direct set prediction: the model emits a fixed set of a hundred predictions in one shot, and training matches them one-to-one against ground truth with the Hungarian algorithm. That matching is the real innovation --- because each ground-truth object can be claimed by exactly one prediction, duplicates are penalized by construction, and anchors and non-max suppression drop out of the pipeline entirely.\par\vspace{3pt}
Mechanically: a CNN backbone produces a feature map, a transformer encoder does global self-attention over it, and a decoder takes a hundred learned object queries and cross-attends to the encoder output. Each query comes out as a class plus a box, with a no-object class for the empty slots. The queries aren't image patches --- they're learned slots that specialize by region and object scale.\par\vspace{3pt}
Why it's a shift: the whole detector is differentiable end to end, with no anchor design and no hand-tuned IoU thresholds, and it extends to panoptic segmentation by adding a mask head.\par\vspace{3pt}
The limitations are real, though. Original DETR needed 500 epochs to converge because attention starts diffuse, against roughly 12 to 36 for Faster R-CNN, and it was weak on small objects. Deformable DETR fixed most of that by attending to a few learned sample points instead of every position, cutting training to about 50 epochs, and DINO-DETR added denoising and contrastive queries to reach state of the art.\par\vspace{3pt}
\flab{Loses the signal}
\begin{itemize}
\item Not mentioning bipartite matching --- it is the innovation
\item Claiming DETR strictly beats Faster R-CNN
\item Confusing object queries with image patch tokens
\end{itemize}
\lab{Then they ask}
\begin{itemize}
\item What are object queries, mechanistically?
\item How does Deformable DETR fix the convergence problem?
\item How would you adapt this for video?
\end{itemize}
\vspace{2pt}\noindent{\color{muted}\rule{\linewidth}{0.3pt}}
\needspace{6\baselineskip}
\qhead{L5}{Conceptual}{What is classifier-free guidance and why does it improve generation quality in diffusion models?}
\lab{Say this}
Classifier-free guidance amplifies the conditioning at sampling time by extrapolating away from the unconditional prediction. During training you randomly drop the text condition, typically ten to twenty percent of the time, so the same network learns both conditional and unconditional denoising. At sampling you run it both ways and take the unconditional prediction plus w times the difference between the two. That difference is the direction the prompt is pushing; scaling it above one pushes further than the model would go on its own.\par\vspace{3pt}
Why it helps: the conditional distribution the model learned is broad, so plain conditional sampling gives diverse but often off-prompt images, and guidance sharpens toward the modes. Ho and Salimans showed the update is equivalent to guiding with an implicit classifier built from the ratio of conditional to unconditional density, which is why you don't need to train a separate one - and that's the difference from classifier guidance, which did.\par\vspace{3pt}
It's a knob, not free quality. Three to seven is balanced, seven to fifteen gives strong adherence with less diversity, Stable Diffusion defaults to 7.5, and past twenty you get oversaturation and artifacts. So CFG trades diversity for prompt adherence, which for text-to-image happens to be the trade users want.\par\vspace{3pt}
\lab{Numbers}
\begin{itemize}
\item Condition dropout during training: \textasciitilde{}10-20\%
\item Guidance scale: 3-7 balanced, 7-15 strong adherence, Stable Diffusion default 7.5
\end{itemize}
\flab{Loses the signal}
\begin{itemize}
\item Confuses it with classifier guidance, which needs a separately trained noisy-image classifier.
\item Doesn't know it requires condition dropout during training; it isn't a pure inference-time trick.
\item ``Higher guidance is always better'' - past the sweet spot you get oversaturation and lose diversity.
\end{itemize}
\lab{Then they ask}
\begin{itemize}
\item How do negative prompts fit into this?
\item What scale would you use for photorealism versus stylized work?
\item How does CFG interact with the number of sampling steps?
\end{itemize}
\vspace{2pt}\noindent{\color{muted}\rule{\linewidth}{0.3pt}}
\needspace{6\baselineskip}
\qhead{L5}{Conceptual}{How does Stable Diffusion work end-to-end? Walk through the architecture from text prompt to generated image.}
\lab{Say this}
Stable Diffusion is a latent diffusion model --- the expensive denoising happens in a compressed latent space instead of in pixels, and that single decision is what made it runnable on consumer GPUs. A pretrained VAE compresses a 512-by-512 image into a 64-by-64-by-4 latent, about 48 times fewer numbers, so each denoising step costs roughly 64 times less and the attention layers become affordable, since attention is quadratic in the number of positions.\par\vspace{3pt}
End to end: CLIP's text encoder turns the prompt into a sequence of embeddings. You sample pure Gaussian noise at the latent shape. Then a U-Net of roughly 860 million parameters iteratively predicts and removes noise. It's conditioned two ways --- the timestep enters through an embedding, and the text enters through cross-attention layers where latent positions attend to the prompt tokens. That cross-attention is the conditioning mechanism; if someone asks how the prompt actually influences the image, that's the answer.\par\vspace{3pt}
At every step you run the U-Net twice, once with the prompt and once with an empty prompt, then push the prediction away from the unconditional one. That's classifier-free guidance, typically around a scale of 7.5, and it's why generation costs two forward passes per step rather than one. After 20 to 50 steps with DDIM or a similar sampler, the VAE decoder turns the clean latent back into pixels.\par\vspace{3pt}
\lab{Numbers}
\begin{itemize}
\item 512$\times$512$\times$3 $\to$ 64$\times$64$\times$4 latent: \textasciitilde{}48$\times$ fewer values
\item U-Net \textasciitilde{}860M params; CFG scale \textasciitilde{}7.5; 20--50 sampling steps
\item CFG costs 2 U-Net passes per step
\end{itemize}
\flab{Loses the signal}
\begin{itemize}
\item Saying the diffusion runs in pixel space
\item Not naming cross-attention as the text conditioning path
\item Forgetting that classifier-free guidance doubles compute per step
\end{itemize}
\lab{Then they ask}
\begin{itemize}
\item What changed in SDXL relative to SD 1.5?
\item What is ControlNet doing architecturally?
\item How would you fine-tune this for one art style?
\end{itemize}
\vspace{2pt}\noindent{\color{muted}\rule{\linewidth}{0.3pt}}
\needspace{6\baselineskip}
\qhead{L5}{Conceptual}{Walk me through the LSTM cell gate by gate. What is each gate's purpose?}
\lab{Say this}
Three gates and a candidate, and all four read the same thing: the previous hidden state concatenated with the current input.\par\vspace{3pt}
The forget gate is a sigmoid, so it's between zero and one per dimension, and it decides what to erase from the cell state --- one means keep this memory slot, zero means wipe it. The input gate, also a sigmoid, decides how much new information to write. The cell candidate is a tanh, so it's between minus one and one, and it proposes what that new content actually is. Then the cell update is the whole point: new cell state equals old cell state times the forget gate, plus the candidate times the input gate. It's additive, not a repeated matrix multiply, so the gradient path along the cell state doesn't shrink geometrically the way a vanilla RNN's does --- that's the entire reason the architecture exists.\par\vspace{3pt}
The output gate then decides which parts of the cell state get exposed. The hidden state is the output gate times tanh of the cell state --- and note that tanh, people drop it when they write it out.\par\vspace{3pt}
The distinction I'd emphasize is cell state versus hidden state. The cell is long-term memory that can pass through many timesteps essentially unchanged; the hidden state is the filtered short-term view that the next layer and the next timestep actually see.\par\vspace{3pt}
\flab{Loses the signal}
\begin{itemize}
\item Listing the gates without saying what each one is for or how they interact.
\item Getting the order wrong --- applying the output gate before the cell update.
\item Dropping the tanh on the cell state in the hidden-state equation.
\item Not distinguishing cell state from hidden state.
\end{itemize}
\lab{Then they ask}
\begin{itemize}
\item What's the parameter count of an LSTM layer?
\item How should you initialize the forget-gate bias, and why?
\item Compare LSTM to GRU.
\end{itemize}
\vspace{2pt}\noindent{\color{muted}\rule{\linewidth}{0.3pt}}
\needspace{6\baselineskip}
\qhead{L5}{Conceptual}{Why did transformers replace LSTMs? What specific limitations did they address?}
\lab{Say this}
Three reasons, and the first is the one that actually decided it: parallelism. An LSTM's hidden state at step t depends on step t minus one, so training is a sequential chain as long as the sequence and you simply can't use the hardware. Self-attention computes all positions at once, so a training step is a few large matrix multiplies that saturate a GPU. That's orders of magnitude in wall-clock, and it's what made scaling possible at all.\par\vspace{3pt}
Second, path length. In an LSTM, information from position one reaches position T through every intermediate state, degrading along the way even with the gating. In a transformer any two positions are one attention hop apart, so a long-range dependency doesn't have to survive the journey.\par\vspace{3pt}
Third, and downstream of the first two, scaling behavior. Transformers keep improving predictably as you add parameters, data and compute - that's what the scaling laws describe - while LSTMs plateau. Once that was visible everything followed: BERT at a few hundred million, GPT-3 at 175 billion.\par\vspace{3pt}
What you give up is the constant-size state. Attention is quadratic in sequence length where an LSTM step is O(1) in memory, which is genuinely better for streaming and very long inputs, and it's why constant-state architectures keep resurfacing.\par\vspace{3pt}
\lab{Numbers}
\begin{itemize}
\item Attention is O(n\textasciicircum{}2) in sequence length; an LSTM step is O(1) state
\end{itemize}
\flab{Loses the signal}
\begin{itemize}
\item ``Transformers are just better'' with no mechanism.
\item Leads with attention quality and never mentions parallel training, which is the practical reason.
\item Claims LSTMs are entirely obsolete; constant-state models still win on streaming and very long inputs.
\end{itemize}
\lab{Then they ask}
\begin{itemize}
\item What's the computational complexity of self-attention, and what does it cost you?
\item Where would you still use an LSTM today?
\item What do transformers do worse than LSTMs?
\end{itemize}
\vspace{2pt}\noindent{\color{muted}\rule{\linewidth}{0.3pt}}
\needspace{6\baselineskip}
\qhead{L5}{Conceptual}{Explain the attention mechanism in seq2seq. Why was it needed?}
\lab{Say this}
Attention was needed because the original seq2seq model funnels the entire source sentence through one fixed-size vector, and that is an information bottleneck. The decoder has to reconstruct everything it needs about a fifty-word sentence from the encoder's last hidden state, and performance fell off as sentences got longer - exactly the symptom you would predict.\par\vspace{3pt}
Bahdanau's fix was to stop making the decoder live off one summary. At each decoding step, the decoder scores its current state against every encoder hidden state, softmaxes those scores into weights, and takes a weighted sum of the encoder states as its context. So the context is recomputed for every output token and can look at whichever part of the source is relevant right now.\par\vspace{3pt}
Three things fall out of that. The decoder gets the full encoder sequence rather than a lossy summary. The weights learn alignment - for translation you can plot them and read off a soft word alignment, which was a genuinely nice interpretability result at the time. And gradients get a much shorter path back to early encoder states, because attention is a direct connection rather than a walk back through every recurrent step.\par\vspace{3pt}
The scoring function is either additive, a small MLP over the two states, or multiplicative, a bilinear dot product. Multiplicative is cheaper, and it is the one that survives into Transformer attention.\par\vspace{3pt}
\flab{Loses the signal}
\begin{itemize}
\item Conflating seq2seq cross-attention with Transformer self-attention.
\item Never naming the fixed-vector bottleneck that motivated it.
\item ``The model pays attention to important words'' with no mechanism - no scores, no softmax, no weighted sum.
\end{itemize}
\lab{Then they ask}
\begin{itemize}
\item Additive versus multiplicative scoring - what is the difference?
\item How does Transformer cross-attention relate to this?
\item What does attention cost computationally?
\end{itemize}
\vspace{2pt}\noindent{\color{muted}\rule{\linewidth}{0.3pt}}
\needspace{6\baselineskip}
\qhead{L5}{Mathematical}{Explain self-attention step by step, including the math.}
\lab{Say this}
Every token produces three vectors from the same input through three learned projections --- a query, a key and a value. The query is what this position is looking for, the key is what a position advertises, and the value is what it hands over if it gets selected. Keeping them separate matters: it lets a position attend on one basis and retrieve something else, which a single shared projection couldn't do.\par\vspace{3pt}
You dot each query against every key, giving an n-by-n matrix of compatibility scores, then divide by the square root of the head dimension. That scaling isn't cosmetic --- the dot product of two random vectors has variance proportional to dimension, so at 128 dimensions raw scores are large, the softmax saturates toward one-hot, and gradients vanish. Scaling keeps the logits order one at initialization. Then you softmax each row, so the weights over positions are non-negative and sum to one, and use them to take a weighted average of the value vectors. Causal masking is just setting the future positions to negative infinity before the softmax.\par\vspace{3pt}
Cost is quadratic in sequence length: n-squared times d compute, and naively n-squared memory for the score matrix --- though that memory term is what FlashAttention removes by tiling and never materializing the full matrix. Multi-head is all of this done in parallel in several lower-dimensional subspaces, then concatenated.\par\vspace{3pt}
\lab{Numbers}
\begin{itemize}
\item Attention is O(n$^2$d) compute; O(n$^2$) attention-matrix memory unless tiled (FlashAttention)
\end{itemize}
\lab{If they hand you the marker}
$\text{Attention}(Q,K,V)=\mathrm{softmax}\!\left(\frac{QK^\top}{\sqrt{d_k}}\right)V$ --- write it once, then point at each piece as you talk: $QK^\top$ is an n$\times$n score matrix, one row per query; $\sqrt{d_k}$ is variance control so the softmax doesn't saturate; the output is a convex combination of the value vectors.\par
\flab{Loses the signal}
\begin{itemize}
\item Can't write the formula when handed the marker
\item Omitting the $1/\sqrt{d_k}$ scale, or unable to say why it exists
\item ``It's just a weighted average'' with no account of where the weights come from
\end{itemize}
\lab{Then they ask}
\begin{itemize}
\item How does causal masking change the computation?
\item Where does the memory saving in FlashAttention come from?
\item What is the relationship between attention and convolution?
\end{itemize}
\vspace{2pt}\noindent{\color{muted}\rule{\linewidth}{0.3pt}}
\needspace{6\baselineskip}
\qhead{L5}{Conceptual}{Why do we need multiple attention heads? What happens with just one head?}
\lab{Say this}
One head can only express one attention pattern per layer --- a single softmax distribution over positions for each query. Multiple heads let the layer attend in several representation subspaces at once: one head tracking syntactic dependencies like subject-verb, another doing something close to positional adjacency, another semantic similarity. You get several relationships per token instead of a weighted compromise between them.\par\vspace{3pt}
The cost argument is the part people get backwards. Multi-head isn't more expensive than single-head at the same model dimension, because you split the dimension across heads --- each head projects into d over h dimensions with its own query, key and value projections, so total QKV compute is unchanged. You're not buying capacity, you're buying diversity out of capacity you already had.\par\vspace{3pt}
The output projection is what makes it a single layer rather than h disjoint ones. Each head produces its own d-over-h output, they're concatenated, and $W_O$ mixes across them so downstream layers see one combined representation. Leave that out and the heads can never interact.\par\vspace{3pt}
Empirically, single-head attention at full dimension performs significantly worse, and the gap widens on tasks that need several kinds of relationship resolved at the same time.\par\vspace{3pt}
\lab{Numbers}
\begin{itemize}
\item Each head projects into d/h dimensions, so h heads cost the same as one full-width head
\end{itemize}
\flab{Loses the signal}
\begin{itemize}
\item Can't say concretely what ``different subspaces'' means.
\item Thinking multi-head attention costs more than single-head at the same model dimension.
\item Never mentioning the output projection $W_O$ and what it does.
\end{itemize}
\lab{Then they ask}
\begin{itemize}
\item How would you decide which heads matter? Can you prune heads?
\item What's the effect of head dimension on quality?
\item How does this relate to GQA and MQA on the serving side?
\end{itemize}
\vspace{2pt}\noindent{\color{muted}\rule{\linewidth}{0.3pt}}
\needspace{6\baselineskip}
\qhead{L5}{Conceptual}{Compare encoder-only, decoder-only and encoder-decoder transformers. Why did decoder-only dominate the LLM era?}
\lab{Say this}
Three shapes, and the difference is entirely in the attention mask and the training objective. Encoder-only, meaning BERT, is bidirectional: every token sees every other token, trained with masked language modeling. Great for classification and understanding, but it can't generate natively. Decoder-only, meaning GPT, is causal: each token sees only what came before it, trained to predict the next token. Encoder-decoder, meaning T5, has a bidirectional encoder feeding a causal decoder through cross-attention, which suits structured input-to-output tasks like translation.\par\vspace{3pt}
Decoder-only won for four reasons. It's one stack instead of two, so there's less to build and less to tune. Next-token prediction scales cleanly on raw internet text - no masking scheme to design, no paired data to collect. In-context learning emerged from it at scale, which meant you stopped needing a separately fine-tuned model per task. And the KV cache works cleanly under causal masking, because previously computed keys and values never change, which is what makes serving tractable.\par\vspace{3pt}
Encoder-decoder still wins on some structured tasks where you genuinely have a fixed input to condition on, and encoder-only models are still the right call for high-throughput classification and embeddings, where you don't need generation and you do want bidirectional context.\par\vspace{3pt}
\flab{Loses the signal}
\begin{itemize}
\item Can't state the attention masking difference; that is the whole distinction.
\item ``GPT is better than BERT'' with no task framing.
\item Never mentions in-context learning as the reason per-task fine-tuning stopped mattering.
\end{itemize}
\lab{Then they ask}
\begin{itemize}
\item When would you still choose BERT over an LLM?
\item What is prefix language modeling?
\item Could encoder-decoder make a comeback?
\end{itemize}
\vspace{2pt}\noindent{\color{muted}\rule{\linewidth}{0.3pt}}
\needspace{6\baselineskip}
\qhead{L5}{Mathematical}{What is the computational complexity of self-attention? Why is it a problem?}
\lab{Say this}
Self-attention is O(n squared times d) in compute and O(n squared) in memory if you materialise the score matrix - quadratic in sequence length, linear in model dimension. It comes from the Q-K-transpose product, which is n-by-d against d-by-n and gives you an n-by-n matrix, and then multiplying that by V costs the same again.\par\vspace{3pt}
The reason it is a problem is that doubling the context quadruples both compute and memory. Put a real number on it: at 128K context with 32 heads, the attention matrix alone in float16 is about a terabyte. There is no GPU where that fits, which is why long-context models cannot just be the naive implementation.\par\vspace{3pt}
There are three ways out and they differ in kind. Flash Attention is exact - it tiles the computation and never writes the n-by-n matrix out to HBM, so memory drops to linear in n while the compute stays quadratic. That is the one to name first, because it costs you nothing in quality. Then there are architectural changes that actually change the maths: sliding-window attention gets you to n times window size by attending only locally, and linear attention schemes get to order n-d by reordering the products, both at some approximation cost.\par\vspace{3pt}
So the honest framing is that Flash Attention fixed the memory wall, not the compute wall. The quadratic FLOPs are still there at long context.\par\vspace{3pt}
\lab{Numbers}
\begin{itemize}
\item 128K context, 32 heads, float16: the attention matrix alone is \textasciitilde{}1 TB
\item Attention: O(n\textasciicircum{}2 d) compute, O(n\textasciicircum{}2) memory naive, O(n) memory with Flash Attention
\end{itemize}
\flab{Loses the signal}
\begin{itemize}
\item Says linear, or cannot state O(n\textasciicircum{}2) at all.
\item Conflates compute complexity with memory complexity - Flash Attention fixes one, not the other.
\item No connection to why context length is capped in practice.
\end{itemize}
\lab{Then they ask}
\begin{itemize}
\item How does Flash Attention get to O(n) memory?
\item What do you give up with linear attention?
\item How does Mamba compare?
\end{itemize}
\vspace{2pt}\noindent{\color{muted}\rule{\linewidth}{0.3pt}}
\chapter{Training}
\addtocontents{toc}{\protect\thispagestyle{fancy}}
\needspace{6\baselineskip}
\qhead{L6}{Trade-off}{Your team has a fixed compute budget. How do you decide model size vs. training data?}
\lab{Say this}
Start from Chinchilla: for a fixed compute budget you scale parameters and tokens together, and the practical rule is about twenty tokens per parameter. Compute is roughly six times parameters times tokens, so once you fix the budget you've fixed both. By that standard most early models were badly undertrained --- GPT-3 was 175 billion parameters on 300 billion tokens, under 2:1, while Chinchilla at 70B on 1.4 trillion tokens matched GPT-3's quality in a model two and a half times cheaper to serve.\par\vspace{3pt}
But I'd push back on treating Chinchilla-optimal as the target, because it minimizes training cost only. If you're going to serve this model to millions of users, the right move is to overtrain a smaller one --- you deliberately waste training compute to buy a permanently cheaper forward pass. That's the LLaMA and Mistral approach; LLaMA-3 8B saw 15 trillion tokens, orders of magnitude past what Chinchilla would prescribe. So my first question back is: how often do you retrain versus how much do you serve?\par\vspace{3pt}
The third case is when data is the binding constraint. If you can't get more tokens, don't scale the model past what your data supports. Spend the budget on data quality instead --- filtering, deduplication, curriculum --- or on synthetic augmentation.\par\vspace{3pt}
\lab{Numbers}
\begin{itemize}
\item Chinchilla: \textasciitilde{}20 tokens per parameter
\item C $\approx$ 6ND FLOPs for a forward + backward pass
\item LLaMA-3 8B: 15T tokens, \textasciitilde{}1875 tokens per parameter
\end{itemize}
\flab{Loses the signal}
\begin{itemize}
\item Quoting the 20:1 ratio with no mention of inference cost
\item ``Just train the biggest model you can'' with no data-availability check
\item Not knowing C $\approx$ 6ND or the rough 20:1 ratio at all
\end{itemize}
\lab{Then they ask}
\begin{itemize}
\item Unlimited data but a tight inference budget --- what changes?
\item How do you estimate the compute cost up front?
\item Do scaling laws transfer to vision or code?
\end{itemize}
\vspace{2pt}\noindent{\color{muted}\rule{\linewidth}{0.3pt}}
\needspace{6\baselineskip}
\qhead{L6}{Estimation}{Estimate the compute-optimal model size for a training budget of $10^{22}$ FLOPs using Chinchilla scaling laws.}
\lab{Say this}
About 9 billion parameters on roughly 180 billion tokens.\par\vspace{3pt}
The arithmetic: compute is about six times parameters times tokens --- six FLOPs per parameter per token, because each parameter does roughly two FLOPs in the forward pass and the backward pass costs about twice the forward. Chinchilla's compute-optimal ratio is about twenty tokens per parameter. Substitute that in and C equals 120 N squared, so N is the square root of ten-to-the-22 over 120 --- square root of 8.3 times ten-to-the-19, call it 9 billion. Twenty times that is 180 billion tokens. Sanity check: six times 9e9 times 1.8e11 is 9.7e21, close enough.\par\vspace{3pt}
For scale, ten-to-the-22 FLOPs is roughly eighteen thousand A100-hours at fifty percent MFU --- a thousand GPUs for eighteen hours, tens of thousands of dollars. That's a modest run by frontier standards.\par\vspace{3pt}
The caveat I'd volunteer without being asked: Chinchilla optimizes training compute only. If you're going to serve this model at scale, inference dominates the lifetime bill, and the right move is deliberately off-optimal --- smaller model, far more tokens, the Llama recipe of a 7B on a trillion-plus. And twenty-to-one is an empirical fit, not a law; the exponents move with data quality and domain.\par\vspace{3pt}
\lab{Numbers}
\begin{itemize}
\item Chinchilla: \textasciitilde{}20 tokens per parameter
\item C \textasciitilde{}= 6ND FLOPs (\textasciitilde{}2 forward + \textasciitilde{}4 backward per parameter per token)
\item 10\textasciicircum{}22 FLOPs \textasciitilde{}= 18,000 A100-hours at 50\% MFU (\textasciitilde{}1,000 GPUs for 18 hours)
\end{itemize}
\flab{Loses the signal}
\begin{itemize}
\item Not knowing the C \textasciitilde{}= 6ND approximation, or where the 6 comes from.
\item Fumbling powers-of-ten arithmetic out loud.
\item Quoting 20:1 as a universal law rather than an empirical fit that moves with data quality.
\end{itemize}
\lab{Then they ask}
\begin{itemize}
\item Where does the factor of 6 come from?
\item How would your answer change for a code model?
\item You're going to serve this at scale --- does that change the answer?
\end{itemize}
\vspace{2pt}\noindent{\color{muted}\rule{\linewidth}{0.3pt}}
\needspace{6\baselineskip}
\qhead{L5}{Conceptual}{When do you use softmax cross-entropy vs. sigmoid BCE?}
\lab{Say this}
It comes down to one question: are the classes mutually exclusive? Softmax cross-entropy assumes exactly one is correct --- the normalizer couples all the outputs, so pushing one logit up necessarily pushes the others down. Sigmoid BCE treats every output independently, each logit through its own sigmoid, so any combination of labels can be on at once.\par\vspace{3pt}
So single-label goes to softmax: is this a cat, a dog, or a bird. Multi-label goes to sigmoid BCE: does this image contain a cat, a dog, a bird. Using softmax on a multi-label problem is a genuinely common bug, and the symptom is that co-occurring labels get suppressed, because you've forced the model to divide probability mass between labels that aren't actually competing.\par\vspace{3pt}
At very large label counts --- ten thousand and up --- softmax still works but the denominator needs every logit computed, so it gets expensive. That's what sampled softmax, noise-contrastive estimation and hierarchical softmax exist for. Sigmoid BCE scales naturally there because the outputs are independent.\par\vspace{3pt}
For hierarchical labels --- animal, then dog, then retriever --- neither handles it cleanly. I'd either use sigmoid BCE with a term that penalizes hierarchy violations, or separate softmax heads per level with consistency constraints, depending on how strict the taxonomy is.\par\vspace{3pt}
\flab{Loses the signal}
\begin{itemize}
\item ``Softmax and sigmoid are interchangeable for binary classification'' --- ignores that softmax couples the outputs.
\item ``I always use softmax because it gives probabilities'' --- sigmoid outputs are probabilities too; the distinction is independence.
\item Reaching for softmax on a multi-label problem without noticing the labels co-occur.
\end{itemize}
\lab{Then they ask}
\begin{itemize}
\item What if you have 10,000 classes --- does softmax still work?
\item How do you handle hierarchical labels?
\item Can you use sigmoid BCE for single-label classification?
\end{itemize}
\vspace{2pt}\noindent{\color{muted}\rule{\linewidth}{0.3pt}}
\needspace{6\baselineskip}
\qhead{L5}{Trade-off}{Cross-entropy vs.\textbackslash\{\} focal loss vs.\textbackslash\{\} class-balanced loss for imbalanced classification---what is your decision framework?}
\lab{Say this}
I pick based on two separate questions: is the problem class frequency, or is it easy-example dominance? Those are different failures with different fixes.\par\vspace{3pt}
If imbalance is mild, say under five to one, plain cross-entropy is usually fine. The minority class still gets enough gradient to learn. Once you're in the five to one through hundred to one range, I reach for class-weighted cross-entropy first: upweight rare classes by inverse frequency, inverse square root, or the effective-number weights if the tail is long. One hyperparameter, and it attacks the thing that's actually wrong.\par\vspace{3pt}
Focal loss solves something else. The one-minus-p-to-the-gamma factor down-weights examples the model already gets right with high confidence, so it's for the case where the majority class isn't just frequent but trivially easy. The canonical one is object detection, where 99.9 percent of anchors are easy background and the gradient signal drowns in them. If your negatives are genuinely hard, focal buys you little and costs you a gamma to tune plus some calibration damage.\par\vspace{3pt}
The point most candidates miss is that these are orthogonal. Weighting controls how much each class contributes, focal controls how much each example contributes by difficulty, and at extreme long-tail imbalance you usually want both: focal loss with effective-number alpha weights.\par\vspace{3pt}
\lab{Numbers}
\begin{itemize}
\item Object detection: \textasciitilde{}99.9\% of anchors are easy background
\item Rough bands: under 5:1 plain CE; 5:1-100:1 weighted CE; beyond 100:1 focal
\end{itemize}
\flab{Loses the signal}
\begin{itemize}
\item ``Always use focal loss for imbalanced data'' - it targets easy-example dominance, not imbalance in general.
\item Answers with resampling when the question was about the loss; related, but not what was asked.
\item Treats class weighting and focal loss as alternatives rather than composable.
\end{itemize}
\lab{Then they ask}
\begin{itemize}
\item What does focal loss do to calibration?
\item How do you choose gamma?
\item When is oversampling the better lever than changing the loss?
\end{itemize}
\vspace{2pt}\noindent{\color{muted}\rule{\linewidth}{0.3pt}}
\needspace{6\baselineskip}
\qhead{L5}{Conceptual}{Why does SimCLR need large batch sizes, and what are the alternatives?}
\lab{Say this}
Because in SimCLR the batch is the negative set. InfoNCE puts every other sample in the batch into the softmax denominator, so batch size literally is the number of negatives you contrast against. At 256 the task is too easy --- the model can separate a handful of unrelated images on colour histograms or low-level texture and drive the loss down without learning anything semantic. At 4096 or 8192 you get enough diverse, genuinely hard negatives that real semantic features are the only way to win. The original paper showed accuracy improving roughly log-linearly with batch size out to 8192.\par\vspace{3pt}
The trap is thinking gradient accumulation rescues you. It doesn't. The softmax denominator needs all the negatives in one forward pass, and you can't split it across micro-batches the way you can split a per-example loss.\par\vspace{3pt}
So if you can't afford the batch, decouple the negative count from it. MoCo keeps a queue of embeddings from earlier batches --- 65 thousand negatives at batch size 256 --- with a momentum encoder so the stale entries stay consistent enough to be usable. BYOL and SimSiam drop negatives entirely and lean on a predictor head plus stop-gradient to avoid collapse. And gradient caching splits embedding computation from loss computation, so you get a large effective batch on limited memory.\par\vspace{3pt}
\lab{Numbers}
\begin{itemize}
\item SimCLR: batch 4096--8192; accuracy roughly log-linear in batch size to 8192
\item MoCo: 65K queued negatives at batch size 256
\end{itemize}
\flab{Loses the signal}
\begin{itemize}
\item ``Use gradient accumulation to fake a large batch'' --- the InfoNCE denominator can't be split
\item ``Bigger is always better'' --- ignores diminishing returns and false negatives from same-class samples
\item Not knowing any method that decouples negatives from batch size
\end{itemize}
\lab{Then they ask}
\begin{itemize}
\item What is the memory cost of a batch that size?
\item How does MoCo's momentum encoder avoid collapse?
\item How does BYOL work with no negatives at all?
\end{itemize}
\vspace{2pt}\noindent{\color{muted}\rule{\linewidth}{0.3pt}}
\needspace{6\baselineskip}
\qhead{L5}{Trade-off}{Your classification model gets 95\% accuracy but the loss curve shows the model is still decreasing slowly. Should you keep training?}
\lab{Say this}
Not necessarily, because loss and accuracy are different objectives. Cross-entropy rewards confidence, so once you're classifying most examples correctly, further loss reduction comes almost entirely from pushing already-correct predictions from 0.95 to 0.99 to 0.999. Accuracy doesn't move; loss does. That curve on its own tells you nothing about whether more training helps.\par\vspace{3pt}
So the first thing I check is validation loss, not training loss. Training loss falling while validation loss rises is the classic divergence signal and you stop there. The second thing is calibration, which quietly degrades as you keep going --- the model starts saying 99 percent confident when it's right 95 percent of the time. That matters anywhere the probability is used rather than the argmax: ranking by confidence, ensembling, thresholded decisions, anything medical or safety-critical.\par\vspace{3pt}
The real answer is to decide what your evaluation metric actually is --- accuracy, F1, calibration error, NDCG --- monitor that on validation data, and early-stop on it rather than on loss. If you need accuracy and calibration both, train for accuracy and then post-hoc calibrate with temperature scaling on a held-out set. That's one parameter, it's cheap, and it doesn't cost you any accuracy because it's a monotone transform of the logits.\par\vspace{3pt}
\flab{Loses the signal}
\begin{itemize}
\item ``Yes, always train until loss converges'' --- ignores overfitting, calibration decay, and the proxy-metric gap.
\item Treating loss and accuracy as measuring the same thing.
\item Early-stopping on training loss rather than on the validation metric you actually care about.
\end{itemize}
\lab{Then they ask}
\begin{itemize}
\item How would you measure and improve calibration?
\item What if you need both accuracy and calibration?
\item How does label smoothing interact with this?
\end{itemize}
\vspace{2pt}\noindent{\color{muted}\rule{\linewidth}{0.3pt}}
\needspace{6\baselineskip}
\qhead{L6}{Debugging}{Your 70B-parameter run in BF16 starts loss-spiking at 100B tokens. Walk me through the activation- and normalization-level causes and mitigations.}
\lab{Say this}
A run that's been healthy for 100 billion tokens doesn't have a code bug --- bugs fire at step ten. Late-onset spikes are slow-growth instabilities: something has drifted upward all run and just crossed a threshold, usually interacting with BF16's coarse precision. BF16 has an eight-bit significand, so relative spacing is about 0.4 percent; range is fine, it shares FP32's exponent. Also separate severity: a spike that recovers is a warning, one that doesn't is divergence, and the same causes produce both.\par\vspace{3pt}
The one I'd check first is attention-logit growth: query and key norms grow all run until softmax saturates into one-hot attention. The tell is in the logs before the spike: max attention logit per head climbing for thousands of steps. If nobody logs that, adding it is step one --- the best early warning you get. Fix is QK-norm at the next restart, or tanh soft-capping mid-run, which doesn't change parameter shapes.\par\vspace{3pt}
Second, output-logit divergence: log Z of the final softmax drifts off zero and cross-entropy turns spike-prone, which PaLM's z-loss fixes for essentially nothing. Third, embedding norms grow for frequent tokens and feed straight into logit scale, so weight-decay the embeddings. And check the dull stuff: norms and softmax in FP32 even in a BF16 run, epsilon around 1e-6, not 1e-12.\par\vspace{3pt}
Then rule out the data: rewind to the checkpoint before the spike and replay the same batches. If it reproduces, skip a few hundred batches, PaLM-style; if not, suspect a flaky node.\par\vspace{3pt}
\lab{Numbers}
\begin{itemize}
\item BF16: 8-bit significand, relative spacing \textasciitilde{}0.4\%; same exponent range as FP32
\item PaLM z-loss: \textasciitilde{}1e-4 * log\textasciicircum{}2 Z
\item RMSNorm epsilon safe band: 1e-6 to 1e-5
\end{itemize}
\flab{Loses the signal}
\begin{itemize}
\item ``Lower the learning rate'' --- trades compute efficiency for nothing without finding what's growing, and the spike usually comes back.
\item Reaching for FP16-era fixes like loss scaling; BF16's problem is precision granularity, not range.
\item Running the single-GPU NaN checklist --- an init or epsilon bug would have failed at step 100, not at 100B tokens.
\end{itemize}
\lab{Then they ask}
\begin{itemize}
\item Why do loss spikes get more frequent as models scale?
\item The spike recovered on its own --- do you intervene anyway?
\item How would muP change this failure surface?
\end{itemize}
\vspace{2pt}\noindent{\color{muted}\rule{\linewidth}{0.3pt}}
\needspace{6\baselineskip}
\qhead{L6}{Estimation}{Estimate the memory required to train a 7B parameter model with batch size 1, sequence length 2048.}
\lab{Say this}
Call it 130 to 140 gigabytes, and I would get there in four buckets.\par\vspace{3pt}
Parameters in BF16 are 7 billion times two bytes, so 14 gigs. Gradients, also BF16, another 14. Then AdamW, which is the part people forget: FP32 master weights at 28 gigs, plus momentum at 28 and variance at 28. That is 84 gigabytes of optimiser state, six times the size of the model, and you are at 112 before doing any real computation.\par\vspace{3pt}
Activations are the fourth bucket and they depend on shapes. For a 32-layer model with hidden size 4096 at sequence length 2048 and batch size one, it works out to roughly a gigabyte per layer, so about 30 gigs total, and the dominant piece of that is the attention score matrices. Flash Attention never materialises those, which brings it down to nine or ten gigs, and gradient checkpointing cuts further at the cost of recompute.\par\vspace{3pt}
So 130 to 140 all in, and the practical consequence is that this does not fit on one 80-gigabyte A100 or H100 - the H100 is faster, not bigger. You need ZeRO-3 to shard optimiser state, gradients and parameters across GPUs, or at minimum ZeRO-2 plus checkpointing on two cards. Even an H200 at 141 gigabytes does not leave comfortable headroom.\par\vspace{3pt}
\lab{Numbers}
\begin{itemize}
\item AdamW mixed precision: 12 bytes/param of optimiser state = 84 GB at 7B
\item Activations at s=2048, b=1, 32 layers: \textasciitilde{}30 GB, \textasciitilde{}9-10 GB with Flash Attention
\item A100 and H100 are 80 GB; H200 is 141 GB
\end{itemize}
\flab{Loses the signal}
\begin{itemize}
\item ``7B times 4 bytes is 28 GB, done'' - forgets gradients, optimiser state and activations.
\item Not knowing Adam keeps two extra full-size copies per parameter on top of the master weights.
\item ``Just use FP16'' without noting the optimiser still holds FP32 state.
\end{itemize}
\lab{Then they ask}
\begin{itemize}
\item What changes at 32K context instead of 2K?
\item How many A100s or H100s do you actually need to train this?
\item What is the difference between ZeRO-1, ZeRO-2 and ZeRO-3?
\end{itemize}
\vspace{2pt}\noindent{\color{muted}\rule{\linewidth}{0.3pt}}
\needspace{6\baselineskip}
\qhead{L5}{Trade-off}{SimCLR vs BYOL vs MAE---give me a decision framework for choosing between them.}
\lab{Say this}
The decision comes down to your backbone and your batch budget, and I would answer it in that order.\par\vspace{3pt}
If you are on a CNN, it is SimCLR or BYOL. SimCLR is contrastive, so negatives come from the batch, which means batches in the four-to-eight-thousand range; it degrades sharply below that, and the augmentation recipe matters a lot - crop, colour jitter, blur. If you cannot afford those batch sizes, MoCo gives you the same idea with a momentum queue, or you go to BYOL, which drops negatives entirely and trains fine at 256 to 512. BYOL's risk is collapse if the momentum coefficient or the predictor is wrong.\par\vspace{3pt}
If you are on a ViT and the downstream task is classification, MAE. It only encodes twenty-five percent of the patches, so it is by far the best compute-to-quality ratio, and it barely needs augmentation. Two caveats: it wants a lot of epochs - the paper runs 1600 - and its features want full fine-tuning rather than linear probing.\par\vspace{3pt}
If you are on a ViT and the task is dense prediction, detection or segmentation, I would reach for DINOv2 instead, because MAE's local features are weaker there.\par\vspace{3pt}
So the short version: CNN means SimCLR, MoCo or BYOL; ViT plus classification means MAE; ViT plus dense prediction means DINOv2; and a small batch budget pushes you to BYOL or MoCo either way.\par\vspace{3pt}
\lab{Numbers}
\begin{itemize}
\item SimCLR batch size: 4096-8192
\item BYOL works at 256-512
\item MAE encodes 25\% of patches (75\% masking); 1600 pretraining epochs
\end{itemize}
\flab{Loses the signal}
\begin{itemize}
\item ``MAE, it is the newest'' - the downstream task decides.
\item Not knowing SimCLR's dependence on very large batches.
\item Getting wrong which of the three need negative samples.
\end{itemize}
\lab{Then they ask}
\begin{itemize}
\item What if you only have four GPUs?
\item Why is MAE weaker under linear probing?
\item What does DINOv2 take from each family?
\end{itemize}
\vspace{2pt}\noindent{\color{muted}\rule{\linewidth}{0.3pt}}
\needspace{6\baselineskip}
\qhead{L5}{Conceptual}{Why do non-contrastive methods like BYOL work without negative samples? Shouldn't they collapse to a trivial solution?}
\lab{Say this}
It should collapse --- mapping every input to the same constant vector is a global minimum of that loss --- and the reason it doesn't is the predictor plus the stop-gradient, not negatives sneaking in somewhere.\par\vspace{3pt}
The essential mechanism is architectural asymmetry. The online branch has a predictor head that the target branch doesn't, so the online network has to predict the target representation rather than simply match it. Collapse is a bad strategy against that: a limited-capacity predictor can't just learn the identity, and a constant output predicts a moving target poorly. The target is an exponential moving average of the online network, momentum around 0.996, so it evolves slowly --- for collapse both networks would have to fall into the same constant simultaneously, and the momentum dynamics resist that.\par\vspace{3pt}
The thing to be careful about is BatchNorm. Early analyses hypothesized the batch statistics were acting as an implicit contrastive term, and that got repeated everywhere. The BYOL authors' own follow-up showed it isn't load-bearing --- BYOL trains successfully with GroupNorm plus weight standardization, which use no batch statistics at all. BatchNorm mainly helps optimization and initialization robustness.\par\vspace{3pt}
SimSiam then showed the momentum encoder isn't required either: predictor plus stop-gradient alone is enough. The clean framing is alternating optimization --- the predictor fits the current targets while the encoder learns to produce targets that are predictable across augmented views.\par\vspace{3pt}
\lab{Numbers}
\begin{itemize}
\item BYOL target EMA momentum: \textasciitilde{}0.996
\end{itemize}
\flab{Loses the signal}
\begin{itemize}
\item ``It works because of the momentum encoder'' --- that misses the predictor, which is the load-bearing piece.
\item Claiming negatives are ``implicitly'' present --- they are not.
\item Repeating the old BatchNorm-as-implicit-contrastive story as settled fact.
\end{itemize}
\lab{Then they ask}
\begin{itemize}
\item What happens if you remove the predictor, or make it too large?
\item How does VICReg prevent collapse explicitly?
\item Is BatchNorm actually necessary for BYOL?
\end{itemize}
\vspace{2pt}\noindent{\color{muted}\rule{\linewidth}{0.3pt}}
\needspace{6\baselineskip}
\qhead{L5}{Trade-off}{Adam vs SGD with momentum---when would you pick each?}
\lab{Say this}
Adam for anything transformer-shaped, SGD with momentum where you have a known vision recipe and generalization is the whole game. The mechanical reason is gradient-scale heterogeneity: in a transformer the embeddings, the attention projections and the layer norms sit at very different gradient magnitudes, and one global step size can't serve all of them. Adam's per-parameter normalization handles that for free, which is why nobody trains an LLM with plain SGD.\par\vspace{3pt}
SGD with momentum still holds up on well-studied convnet tasks with a tuned cosine schedule. The uniform step size acts as implicit regularization --- the extra gradient noise helps you settle in flatter minima --- and if you have the tuning budget and a recipe someone already validated, it's hard to beat. I'd also reach for Adam for fast prototyping, for GANs, and anywhere extensive learning-rate tuning isn't realistic.\par\vspace{3pt}
The nuance I'd volunteer is that the gap has largely closed, and the old ``SGD generalizes better'' folklore mostly came from comparing SGD against Adam with plain L2 rather than decoupled weight decay. In Adam, L2 added into the gradient gets divided by the same running second-moment estimate, so it's attenuated exactly on the parameters with the largest gradients and you end up under-regularized without noticing. AdamW decouples the decay from the adaptive step and mostly erases the difference --- ConvNeXt is the clean example of a modern convnet trained that way.\par\vspace{3pt}
\flab{Loses the signal}
\begin{itemize}
\item ``Adam always, it converges faster'' --- ignores the generalization story on vision
\item ``SGD is obsolete''
\item Not distinguishing L2 in the loss from decoupled weight decay in Adam
\end{itemize}
\lab{Then they ask}
\begin{itemize}
\item Why might Adam generalize worse?
\item How do adaptive learning rates relate to sharp minima?
\item How does weight decay interact differently with each optimizer?
\end{itemize}
\vspace{2pt}\noindent{\color{muted}\rule{\linewidth}{0.3pt}}
\needspace{6\baselineskip}
\qhead{L5}{Debugging}{Training loss is NaN after 1000 steps. Walk through your debugging process.}
\lab{Say this}
NaN means an Inf or an invalid operation got into the graph, and the fact that it's 1000 steps in already tells me something: a bad epsilon or corrupt data usually fires at step one, so I'd start with what grows.\par\vspace{3pt}
First move is to look at gradient norms in the steps before the NaN. If they're climbing --- 1e2, then 1e5, then 1e12, then NaN --- it's exploding gradients, and the fix is clipping at 1.0, warmup, and probably a lower learning rate. If gradient norms are flat and then it's suddenly NaN, that's a different animal: a single bad op or a single bad batch.\par\vspace{3pt}
Then in order of how often it's actually the cause. Check FP16 overflow --- FP16 tops out around 65,500, and anything past that is Inf and then NaN on the next op; switching to BF16 and watching the problem vanish is a one-experiment diagnosis. Check log of zero and divide by zero: use the framework's fused cross-entropy rather than hand-rolling log-softmax, and make sure every normalization has its epsilon. Assert on the input batch, because corrupt data happens --- just less often than people assume.\par\vspace{3pt}
If none of that isolates it, register forward hooks and find the first module that emits a NaN. That's mechanical and always works; I don't start there because the gradient-norm trajectory usually answers it in two minutes.\par\vspace{3pt}
\lab{Numbers}
\begin{itemize}
\item FP16 max \textasciitilde{}65,504; BF16 range \textasciitilde{}3.4e38 (same exponent as FP32)
\end{itemize}
\flab{Loses the signal}
\begin{itemize}
\item ``Set the LR to 1e-6 and hope'' --- that's guessing, and a low enough LR hides the bug while making training useless.
\item Assuming data first; exploding gradients and numerical instability are more common in practice.
\item Never looking at the gradient-norm trajectory before the NaN, which is where the diagnosis actually lives.
\end{itemize}
\lab{Then they ask}
\begin{itemize}
\item What if it only NaNs in mixed precision and never in FP32?
\item How do you debug a NaN that only some ranks see in distributed training?
\item What's the difference in root cause between NaN and Inf?
\end{itemize}
\vspace{2pt}\noindent{\color{muted}\rule{\linewidth}{0.3pt}}
\needspace{6\baselineskip}
\qhead{L6}{Estimation}{How much GPU memory is needed to fine-tune a 13B parameter model with LoRA vs.\textbackslash\{\} full fine-tuning?}
\lab{Say this}
Full fine-tuning a 13B model runs about 240 to 270 gigabytes, LoRA gets you to roughly 60 to 85, and QLoRA to under 40. The gap is not about parameter count, it is about what the optimiser has to hold.\par\vspace{3pt}
Full fine-tuning in BF16 mixed precision with AdamW is 26 gigs of BF16 weights, 52 of FP32 master weights, 104 of momentum and variance, and 26 of gradients. That is 208 before activations, which at batch four, sequence 2048 and forty layers add another 30 to 60 with checkpointing. Four A100-80s with FSDP or ZeRO-3, minimum.\par\vspace{3pt}
With LoRA at rank 16 on the Q and V projections you are training about 13 million parameters, roughly a tenth of a percent. The base model is frozen, so it costs 26 gigs of BF16 weights and nothing else: no master copy, no optimiser state, no gradients for it. The adapter's own weights, gradients and optimiser state together are under a quarter of a gig. Activations are essentially unchanged, because you still run the full forward pass. Total: one A100-80.\par\vspace{3pt}
QLoRA quantises the frozen base to four-bit, so 6.5 gigs instead of 26, landing around 37 total - a single 40-gigabyte A100 or a 48-gigabyte workstation card.\par\vspace{3pt}
The thing to say out loud is that the saving comes from the frozen base needing no optimiser state. The small parameter count on its own would only be worth about 0.2 gigabytes.\par\vspace{3pt}
\lab{Numbers}
\begin{itemize}
\item Full FT 13B: 26 + 52 + 104 + 26 = \textasciitilde{}208 GB before activations
\item LoRA r=16 on Q/V of 40 layers: \textasciitilde{}13M trainable params, \textasciitilde{}0.1\% of the model
\item QLoRA: INT4 base 6.5 GB, \textasciitilde{}37 GB total
\end{itemize}
\flab{Loses the signal}
\begin{itemize}
\item ``LoRA is cheaper because it has fewer parameters'' - true but not the mechanism; the frozen base carrying no optimiser state is.
\item Forgetting activation memory, which is roughly the same for LoRA and full fine-tuning.
\item Not knowing AdamW holds master weights plus two moments per trainable parameter.
\end{itemize}
\lab{Then they ask}
\begin{itemize}
\item What LoRA rank would you start at for 13B?
\item When is full fine-tuning worth the extra hardware?
\item How does LoRA interact with quantisation at inference time?
\end{itemize}
\vspace{2pt}\noindent{\color{muted}\rule{\linewidth}{0.3pt}}
\needspace{6\baselineskip}
\qhead{L5}{Estimation}{Estimate the number of trainable parameters for LoRA with rank 16 on a 7B model.}
\lab{Say this}
About 17 million trainable parameters --- roughly a quarter of a percent of the model.\par\vspace{3pt}
The arithmetic: a 7B like Llama-7B is 32 layers with hidden dimension 4096. Target the attention projections --- Q, K, V, O --- and each is 4096 by 4096. LoRA replaces the update to each with an A and a B, so at rank 16 that's 16 times 4096, twice, about 131 thousand parameters per matrix. Four matrices per layer, 32 layers: 32 times 4 times 131K is 16.8 million. Against a 6.7 billion parameter base that's 0.25 percent. In memory it's about 34 megabytes at FP16 --- nothing.\par\vspace{3pt}
If you also target the feed-forward projections, those aren't square; for Llama they're 4096 by 11008, so each LoRA pair is 16 times the sum of those, about 242 thousand. Three of them per layer over 32 layers adds roughly 23 million, so you land near 40 million total --- still well under one percent.\par\vspace{3pt}
The thing I'd flag is that the parameter count isn't where the memory saving actually comes from. What matters is that you're not holding Adam's two moment tensors for 7 billion parameters. Trainable-parameter count is the headline number; the optimizer state is the reason it fits on the GPU.\par\vspace{3pt}
\lab{Numbers}
\begin{itemize}
\item LoRA r=16, attention only, 7B: 32 layers x 4 matrices x 2 x 4096 x 16 \textasciitilde{}= 16.8M (\textasciitilde{}0.25\% of base)
\item Adding FFN (4096 x 11008) adds \textasciitilde{}23M, for \textasciitilde{}40M total --- still under 1\%
\end{itemize}
\flab{Loses the signal}
\begin{itemize}
\item Forgetting the factor of two --- there is both an A and a B matrix.
\item Not knowing the shape of a 7B: 32 layers, hidden dimension 4096.
\item Quoting total parameters when asked for trainable ones.
\end{itemize}
\lab{Then they ask}
\begin{itemize}
\item How would you decide whether rank 16 is enough?
\item What changes for a 70B?
\item How does this compare with full fine-tuning on memory?
\end{itemize}
\vspace{2pt}\noindent{\color{muted}\rule{\linewidth}{0.3pt}}
\needspace{6\baselineskip}
\qhead{L6}{Debugging}{Your LoRA fine-tuned model performs well on your test set but catastrophically forgets general capabilities. How do you fix it?}
\lab{Say this}
First I'd confirm the forgetting is real rather than a vibe from a handful of prompts --- run general benchmarks like MMLU and HellaSwag before and after, alongside the domain eval. Assuming it's real, the frame is that LoRA reduces forgetting but doesn't prevent it: the base weights are frozen, but the adapter sits in the forward path and can absolutely learn to override base behavior.\par\vspace{3pt}
Then four levers, roughly in that order. Adapter capacity: rank 64 is usually more than the task needs, so drop to 8 or 16 and keep alpha over rank at or below one, which scales the update down. Learning rate: LoRA wants 1e-4 to 3e-4, and if someone has pushed well past that it's doing surgery rather than adaptation. Training length: too many epochs on narrow data overfits the adapter to that distribution, so early-stop on a held-out set that contains general examples as well as domain ones --- domain-only validation will never show you this. And the data mix: fold in ten to twenty percent general-purpose data so the adapter never sees a world made only of your domain.\par\vspace{3pt}
What I'd actually put in place for production is the gate: a general-capability eval suite that runs beside the domain eval, with deployment blocked unless both pass. And I wouldn't reach for full fine-tuning as the fix --- that makes forgetting worse, not better.\par\vspace{3pt}
\lab{Numbers}
\begin{itemize}
\item LoRA LR: 1e-4 to 3e-4 (vs \textasciitilde{}1e-5 for full fine-tuning)
\item Reduce rank 64 $\to$ 8--16, keep $\alpha$/r $\leq$ 1
\item Mix in 10--20\% general-purpose data
\end{itemize}
\flab{Loses the signal}
\begin{itemize}
\item ``LoRA can't cause forgetting because the base weights are frozen'' --- the adapter can override base behavior.
\item Proposing full fine-tuning as the fix, which makes it worse.
\item Never measuring general capability --- no MMLU-style benchmark before and after.
\end{itemize}
\lab{Then they ask}
\begin{itemize}
\item How do you quantify catastrophic forgetting?
\item What if you need aggressive domain adaptation and general capability?
\item Can you combine multiple LoRA adapters?
\end{itemize}
\vspace{2pt}\noindent{\color{muted}\rule{\linewidth}{0.3pt}}
\needspace{6\baselineskip}
\qhead{L6}{Debugging}{Your fine-tuned model echoes the prompt before answering---why?}
\lab{Say this}
Almost always it's missing loss masking on the prompt tokens. If the training loop set labels equal to \texttt{input\_ids} without setting the prompt span to the ignore index, you trained the model to predict prompt tokens as well as response tokens - and since prompts usually outnumber responses, most of your gradient taught it to produce user-turn text. The symptom is exactly what that asks for: restate the prompt, answer, and often invent a fresh user turn afterwards.\par\vspace{3pt}
Verifying takes ten minutes. Pull one collated batch out of the real training pipeline, decode it, and print which positions carry loss. If prompt positions do, that's the bug. The fix is a completion-only collator - TRL's DataCollatorForCompletionOnlyLM, or \texttt{train\_on\_inputs} false in Axolotl - and a retrain, which for LoRA SFT is a single-GPU job measured in hours.\par\vspace{3pt}
Two things I'd rule out first. Template mismatch: if serving renders a different chat template than training, or drops the generation prompt, the model is completing the user's turn instead of starting the assistant's - render the exact inference-time token sequence and diff it against a training example. And base-model residue: a non-instruct base tuned on a few hundred examples can echo from pretraining habits, which is a data problem, not a masking one.\par\vspace{3pt}
The discriminating test is to replay a literal training example through the training-time template offline. A masking bug still echoes. A template mismatch disappears.\par\vspace{3pt}
\flab{Loses the signal}
\begin{itemize}
\item Blames the base model or ``not enough data'' without ever inspecting a collated training batch.
\item Doesn't know the labels and ignore-index mechanics at all.
\item Proposes to hide it - a system prompt telling the model not to echo, or stripping the echo in postprocessing.
\end{itemize}
\lab{Then they ask}
\begin{itemize}
\item You turned on sequence packing and quality dropped - what happened?
\item The model never stops generating - what do you check?
\item How would you catch this in CI before it ships?
\end{itemize}
\vspace{2pt}\noindent{\color{muted}\rule{\linewidth}{0.3pt}}
\needspace{6\baselineskip}
\qhead{L6}{Debugging}{Your training job shows the GPU at 20\% utilization. Diagnose it.}
\lab{Say this}
First, distrust the metric. The nvidia-smi utilization number is the fraction of time any kernel is resident --- one tiny kernel per sampling window reads as busy --- so twenty percent neither proves nor localizes anything. Get a torch.profiler or Nsight trace, then walk the ladder in order of prior probability, each rung with a cheap discriminating test.\par\vspace{3pt}
Data loading is the most common cause. Swap in synthetic pre-loaded batches; if step time collapses it's the input pipeline --- worker count, pinned memory and async prefetch, pre-tokenization, storage throughput. Next, CPU-bound launches: the trace shows sliver kernels separated by gaps with a saturated Python thread, because eager dispatch costs ten to thirty microseconds per op and your kernels are shorter than that. Fix with torch.compile, fused ops, CUDA graphs. Third, kernels too small to fill 132 SMs --- raise the true batch size or pack sequences, and note that gradient accumulation fixes memory, not per-kernel parallelism. Fourth, sync points: recurring cudaStreamSynchronize lined up with the gaps, usually .item() on the loss every step; make logging asynchronous and periodic. Fifth, communication --- exposed NCCL kernels, or per-rank timelines that disagree.\par\vspace{3pt}
And I'd close by renaming the target. Healthy LLM pretraining sits around forty to fifty-five percent MFU. Utilization percent isn't the metric; achieved FLOP/s or tokens per second is.\par\vspace{3pt}
\lab{Numbers}
\begin{itemize}
\item Healthy LLM pretraining: \textasciitilde{}40--55\% MFU
\item Eager-mode dispatch: \textasciitilde{}10--30 $\mu$s per op
\item H100: 132 SMs to fill
\end{itemize}
\flab{Loses the signal}
\begin{itemize}
\item Jumping to kernel optimization or bigger GPUs before ruling out the input pipeline.
\item Taking nvidia-smi utilization at face value, or optimizing it instead of MFU.
\item Listing causes with no discriminating experiment to separate them.
\item Believing gradient accumulation increases per-kernel parallelism.
\end{itemize}
\lab{Then they ask}
\begin{itemize}
\item Utilization is 95\% but MFU is 15\% --- now what?
\item How do you find one straggler among 512 GPUs?
\item Which of these does torch.compile subsume, and which not?
\end{itemize}
\vspace{2pt}\noindent{\color{muted}\rule{\linewidth}{0.3pt}}
\needspace{6\baselineskip}
\qhead{L5}{Conceptual}{What is catastrophic forgetting? How do you mitigate it during fine-tuning?}
\lab{Say this}
Catastrophic forgetting is when fine-tuning on a downstream task overwrites what the model learned in pre-training --- the parameters move to minimize the new loss and destroy the representations that encoded general language understanding. It's worst when the fine-tuning set is small and the learning rate is high, because then you're taking large steps on a narrow slice of the distribution.\par\vspace{3pt}
The single biggest lever is the learning rate: 1e-5 to 5e-5 for fine-tuning a pre-trained model. After that, early stopping --- fine-tuning usually peaks in two or three epochs and everything past that is mostly forgetting. Gradual unfreezing, the ULMFiT recipe, trains the task head first and then unfreezes layers top down, so the lower layers holding general features move least. Weight decay keeps parameters near their pre-trained values, and EWC does that selectively by penalizing changes to parameters that mattered for pre-training. Data replay --- mixing pre-training-style data into your fine-tuning batches --- works and is usually the cheapest thing to try first.\par\vspace{3pt}
But the structural answer is PEFT. With LoRA or adapters the base weights are literally frozen, so the great majority of pre-trained knowledge is preserved by construction rather than by tuning discipline. That's why it's the default now --- though it reduces forgetting rather than eliminating it, since the adapter still sits in the forward path.\par\vspace{3pt}
\lab{Numbers}
\begin{itemize}
\item Fine-tuning LR: 1e-5 to 5e-5
\item Fine-tuning typically peaks in 2--3 epochs
\end{itemize}
\flab{Loses the signal}
\begin{itemize}
\item Cannot actually define catastrophic forgetting when asked.
\item ``Use a small learning rate'' and nothing else.
\item Not connecting PEFT to forgetting prevention.
\end{itemize}
\lab{Then they ask}
\begin{itemize}
\item What is Elastic Weight Consolidation doing?
\item Does forgetting differ for encoder-only versus decoder-only models?
\item Can you recover from forgetting after the fact?
\end{itemize}
\vspace{2pt}\noindent{\color{muted}\rule{\linewidth}{0.3pt}}
\chapter{Inference \& Serving}
\addtocontents{toc}{\protect\thispagestyle{fancy}}
\needspace{6\baselineskip}
\qhead{L6}{Estimation}{Calculate the KV cache memory for serving LLaMA 70B at 128K context length.}
\lab{Say this}
About 42 gigabytes for a single sequence, and that number is the punchline. You store a key and a value per layer per KV head, so it's two, times 80 layers, times 8 KV heads under grouped-query attention, times head dimension 128, times two bytes for BF16 --- roughly 330 kilobytes per token. Multiply by 128 thousand tokens and you land at 42 gigabytes.\par\vspace{3pt}
The comparison is what matters: the cache for one user's context is larger than the entire model quantized to INT4, which is about 35 gigabytes. That's why long-context serving is a memory problem rather than a compute problem --- you can't fit batch, and without batch you can't get utilization out of the GPU.\par\vspace{3pt}
Two things I'd flag. It's exactly linear in context length, so 4K is about 1.3 gigabytes, 32K about 10, 128K about 42 --- no cliff, just a line you eventually run into, and it's per concurrent sequence. And this is already after grouped-query attention saved you: with 64 full KV heads instead of 8 you'd be at over 300 gigabytes for the same request. If asked how to shrink it, I'd go to KV cache quantization to INT8 or FP8, paged attention so fragmentation doesn't waste capacity, and prefix sharing when many requests start from the same system prompt.\par\vspace{3pt}
\lab{Numbers}
\begin{itemize}
\item LLaMA 70B at 128K: \textasciitilde{}42 GB KV cache per sequence in BF16
\item \textasciitilde{}330 KB of KV cache per token (80 layers, 8 KV heads, $d_{head}$ 128)
\item Same model's INT4 weights: \textasciitilde{}35 GB
\end{itemize}
\flab{Loses the signal}
\begin{itemize}
\item Dropping the factor of 2 for keys and values
\item Using all 64 attention heads instead of the 8 GQA KV heads
\item Not seeing that it grows linearly with context and per concurrent sequence
\end{itemize}
\lab{Then they ask}
\begin{itemize}
\item How would you reduce this KV cache?
\item What batch size fits on 8$\times$H100 at 80 GB each?
\item How does prefix caching help a chatbot workload?
\end{itemize}
\vspace{2pt}\noindent{\color{muted}\rule{\linewidth}{0.3pt}}
\needspace{6\baselineskip}
\qhead{L6}{Trade-off}{INT8 vs INT4 vs FP8 quantization---what is your decision framework?}
\lab{Say this}
I choose on hardware first, then accuracy budget. On Hopper and newer, FP8 in E4M3 is the default: native tensor-core support, per-tensor dynamic scaling, half the memory of FP16, near-zero quality loss, almost no kernel engineering. FP8 beats INT8 at the same bit width because of dynamic range - floating point absorbs activation outliers, where INT8's fixed scale and zero point get stretched by them.\par\vspace{3pt}
On A100-class or older hardware, or on CPU, INT8 is the answer, and the technique is SmoothQuant. It migrates quantization difficulty out of activations, which have nasty outlier channels, into weights, which are well behaved, and that's what makes W8A8 work at all. Expect under half a percent perplexity change.\par\vspace{3pt}
INT4 is for when you must fit a bigger model on fewer GPUs and you'll pay for it. GPTQ quantizes layer by layer using approximate second-order information, compensating with the weights it hasn't done yet; AWQ is faster and protects the roughly one percent of weight channels that see large activations. Both give usable models, but degradation lands on reasoning, math and code rather than general text, so benchmark on your task, not perplexity.\par\vspace{3pt}
Two places this has moved: FP8 is no longer H100-only, since Ada-generation cards have it too, and Blackwell adds native FP4, which makes four-bit a first-class hardware path.\par\vspace{3pt}
\lab{Numbers}
\begin{itemize}
\item INT8: <0.5\% perplexity loss; INT4: \textasciitilde{}1-3\%
\item AWQ protects the \textasciitilde{}1\% of weight channels with the largest activations
\end{itemize}
\flab{Loses the signal}
\begin{itemize}
\item ``INT4 always, it uses least memory'' - the degradation lands on reasoning and code, and it needs a stated accuracy budget.
\item ``FP8 and INT8 are the same eight bits'' - FP8 has floating-point dynamic range, which is the whole point for activations.
\item Validates only on perplexity instead of the task the model actually does.
\end{itemize}
\lab{Then they ask}
\begin{itemize}
\item How does SmoothQuant move difficulty from activations into weights?
\item Why does INT4 hurt math and code more than general text?
\item Can you combine INT4 weights with INT8 activations?
\end{itemize}
\vspace{2pt}\noindent{\color{muted}\rule{\linewidth}{0.3pt}}
\needspace{6\baselineskip}
\qhead{L5}{Conceptual}{Explain speculative decoding. When would you use it and when would you not?}
\lab{Say this}
Speculative decoding is a latency trick that gives you the exact same output distribution as the target model. A small draft model proposes K tokens and the target verifies all K in one forward pass. The insight is that verification is parallel --- prefill-shaped --- while ordinary decoding is strictly sequential. So if the draft is usually right you get several tokens for roughly the cost of one target pass.\par\vspace{3pt}
The acceptance step is what makes it lossless. If the target assigns the drafted token at least as much probability as the draft did, accept it. Otherwise accept with probability target over draft, and on rejection resample from the normalized positive difference between the two distributions. That's provably equivalent to sampling from the target directly, so quality is identical --- not better, identical. That's the thing people get wrong in interviews.\par\vspace{3pt}
I'd use it for single-request latency on a large model, for structured output like JSON or code where draft agreement is high, and generally in low-batch serving where the GPU is sitting mostly idle during decode anyway.\par\vspace{3pt}
I would not use it at high batch, where the GPU is already compute-saturated and the extra verification work makes you slower. Not for high-temperature creative generation, where acceptance collapses. And it's pointless without a draft model agreeing on more than roughly half the tokens.\par\vspace{3pt}
\lab{Numbers}
\begin{itemize}
\item Rule of thumb: need >50\% per-token draft agreement before it pays
\end{itemize}
\lab{If they hand you the marker}
$$\text{accept } \tilde{x}\sim q \text{ w.p. } \min\!\left(1, \tfrac{p(\tilde{x})}{q(\tilde{x})}\right), \quad\text{else resample from } \tfrac{\max(0,\,p-q)}{\lVert\max(0,\,p-q)\rVert_1}$$ Write the accept rule first, then the corrected residual distribution underneath, and say out loud that this pair is exactly why the sampled distribution is the target's and not an approximation of it.\par
\flab{Loses the signal}
\begin{itemize}
\item Saying it improves quality --- it guarantees identical quality, that's the whole point.
\item ``It always speeds things up'' --- at high batch it can be a net slowdown.
\item Confusing it with beam search or parallel sampling.
\end{itemize}
\lab{Then they ask}
\begin{itemize}
\item How do you pick or train the draft model?
\item What acceptance rate do you need for a 2$\times$ speedup?
\item Does this help throughput or only latency?
\end{itemize}
\vspace{2pt}\noindent{\color{muted}\rule{\linewidth}{0.3pt}}
\needspace{6\baselineskip}
\qhead{L6}{Debugging}{KV cache memory is your bottleneck---you are running out of GPU memory and dropping requests. What do you do?}
\lab{Say this}
I'd order this by effort, and the first thing to be clear about is that this is the KV pool, not the weight pool --- quantizing weights doesn't buy you a byte here.\par\vspace{3pt}
Cheapest first, no model changes. PagedAttention if the stack doesn't have it: static pre-allocation loses a large fraction of KV to fragmentation and paging typically recovers 40 to 60 percent of that waste. Prefix caching if requests share a long system prompt. Admission control, so you queue against available KV rather than OOMing --- that converts dropped requests into slower ones. And a max context length, so a few enormous requests can't starve everyone.\par\vspace{3pt}
Next tier is compression. Quantize the KV cache itself to FP8 or INT8, separately from the weights --- the cache tolerates low precision well because softmax is fairly robust to small perturbations, and that's a clean 2x. If you get to choose the model, GQA cuts the cache by the ratio of attention heads to KV heads, and DeepSeek-style MLA goes much further by caching a small compressed latent, around 57 times smaller than the MHA-equivalent. Sliding-window models like Mistral bound cache growth outright.\par\vspace{3pt}
Beyond that: offload cold KV to host memory, evict low-attention tokens with something like H2O --- risky, it silently breaks long-range retrieval --- or add GPUs. My order: paging, KV quantization, prefix caching, then architecture.\par\vspace{3pt}
\lab{Numbers}
\begin{itemize}
\item PagedAttention: recovers \textasciitilde{}40--60\% of KV memory lost to fragmentation
\item KV quantization to INT8/FP8: \textasciitilde{}2$\times$ cache reduction, minimal quality impact
\item MLA (DeepSeek-V3): cached latent \textasciitilde{}57$\times$ smaller than the MHA equivalent
\end{itemize}
\flab{Loses the signal}
\begin{itemize}
\item Proposing weight quantization as a fix for KV pressure --- different memory pools.
\item ``Add more GPUs'' as the first move, before any software mitigation.
\item No mention of admission control, so the failure mode stays OOM instead of queueing.
\end{itemize}
\lab{Then they ask}
\begin{itemize}
\item How much KV does your model use at 32K context? Do the arithmetic.
\item What's the quality cost of KV cache quantization?
\item Could RAG cut the context length instead?
\end{itemize}
\vspace{2pt}\noindent{\color{muted}\rule{\linewidth}{0.3pt}}
\needspace{6\baselineskip}
\qhead{L6}{Estimation}{How much memory savings does INT4 quantization give for a 70B model? What is the quality trade-off?}
\lab{Say this}
Roughly 4x on weights. Seventy billion parameters at FP16 is 140 gigabytes; at INT4 it's 35. That takes you from needing two H100s to fitting on one with room to spare. Add the overhead honestly. GPTQ and AWQ store a scale and zero-point per group, usually groups of 128 weights --- about four extra bytes per group, call it 6 percent. So realistically 37 gigabytes, not 35. Still one GPU.\par\vspace{3pt}
Then the part that trips people up: this shrinks weights only. The KV cache is a separate pool, still FP16 unless you quantize it separately, and prefill activations are FP16 too. At 32K-plus context the KV cache can exceed the weight memory outright, so INT4 weights alone do not solve memory pressure on a long-context serving stack.\par\vspace{3pt}
On quality, 4-bit 70B typically lands within about a point of FP16 on MMLU, HumanEval and GSM8K. But degradation isn't uniform, which is the part worth saying: reasoning-heavy and multi-step math degrade more than factual recall, and long-context tasks degrade because quantization error compounds across layers. So evaluate on your task, not the aggregate.\par\vspace{3pt}
Two things I'd land on. A 4-bit 70B beats an FP16 7B almost every time --- bigger and quantized beats smaller and exact. And 4x smaller is not 4x faster: dequantization kernels cost something and you may be bandwidth-bound regardless.\par\vspace{3pt}
\lab{Numbers}
\begin{itemize}
\item 70B: 140 GB at FP16 $\to$ 35 GB at INT4, \textasciitilde{}37 GB with group scales
\item Group-128 scale + zero-point overhead: \textasciitilde{}6.25\%
\item 4-bit 70B: typically <1\% drop on MMLU / HumanEval / GSM8K
\end{itemize}
\flab{Loses the signal}
\begin{itemize}
\item ``4$\times$ smaller, no quality loss'' --- there is always some loss; the question is whether it lands on your task.
\item Conflating weight quantization with KV cache quantization.
\item Claiming a 4$\times$ speedup follows from a 4$\times$ memory reduction.
\end{itemize}
\lab{Then they ask}
\begin{itemize}
\item When would you pick FP8 over INT4?
\item How would you evaluate quantization quality for your specific workload?
\item Where does sub-4-bit fall off a cliff?
\end{itemize}
\vspace{2pt}\noindent{\color{muted}\rule{\linewidth}{0.3pt}}
\needspace{6\baselineskip}
\qhead{L6}{Trade-off}{Tensor parallelism vs pipeline parallelism for inference---what is your decision framework?}
\lab{Say this}
The decision is mostly about interconnect. Tensor parallelism splits every layer across GPUs and needs an all-reduce at each one --- two per transformer block, after the attention output projection and after the MLP down projection. Pipeline parallelism gives whole layer groups to each GPU and only communicates at stage boundaries, but pays for that with bubbles.\par\vspace{3pt}
Inside a node with NVLink, always TP. Around 900 gigabytes a second per H100, and at batch one each all-reduce is only about 16 kilobytes, so cost is per-collective latency rather than bandwidth --- ten microseconds each, twenty per layer, about 1.6 milliseconds across 80 layers against a decode step of maybe five. Tolerable, and TP is what actually cuts per-token latency. TP=8 on one node is the standard way to serve a 70B.\par\vspace{3pt}
Across nodes, don't stretch TP. The bandwidth math looks deceptively cheap for 16-kilobyte messages, but InfiniBand's per-collective latency is ten to twenty microseconds and you're now paying it 160 times per token --- two to four milliseconds of pure communication per token, which dominates everything. So TP inside each node, PP between nodes.\par\vspace{3pt}
Two refinements. Latency-sensitive chatbot: maximize TP degree, accept the overhead. Throughput-sensitive batch work: prefer data parallelism over PP where the model fits --- replicate, keep each replica busy, avoid bubbles entirely. PP is really the answer when the model doesn't fit any other way.\par\vspace{3pt}
\lab{Numbers}
\begin{itemize}
\item NVLink: \textasciitilde{}900 GB/s per H100; per-collective latency \textasciitilde{}10 $\mu$s
\item TP=8 over 80 layers: \textasciitilde{}1.6 ms/token of comms on NVLink vs \textasciitilde{}2--4 ms cross-node
\end{itemize}
\flab{Loses the signal}
\begin{itemize}
\item ``TP is always better'' --- it needs NVLink-class interconnect; over InfiniBand high-degree TP kills you.
\item Forgetting pipeline bubbles, which waste 20--40\% of cycles without careful micro-batching.
\item Never considering data parallelism when the model fits in one replica.
\end{itemize}
\lab{Then they ask}
\begin{itemize}
\item At what TP degree does communication start to dominate?
\item What do you do when the model doesn't divide evenly across GPUs?
\item What about expert parallelism for an MoE?
\end{itemize}
\vspace{2pt}\noindent{\color{muted}\rule{\linewidth}{0.3pt}}
\needspace{6\baselineskip}
\qhead{L6}{Trade-off}{Your model needs to run in $<$10ms. Current latency is 50ms. Walk through the optimization decision tree.}
\lab{Say this}
Five times, and the first thing I'd do is refuse to optimize before profiling. Break the fifty milliseconds into feature computation, model forward, post-processing and network. If features are thirty of it, every model trick in the world caps you at twenty --- you need precomputation or caching, and that becomes the whole project. Forgetting the feature stage is the classic mistake on this question.\par\vspace{3pt}
Assuming the model is the bottleneck, cheap wins first. INT8 quantization is typically two to four times for under a point of quality. Graph compilation --- TensorRT or ONNX Runtime --- fuses operators for another one and a half to three. Flash attention if attention dominates. Together those plausibly take fifty down to fifteen or twenty in a day.\par\vspace{3pt}
If that's not enough, distillation is the next rung: train a small student on the large model's outputs. Twelve layers down to three is roughly four times at ninety-five to ninety-seven percent of quality, one to two weeks of work. Structured pruning of heads and layers is the alternative.\par\vspace{3pt}
Below that: swap in a genuinely smaller architecture, cache predictions so repeated inputs skip inference entirely, or early-exit when the model is already confident. One caution --- dynamic batching helps throughput and hurts tail latency, so it's the wrong lever here. My plan: profile in half an hour, INT8 plus a compiler for three to four times, then distillation if we're short.\par\vspace{3pt}
\lab{Numbers}
\begin{itemize}
\item INT8: typically 2--4$\times$ speedup at <1\% quality loss
\item TensorRT / ONNX graph fusion: 1.5--3$\times$
\item 12-layer $\to$ 3-layer distillation: \textasciitilde{}4$\times$ at 95--97\% of quality
\end{itemize}
\flab{Loses the signal}
\begin{itemize}
\item ``Use a smaller model'' without profiling where the 50 ms actually goes.
\item Not ordering the options by effort-to-impact.
\item Forgetting feature computation, which is often the real bottleneck.
\end{itemize}
\lab{Then they ask}
\begin{itemize}
\item What if quantization costs you 3\% accuracy?
\item How do you choose between distillation and pruning?
\item What if the model is already INT8?
\end{itemize}
\vspace{2pt}\noindent{\color{muted}\rule{\linewidth}{0.3pt}}
\needspace{6\baselineskip}
\qhead{L6}{Estimation}{I give you an op---say, a LayerNorm over a $[16{,}384 \times 8{,}192]$ BF16 tensor, or a $4096^3$ GEMM. Is it compute-bound or bandwidth-bound on an H100? Show your method.}
\lab{Say this}
Method first: count FLOPs, count HBM bytes moved, take the ratio, and compare it to the machine's ridge point. On an H100 that's about 989 teraFLOP/s dense BF16 over 3.35 terabytes a second, so the ridge is around 295 FLOPs per byte. Runtime is the max of compute time and memory time, never the sum.\par\vspace{3pt}
LayerNorm on that tensor: 134 million elements at a handful of FLOPs each, so order a gigaFLOP, against a read plus a write of about half a gigabyte. Intensity around two against a ridge of 295 --- hopelessly bandwidth-bound. Time is bytes over bandwidth, roughly 160 microseconds, and no kernel cleverness beats that; the only win is fusing it with its neighbours so the tensor gets touched once.\par\vspace{3pt}
The 4096-cubed GEMM: about 140 gigaFLOPs over roughly a hundred megabytes, intensity around thirteen hundred, well past the ridge. Compute-bound, about 140 microseconds at peak, call it 0.2 milliseconds at realistic efficiency.\par\vspace{3pt}
Three reflexes I'd add. For matmuls, intensity is governed by the smallest dimension --- a batch-eight decode GEMM has intensity around eight and is bandwidth-bound despite being a matmul, which is the whole story of why LLM decoding is slow. Count HBM bytes, not instruction-level loads, since cache reuse is the entire point. And precision moves both axes: FP8 doubles the compute roof and halves the bytes, pushing the ridge to about 590.\par\vspace{3pt}
\lab{Numbers}
\begin{itemize}
\item H100: \textasciitilde{}989 TFLOP/s dense BF16, \textasciitilde{}3.35 TB/s HBM $\to$ ridge $\approx$ 295 FLOP/byte
\item LayerNorm [16,384 $\times$ 8,192] BF16: AI $\approx$ 2, \textasciitilde{}160 $\mu$s, bandwidth-bound
\item 4096$^3$ GEMM: AI $\approx$ 1,365, compute-bound, \textasciitilde{}0.2 ms realistic
\end{itemize}
\lab{If they hand you the marker}
$$\text{AI}=\frac{\text{FLOPs}}{\text{HBM bytes}},\quad \text{ridge}=\frac{989\ \text{TFLOP/s}}{3.35\ \text{TB/s}}\approx 295,\quad t=\max\!\left(\frac{\text{FLOPs}}{989\text{T}},\ \frac{\text{bytes}}{3.35\text{T}}\right)$$ Write the two anchor numbers first and divide them to get the ridge, then put the op's intensity on the same line and point at which side of 295 it lands. Say ``max, not sum'' out loud as you write the last term.\par
\flab{Loses the signal}
\begin{itemize}
\item Not knowing the anchor numbers --- 989 TFLOP/s dense BF16 and \textasciitilde{}3.35 TB/s --- the method dies without them.
\item Quoting 1,979 TFLOP/s as dense BF16 peak; that's the 2:4-sparsity figure, and \textasciitilde{}1,979 is legitimately FP8 dense.
\item ``It's a matmul, so it's compute-bound'' without checking the smallest dimension.
\item Adding compute time and memory time instead of taking the max.
\end{itemize}
\lab{Then they ask}
\begin{itemize}
\item Fuse the LayerNorm with the residual add --- what changes?
\item At what batch size does a d=8192 decode GEMM cross over?
\item Is attention prefill compute- or bandwidth-bound? And decode?
\end{itemize}
\vspace{2pt}\noindent{\color{muted}\rule{\linewidth}{0.3pt}}
\needspace{6\baselineskip}
\qhead{L7}{Estimation}{Estimate the maximum tokens/second for a 70B dense model on 8$\times$H100---for training, and for inference.}
\lab{Say this}
Training first. Six times parameters per token, so 4.2e11 FLOPs per token. Eight H100s is about 7.9 petaFLOP/s dense BF16 peak; at a well-optimized forty to fifty percent MFU you get three to four effective, so roughly seventy-five hundred to ninety-five hundred tokens a second. Then the caveat that separates a Staff answer: eight H100s is 640 gigabytes, and full training state for a 70B is about 1.1 terabytes before activations --- sixteen bytes a parameter for BF16 weights, FP32 master and Adam moments. The arithmetic is right but the recipe isn't feasible as stated; you'd need ZeRO-3 with offload, which erodes MFU, or LoRA.\par\vspace{3pt}
Inference splits into two regimes, and one number is the classic failure. At batch one you're bandwidth-bound: you read all 140 gigabytes of weights per token, 17.5 per GPU at 3.35 terabytes a second is 5.2 milliseconds, plus 1.6 of tensor-parallel all-reduce latency --- six to seven milliseconds a token, or 140 to 170 tokens a second roofline. Production stacks land at fifty to ninety.\par\vspace{3pt}
At large batch, weights are read once per step regardless of batch, intensity grows with batch, and you turn compute-bound somewhere around batch 150. Past that it's effective FLOP/s over 2P per token: roughly fifteen to twenty-five thousand output tokens a second aggregate. Sanity check --- that's within a small factor of training, as it must be, since training pays three times the FLOPs per token.\par\vspace{3pt}
\lab{Numbers}
\begin{itemize}
\item Training: 6P = 4.2e11 FLOPs/token; 8$\times$H100 $\approx$ 7.9 PFLOP/s peak $\to$ \textasciitilde{}8.5K tok/s at 45\% MFU
\item Batch-1 decode: \textasciitilde{}6--7 ms/token $\to$ \textasciitilde{}140--170 tok/s roofline, \textasciitilde{}50--90 observed
\item Large batch: \textasciitilde{}15K--25K output tok/s aggregate, compute-bound past B $\approx$ 150
\end{itemize}
\lab{If they hand you the marker}
$$\text{train}:\ \frac{8\times 989\text{T}\times \text{MFU}}{6P}\ \text{tok/s} \qquad \text{decode }(B{=}1):\ t \approx \frac{2P\ \text{bytes}}{N_{\text{GPU}}\times 3.35\ \text{TB/s}} + t_{\text{comm}}$$ Write the training line, then the batch-1 line beneath it, and label them ``compute-bound'' and ``bandwidth-bound'' before computing anything --- the interviewer is grading whether you split the regimes, not your arithmetic.\par
\flab{Loses the signal}
\begin{itemize}
\item One inference number with no batch-size regime split --- the single most common failure here.
\item Using 1,979 TFLOP/s as dense BF16 peak, or assuming 100\% MFU.
\item No memory-feasibility check on the training side.
\item Forgetting that at long context the KV read joins the weight read and erodes the batch-1 number.
\end{itemize}
\lab{Then they ask}
\begin{itemize}
\item Which numbers change if it's an 8$\times$22B MoE instead?
\item Add a 32K prompt --- what are the new TTFT and decode rate?
\item Why do training MFU and decode-GEMM efficiency differ so much?
\end{itemize}
\vspace{2pt}\noindent{\color{muted}\rule{\linewidth}{0.3pt}}
\chapter{NLP \& LLMs}
\addtocontents{toc}{\protect\thispagestyle{fancy}}
\needspace{6\baselineskip}
\qhead{L5}{Conceptual}{How does subword tokenization (BPE) interact with embedding quality?}
\lab{Say this}
Tokenization sets the granularity at which embeddings are learned, so it decides which words get a trained vector of their own and which have to be reassembled from pieces. Frequent words get a single token and an embedding updated millions of times. Rare words get split - photosynthesis becomes three or four fragments - and the model composes their meaning through attention and the feedforward layers. That works, but it's noisier, because it has seen that particular combination far less often.\par\vspace{3pt}
So vocabulary size is a real trade-off, not a free parameter. Llama 3 went to 128K from Llama 2's 32K, and GPT-2 was 50K. Bigger vocabulary means more words get their own token, better per-token semantics, shorter sequences, more text per context window, cheaper inference. The cost is a much larger embedding table and each token seen less often in training, so the right size depends on how much data you have.\par\vspace{3pt}
Where it bites is domain and language. English-centric tokenizers shred other languages and code identifiers into fragments, so you get worse per-token embeddings and longer sequences at once. For retrieval over medical terms or product SKUs, expect fragmentation to hurt, and expect domain fine-tuning to help mainly because the model learns to compose those fragments correctly.\par\vspace{3pt}
\lab{Numbers}
\begin{itemize}
\item Vocabulary sizes: GPT-2 50K, Llama 2 32K, Llama 3 128K
\end{itemize}
\flab{Loses the signal}
\begin{itemize}
\item ``BPE is just OOV handling'' - it shapes the structure of the whole embedding space.
\item ``Bigger vocabulary is always better'' - the optimum depends on training data size.
\item Doesn't know English-centric tokenizers disproportionately fragment other languages and code.
\end{itemize}
\lab{Then they ask}
\begin{itemize}
\item How would you pick a vocabulary size for a new language model?
\item What does tokenization do to a multilingual model?
\item What problem are Matryoshka embeddings solving?
\end{itemize}
\vspace{2pt}\noindent{\color{muted}\rule{\linewidth}{0.3pt}}
\needspace{6\baselineskip}
\qhead{L5}{Trade-off}{BERT vs GPT pre-training objectives---what are the fundamental tradeoffs between masked language modeling and autoregressive language modeling?}
\lab{Say this}
The trade is bidirectional context against loss coverage, and the honest version of it is subtler than the story people usually tell. BERT masks about 15 percent of tokens and predicts them from both sides, so each prediction is better informed while gradient arrives at only 15 percent of the positions you paid to process. GPT predicts every token from left context alone, which is weaker information per position but a loss term everywhere --- call it six or seven times more loss terms per sequence.\par\vspace{3pt}
I would not oversell that as data efficiency, because loss coverage is not sample efficiency. At matched scale the masked models were actually more sample-efficient on understanding benchmarks, RoBERTa against GPT-2, and masking rates well above 15 percent turn out to work better than the original recipe. The defensible claim is more loss terms per FLOP, and that compounds with the two things that really settled it: a prompt-to-completion interface that covers any task without a task-specific head, and in-context learning, which only shows up at scale. That is why every frontier model is decoder-only.\par\vspace{3pt}
Masked modelling also carries a train-test mismatch --- the mask token exists in pretraining and never again. The 80/10/10 rule, where of the selected positions 80 percent are masked, 10 percent replaced at random and 10 percent left alone, patches that without curing it. Causal modelling has no such gap, and encoders now survive mainly as cheap fixed-size embedders.\par\vspace{3pt}
\lab{Numbers}
\begin{itemize}
\item BERT masks 15\% of tokens; of those, 80/10/10 mask/random/unchanged
\item MLM gets gradient from \textasciitilde{}15\% of positions; next-token prediction from every position
\item Next-token prediction computes \textasciitilde{}6.7x more loss terms per sequence than 15\% MLM
\end{itemize}
\flab{Loses the signal}
\begin{itemize}
\item ``Bidirectional is always better'' as the whole answer
\item Not knowing the signal-density difference --- 15 percent of positions versus all of them
\item Missing the mask-token train/test mismatch and what 80/10/10 is for
\end{itemize}
\lab{Then they ask}
\begin{itemize}
\item Why not just use bidirectional attention inside a decoder?
\item Does T5's span corruption get the best of both?
\item How does ELECTRA fix the signal-density problem?
\end{itemize}
\vspace{2pt}\noindent{\color{muted}\rule{\linewidth}{0.3pt}}
\needspace{6\baselineskip}
\qhead{L6}{First Principles}{Next-token prediction vs masked language modeling---why did GPT's autoregressive approach win for generative AI while BERT's bidirectional approach initially dominated NLU benchmarks?}
\lab{Say this}
Three reasons: task alignment, the scaling story, and interface universality.\par\vspace{3pt}
Task alignment is the big one. An autoregressive model trains on exactly the distribution it's asked to sample from at inference --- predict the next token, generate the next token, no gap. BERT trains to fill in blanks, which is per-position classification. To generate with it you need some iterative unmasking scheme it was never trained for, and it comes out incoherent.\par\vspace{3pt}
On scaling, the honest version of the signal-density argument is that causal LM puts a loss term on every token where masked LM scores only the fifteen percent it masks --- about seven times more loss terms per sequence. I'd be careful not to oversell that as sample efficiency, because at matched scale the masked models were actually more sample-efficient on NLU, and higher masking rates narrow the gap. The defensible claim is more loss terms per FLOP. What compounds it is that perplexity is a smooth scaling metric, and the emergent behaviors --- in-context learning, chain of thought --- appear only through the generative interface.\par\vspace{3pt}
That interface is the third reason. Prompt in, completion out expresses every NLP task, so one model replaces a rack of task-specific encoders.\par\vspace{3pt}
And why BERT won first: bidirectional context genuinely gives strictly more information for understanding tasks, and at a few hundred million parameters that's decisive. Decoders caught up on NLU only around ten billion parameters.\par\vspace{3pt}
\lab{Numbers}
\begin{itemize}
\item MLM scores \textasciitilde{}15\% of tokens per sequence, causal LM 100\% --- \textasciitilde{}6.7x more loss terms
\end{itemize}
\flab{Loses the signal}
\begin{itemize}
\item ``GPT is just better'' with no structural reason.
\item Overselling loss-term density as sample efficiency --- at matched scale MLM models were more sample-efficient on NLU.
\item Not conceding that encoders remain the better tool for some understanding tasks at small scale.
\end{itemize}
\lab{Then they ask}
\begin{itemize}
\item Could you design a bidirectional model that generates well?
\item Why didn't T5's encoder-decoder win the scaling race?
\item What about diffusion language models?
\end{itemize}
\vspace{2pt}\noindent{\color{muted}\rule{\linewidth}{0.3pt}}
\needspace{6\baselineskip}
\qhead{L6}{Debugging}{Your fine-tuned chat model performs worse than the base model it was tuned from. Walk me through your diagnosis.}
\lab{Say this}
When fine-tuning makes a model worse, it's almost always a bookkeeping bug rather than a modelling problem, so I check the artifact before I touch the algorithm. First suspect is chat-template mismatch: training serialized conversations one way and serving does it another. The test takes fifteen minutes --- take one training example, push it through the serving stack, and diff the exact token IDs against what the training pipeline produced. You're looking for hand-rolled prompt strings, a doubled BOS because the template inserts one and the framework prepends another, and missing role headers. And note that training loss looks perfectly healthy under a wrong template; the model learns the malformed format just fine.\par\vspace{3pt}
Second, stop tokens. If generations ramble forever or truncate mid-sentence, either the turn-end token was masked out of the loss so the model never learned to emit it, or serving stops on document EOS while the model emits a turn-end token. Check whether the model ever produces the configured stop token at all.\par\vspace{3pt}
Third, tokenizer drift --- a library version bump, special tokens added that shifted IDs, or an embedding matrix resized in training while the old tokenizer files shipped.\par\vspace{3pt}
Only once those are excluded do I call it forgetting from too high a learning rate or too many epochs, and I'd want a broad, uniform drop across general benchmarks to believe it. The unifying point is that the deployable artifact is weights plus tokenizer plus template plus stop config, versioned together.\par\vspace{3pt}
\flab{Loses the signal}
\begin{itemize}
\item Reaching for ``more data, bigger model'' before checking serialization
\item Not knowing training loss looks fine under a wrong chat template
\item Treating EOS as one token with one job --- document end versus turn end
\end{itemize}
\lab{Then they ask}
\begin{itemize}
\item Why does training loss still go down with the wrong template?
\item How would you detect a doubled BOS in production?
\item Does LoRA change the forgetting diagnosis?
\end{itemize}
\vspace{2pt}\noindent{\color{muted}\rule{\linewidth}{0.3pt}}
\needspace{6\baselineskip}
\qhead{L6}{Debugging}{Your fine-tuned LLM generates fluent but factually wrong answers. Diagnose the issue and propose solutions.}
\lab{Say this}
That's hallucination, and I'd separate three causes before proposing anything, because the fix is different for each.\par\vspace{3pt}
First, is the knowledge actually in the model? Compare the base model against the fine-tuned one on a factual QA benchmark. If the base was right and the fine-tune is wrong, you degraded pretrained knowledge --- learning rate too high, too many epochs, catastrophic forgetting. It's also worth hand-sampling the fine-tuning data, because if it contains factual errors the model faithfully memorized them, and if your reward signal only scored style, you explicitly trained fluency over accuracy.\par\vspace{3pt}
Second, is it decoding? High temperature or wide top-p pushes toward fluent but low-probability continuations. Rerun the eval at temperature zero. If accuracy jumps, the knowledge is there and sampling is introducing the errors --- for factual tasks I'd run greedy or keep temperature at 0.3 or below.\par\vspace{3pt}
Third, is it a knowledge boundary? Models don't know what they don't know and answer confidently regardless. Check whether the errors concentrate on recent events, niche domains, or numbers.\par\vspace{3pt}
On fixes, retrieval is far and away the highest-leverage one --- put the actual documents in context and the answer is grounded rather than recalled. Layered on top: train the model to abstain when it can't support an answer, require citations so hallucinations become detectable rather than invisible, and use self-consistency across samples as a cheap uncertainty signal for the cases where you can't retrieve.\par\vspace{3pt}
\flab{Loses the signal}
\begin{itemize}
\item ``Train on more data'' --- volume alone doesn't fix hallucination without a targeted intervention.
\item Never checking decoding parameters before blaming the model or the data.
\item Not reaching for retrieval, which is the highest-leverage practical fix.
\end{itemize}
\lab{Then they ask}
\begin{itemize}
\item How do you measure hallucination rate?
\item When does RAG fail?
\item How does RLHF contribute to hallucination?
\end{itemize}
\vspace{2pt}\noindent{\color{muted}\rule{\linewidth}{0.3pt}}
\needspace{6\baselineskip}
\qhead{L6}{Conceptual}{Why do decoder-only models dominate modern NLP despite encoder-decoder being theoretically more flexible?}
\lab{Say this}
They dominate for interface and infrastructure reasons, not because they're better models. The honest version is that at matched FLOPs, the T5 and UL2 era comparisons show encoder-decoder is compute-competitive on plenty of tasks, so I wouldn't claim decoder-only uses compute more efficiently --- that isn't settled, and claiming it is invites a bad follow-up.\par\vspace{3pt}
What is solid is three things. Interface: one text-in, text-out model handles every task through prompting --- classification, translation, summarization, code, conversation --- where encoder-decoder wants structured input-output pairs and is awkward for open-ended chat. One decoder-only model replaces a shelf of specialized ones. Data: decoder-only pretrains on raw text with no formatting at all, while encoder-decoder needs paired data or a corruption scheme like T5's span corruption, and raw text is vastly more abundant. Infrastructure: one forward path, one KV cache, symmetric batching. Encoder-decoder means two forward passes, a cross-attention cache and asymmetric batching, and that complexity compounds across training, serving and optimization.\par\vspace{3pt}
Encoder-decoder still wins where the input genuinely needs bidirectional encoding and the output is short. Machine translation is the classic case, Whisper is encoder-decoder for speech, and long-input short-output summarization. And a prefix LM sits in between --- bidirectional attention over the prompt, causal generation after it --- which is the thing to raise if the interviewer pushes on whether it has to be either-or.\par\vspace{3pt}
\flab{Loses the signal}
\begin{itemize}
\item ``Decoder-only is just objectively better'' --- the dominance is about scaling and infrastructure, not modeling power.
\item Asserting decoder-only is more compute-efficient at matched FLOPs as if it were settled.
\item Unable to name a single case where encoder-decoder still wins.
\end{itemize}
\lab{Then they ask}
\begin{itemize}
\item What's a prefix LM, and does it get the best of both?
\item Why is Whisper encoder-decoder while LLMs aren't?
\item Could a decoder-only model ever match encoder-decoder at translation?
\end{itemize}
\vspace{2pt}\noindent{\color{muted}\rule{\linewidth}{0.3pt}}
\needspace{6\baselineskip}
\qhead{L6}{Trade-off}{Your 8K-trained model must serve 64K contexts next month. What are your options, and what do they cost?}
\lab{Say this}
Five options, and the honest answer ships two of them in parallel.\par\vspace{3pt}
Training-free RoPE rescaling --- dynamic NTK or dual-chunk attention --- is shippable this week: change the position mapping at inference so no rotary angle leaves the trained range. Zero training cost, quality degrades toward the far end of the window, and 8x is about the edge of what training-free methods deliver. A flagged stopgap, not the endpoint. What actually fits a month is YaRN plus a short fine-tune on long sequences --- a few hundred million to a few billion tokens, a rounding error against pretraining. Two risks: fine-tuning only on long data regresses short-context quality, so mix lengths; and perplexity can stay flat while in-context retrieval fails, so gate the launch on needle-in-a-haystack plus a RULER-style suite at 64K, never on loss curves.\par\vspace{3pt}
Continued pretraining with progressive length staging is the robust answer and a roadmap item, not a one-month plan. And two non-model options belong on the table: if ``64K support'' really means ``answer questions over long documents,'' RAG inside the existing window does it with no model work; or just serve a model already trained to 64K for that traffic slice.\par\vspace{3pt}
The serving bill arrives regardless, and I'd do this arithmetic unprompted. KV per request goes up 8x --- a 70B GQA model at 64K holds roughly 20 gigabytes against 2.5 at 8K, so concurrency per GPU collapses. Prefill's quadratic term grows about 64x, so chunked prefill and prefix caching stop being optional.\par\vspace{3pt}
\lab{Numbers}
\begin{itemize}
\item 70B GQA at 64K: \textasciitilde{}20 GB KV per request vs \textasciitilde{}2.5 GB at 8K
\item Prefill quadratic term grows \textasciitilde{}64$\times$ going 8K $\to$ 64K
\end{itemize}
\flab{Loses the signal}
\begin{itemize}
\item ``Set \texttt{max\_position\_embeddings} to 65536 and redeploy'' --- positions past 8K are out-of-distribution rotary angles.
\item Quoting perplexity as evidence the extension worked.
\item No mention of serving cost; the KV and prefill bill is usually the binding constraint, not quality.
\end{itemize}
\lab{Then they ask}
\begin{itemize}
\item Why does NTK-aware scaling preserve short-range quality better than position interpolation?
\item How much long data does the continued-pretraining stage need?
\item What breaks first if you serve 64K with the KV cache at INT4?
\end{itemize}
\vspace{2pt}\noindent{\color{muted}\rule{\linewidth}{0.3pt}}
\needspace{6\baselineskip}
\qhead{L6}{Estimation}{Estimate the cost per query for a RAG system vs a long-context LLM (128K).}
\lab{Say this}
Long context runs roughly thirty times the cost per query, and caching closes most of that gap. I'd use illustrative prices and say so: ten dollars per million input tokens, thirty per million output, cached input at a tenth of fresh. Prices fall every year; the durable point is that input tokens dominate.\par\vspace{3pt}
Stuffing a 128K window: 128,000 input tokens plus a short query, 500 tokens out --- a dollar twenty-eight in, a cent and a half out, call it a dollar thirty. RAG retrieves five 512-token chunks, so about 2,800 input tokens: under three cents of model cost plus a fraction of a cent of retrieval infrastructure. Call it four and a half cents. At ten thousand queries a day that's thirteen thousand dollars against four hundred and fifty --- about four point six million a year, which is the number that ends the meeting.\par\vspace{3pt}
Latency favours RAG where users feel it: retrieval and reranking cost around fifty milliseconds, while a 128K prefill is compute-bound and real APIs commonly show ten to thirty seconds to first token.\par\vspace{3pt}
Then the caching correction, because any caching argument that ignores the cached-input price is hand-waving. Over the same documents, turn one pays full price and writes the cache; later turns pay a tenth and skip prefill. Amortized over ten turns you're near twenty-five cents --- a gap of five or six times, close enough that cross-document reasoning quality decides rather than cost.\par\vspace{3pt}
\lab{Numbers}
\begin{itemize}
\item Illustrative: 128K long-context \textasciitilde{}$1.30/query vs RAG ~$0.045/query
\item At 10K queries/day: \textasciitilde{}\$4.6M/year difference
\item Cached input at \textasciitilde{}0.1$\times$: 10-turn session $\approx$ \$0.25/turn, gap narrows to \textasciitilde{}5--6$\times$
\end{itemize}
\flab{Loses the signal}
\begin{itemize}
\item Pricing only model inference and forgetting RAG's embedding, vector DB and reranking costs.
\item Ignoring prompt caching, which reprices the whole long-context arm for multi-turn sessions.
\item Not converting to an annual dollar figure --- the point of the estimate is the business decision.
\end{itemize}
\lab{Then they ask}
\begin{itemize}
\item How does caching change this for multi-turn conversations?
\item When is the long-context premium actually justified?
\item How would the answer change if answers were 5,000 tokens instead of 500?
\end{itemize}
\vspace{2pt}\noindent{\color{muted}\rule{\linewidth}{0.3pt}}
\needspace{6\baselineskip}
\qhead{L5}{Conceptual}{Walk me through what happens between ``we have a Common Crawl snapshot'' and the first training step.}
\lab{Say this}
In pipeline order, and the ordering is most of the answer. Start from WARC, the raw HTML, not the convenience WET text --- re-extract main content with a trafilatura-class extractor. That's the single highest-leverage decision, because boilerplate that survives extraction poisons everything downstream, and FineWeb's ablations showed the extractor choice mattered more than most of the filters people obsess over.\par\vspace{3pt}
Then language ID, routing into per-language pipelines with per-language thresholds. Then cheap heuristic filtering: URL and domain blocklists, then C4- and Gopher-style rules on length, symbol ratio, repetition, stop words. Interpretable, and it removes the bulk of obvious junk so the expensive stages are affordable at all.\par\vspace{3pt}
Then model-based quality filtering --- KenLM perplexity buckets and a fastText-class classifier trained on a small seed set, the DCLM and FineWeb-Edu pattern. The LLM only ever labels that seed set; you never run one over the corpus. Then dedup, exact hashing then MinHash near-dup, scoped per snapshot. Then decontamination by n-gram overlap against every benchmark you intend to report --- after dedup, so no surviving copy of a contaminated page slips through. Then mixture assembly, weights validated by proxy-run ablations. Then freeze the tokenizer, tokenize offline into fixed-width shards with checksums and version stamps, shuffle globally, fix a deterministic data order.\par\vspace{3pt}
The one-line version interviewers want: cheap filters before expensive ones, dedup before decontamination, quality defined by ablation, and everything frozen and versioned before the fleet spins up.\par\vspace{3pt}
\flab{Loses the signal}
\begin{itemize}
\item Starting at ``download a dataset from Hugging Face'' --- no awareness of what's in the corpus or how it got there.
\item Wrong ordering: decontamination before dedup, or expensive model filters before cheap heuristics.
\item No extraction stage at all, assuming Common Crawl is already text.
\end{itemize}
\lab{Then they ask}
\begin{itemize}
\item Why is WET a trap?
\item Where does an LLM legitimately appear in this pipeline?
\item What breaks if you skip dedup?
\end{itemize}
\vspace{2pt}\noindent{\color{muted}\rule{\linewidth}{0.3pt}}
\needspace{6\baselineskip}
\qhead{L6}{Trade-off}{How would you decide the code:web:math mixture for a pretraining run?}
\lab{Say this}
I'd anchor on a public reference and then describe what would move me off it. Llama 3 landed around fifty percent general knowledge, twenty-five percent math and reasoning, seventeen percent code, eight percent multilingual --- so code plus math is over forty percent, wildly above their share of raw web text, reflecting an ablation-supported belief that both transfer to general reasoning.\par\vspace{3pt}
The method has four parts. Feasibility first: multiply each domain's unique-token supply by a sane epoch cap --- roughly four repeats free, sixteen is the wall --- and that bounds each weight before you optimize anything. High-quality math is almost always the scarcest pool and the binding constraint, which is why synthetic math budgets exist. Then search: proxy runs at fixed compute over a mixture grid, or cheaper, many small randomized-mixture runs with a regression from weights to eval scores, RegMix-style. Validate at two proxy scales, because the optimum is scale-dependent --- what you rely on is candidate ranking staying stable across scale, not the optimum transferring. And exploit staging: a marginal math token is worth more in the anneal than the main phase, so part of the question is when, not just how much.\par\vspace{3pt}
The trade-off to name out loud: more code and math buys reasoning and tool use at the cost of world-knowledge breadth at fixed compute, and your eval suite sets the exchange rate. A suite over-indexed on GSM8K will happily push you into over-weighting math --- that's a Goodhart risk, not a discovery.\par\vspace{3pt}
\lab{Numbers}
\begin{itemize}
\item Llama 3 mix: \textasciitilde{}50\% general knowledge, 25\% math/reasoning, 17\% code, 8\% multilingual
\item Repeat budget: \textasciitilde{}4 epochs roughly free, \textasciitilde{}16 is the wall
\end{itemize}
\flab{Loses the signal}
\begin{itemize}
\item Quoting a ratio as a universal constant with no method for deriving or updating it.
\item Ignoring supply --- proposing 25\% math without asking whether that many unique math tokens exist.
\item Optimizing per-domain loss instead of downstream evals; low code loss can mean boilerplate, not capability.
\end{itemize}
\lab{Then they ask}
\begin{itemize}
\item You run out of unique math tokens --- now what?
\item Does code data actually help non-code tasks?
\item How does this change for a 1B edge model versus a frontier run?
\end{itemize}
\vspace{2pt}\noindent{\color{muted}\rule{\linewidth}{0.3pt}}
\needspace{6\baselineskip}
\qhead{L5}{Conceptual}{What makes an LLM system an ``agent''? Walk me through the agent loop and where the engineering difficulty actually lives.}
\lab{Say this}
An agent is a system where the model decides its own control flow at runtime. It reads its context, picks an action --- usually a tool call --- the environment executes it, the observation goes back into context, and it loops until the model produces a final answer or a budget stops it. The contrast class is a workflow, where code fixes the sequence at design time and LLM calls just fill slots. The distinguishing property is not ``uses tools'' --- a single forced function call isn't agency --- it's who decides what happens next.\par\vspace{3pt}
Three structural facts drive the engineering. The model is stateless: all state is the trajectory, replayed through the model every step, which makes context management, token cost and prefix caching first-order concerns rather than optimizations. Errors compound multiplicatively: per-step reliability p over n steps is roughly p to the n, and 0.98 to the fiftieth is about thirty-six percent --- so the difficulty lives in per-step reliability and recovery, not the loop's plumbing. And the action space is the tool set, which is why most agent failures trace to tool design --- overlapping tools, vague descriptions, bloated observations --- rather than model incapacity.\par\vspace{3pt}
I'd finish by placing the system on the agency spectrum: a workflow with LLM steps, a single tool loop, or autonomous multi-step with plans and sub-agents. The right design is the lowest level on that ladder that solves the task, because every rung up costs you predictability.\par\vspace{3pt}
\lab{Numbers}
\begin{itemize}
\item 0.98 per-step reliability over 50 steps $\approx$ 36\% task success
\end{itemize}
\flab{Loses the signal}
\begin{itemize}
\item Defining an agent as ``an LLM with tools'' --- that misses runtime control flow, which is the actual distinction.
\item Framework-first answers naming LangChain or AutoGen components with no grasp of the underlying loop.
\item No mention of termination or budgets; an agent that can't stop is a design bug, not an edge case.
\end{itemize}
\lab{Then they ask}
\begin{itemize}
\item Where does the ReAct pattern fit, and what replaced it?
\item Why is the trajectory replayed every step, and what does that cost?
\item What state lives outside the context window?
\end{itemize}
\vspace{2pt}\noindent{\color{muted}\rule{\linewidth}{0.3pt}}
\needspace{6\baselineskip}
\qhead{L6}{Trade-off}{When should you NOT build an agent? Your PM wants ``an agent'' for a document-processing product---how do you decide?}
\lab{Say this}
The test is whether the task decomposes statically. If I can write the steps down before seeing the input, and intermediate results don't change which steps run, then a deterministic pipeline with LLM calls in the slots beats an agent on cost, latency, debuggability and reliability. Document processing is usually exactly that: ingest, classify, extract fields, validate, route. Fixed DAG, LLM in the slots, code in charge.\par\vspace{3pt}
The argument I'd make to the PM is economic, not aesthetic. A pipeline runs a known number of calls with bounded context, so unit economics and p95 latency are predictable, each stage is unit-testable, and a failure localizes to one stage. An agent re-decides the workflow at runtime on every input: variable cost, variable latency, and every dynamic decision is a fresh opportunity to be wrong, with errors compounding across steps. You pay all that to buy adaptivity the task doesn't need. Debuggability settles it at scale --- pipeline failures reproduce, agent failures are points in a trajectory space where every run takes a different path.\par\vspace{3pt}
When the agent is right: the step structure is genuinely input-dependent and the number and identity of actions can't be known in advance --- debugging code, open-ended research, multi-turn support. And the honest answer is usually a hybrid: a static pipeline whose one genuinely dynamic stage, say resolving documents that fail validation, is a small budgeted agent loop. Spend agency exactly where the environment's response changes what to do next, and nowhere else.\par\vspace{3pt}
\flab{Loses the signal}
\begin{itemize}
\item Treating ``agent'' as the more advanced architecture rather than a cost you have to justify.
\item Hedging ``it depends'' without naming static decomposability as the actual test.
\item Ignoring debuggability and unit economics --- the two arguments that win this conversation with a PM.
\end{itemize}
\lab{Then they ask}
\begin{itemize}
\item What fraction of traffic would need the agent path to justify building it?
\item How would you migrate if the task later turns out to need adaptivity?
\item Your pipeline's classify stage is the bottleneck --- agent, or a better prompt?
\end{itemize}
\vspace{2pt}\noindent{\color{muted}\rule{\linewidth}{0.3pt}}
\needspace{6\baselineskip}
\qhead{L7}{Debugging}{Your agent succeeds on 60\% of an internal task suite. Leadership wants 95\%. Take it there.}
\lab{Say this}
I'd refuse to treat forty percent failures as one problem. The method: error taxonomy, biggest bucket first, cheapest adequate fix, every fix becomes a permanent regression case, repeat.\par\vspace{3pt}
The taxonomy comes from actually reading failed trajectories and labelling the first wrong step and its category --- wrong tool, bad arguments, misread observation, planning error, constraint lost after compaction, tool bug, grader bug where the eval marks correct behaviour wrong, and genuinely beyond the model. I'd also split consistent from flaky failures by rerunning with different seeds; a capability problem and a variance problem have different fixes.\par\vspace{3pt}
Then scaffold levers before model levers. Tool-interface confusion is usually the biggest and cheapest bucket: rewrite descriptions, merge overlapping tools, tighten schemas. Missing verification is the highest-leverage structural fix because it changes the arithmetic --- at 0.9 per-step success, a verifier catching ninety percent of failures with one retry lifts effective per-step success to about 0.98, orders of magnitude at task level. Lost constraints get compaction with a pinned constraint block; runaway trajectories get budgets and loop detection. Only then a stronger model on planning steps.\par\vspace{3pt}
And be honest about the last stretch. Sixty to eighty-five is scaffold engineering; eighty-five to ninety-five is auditing the residue --- under-specified tasks, and tasks needing capabilities the model lacks, which should escalate to a human. Ninety-five percent pass@1, ninety-five with verified retry, and ninety-five including graceful handoff are three different products, and committing to a number without defining the metric is the trap in the question.\par\vspace{3pt}
\lab{Numbers}
\begin{itemize}
\item p=0.9 per step + a verifier catching 90\% of failures with one retry $\to$ \textasciitilde{}0.98 effective per-step
\end{itemize}
\flab{Loses the signal}
\begin{itemize}
\item Reaching for better prompts or the next model as the first move, with no diagnosis.
\item Treating the 40\% as homogeneous and never reading actual trajectories.
\item No verification-and-retry, or no grasp of the compounding arithmetic that explains why it works.
\item Tuning against the full eval set and calling the resulting number progress.
\end{itemize}
\lab{Then they ask}
\begin{itemize}
\item The biggest bucket is ``model misreads tool output'' --- what do you change?
\item You get to 90\% and plateau --- now what?
\item How do you keep 95\% from regressing when the underlying model is upgraded?
\end{itemize}
\vspace{2pt}\noindent{\color{muted}\rule{\linewidth}{0.3pt}}
\needspace{6\baselineskip}
\qhead{L5}{Conceptual}{What is TF-IDF? Derive the formula. When would you use it over neural embeddings?}
\lab{Say this}
TF-IDF scores a term high when it's frequent in this document and rare across the corpus. Term frequency is the count of the term in the document, usually dampened as one plus log count so a term appearing fifty times doesn't count fifty times as much as appearing once. Inverse document frequency is the log of total documents over documents containing that term. Multiply the two and that's the weight of that term in that document.\par\vspace{3pt}
The log is doing real work. Document frequency spans orders of magnitude, so a raw ratio would let one rare term swamp everything; the log turns it into an additive weight on a sane scale. And the reason ``the'' scores near zero falls straight out --- it appears in nearly every document, so N over df is about one and its log is about zero. You also normalize by document length, or long documents win by default.\par\vspace{3pt}
I'd pick it over neural embeddings when I need interpretable features, because every dimension is a literal term; when latency matters and I don't want a GPU in the path; when there's no labeled data, since it needs only corpus statistics; or when the task is genuinely lexical --- keyword search, duplicate detection, matching identifiers. In production I'd reach for BM25 over raw TF-IDF and run it beside a dense retriever, because the failure modes are complementary: lexical misses paraphrase, dense misses exact rare terms.\par\vspace{3pt}
\flab{Loses the signal}
\begin{itemize}
\item ``TF-IDF is outdated and should never be used'' --- it's still the right tool for lexical, low-latency, unlabeled settings.
\item Not knowing the IDF formula, or why the log is there.
\item Forgetting sublinear TF and document-length normalization.
\end{itemize}
\lab{Then they ask}
\begin{itemize}
\item How does BM25 improve on TF-IDF?
\item What is the vocabulary mismatch problem?
\item How would you combine TF-IDF with neural retrieval in a hybrid system?
\end{itemize}
\vspace{2pt}\noindent{\color{muted}\rule{\linewidth}{0.3pt}}
\needspace{6\baselineskip}
\qhead{L5}{Conceptual}{Explain the difference between rule-based, statistical, and neural NLP. When would you still use rule-based methods?}
\lab{Say this}
Three eras, and the difference is what gets learned versus what gets written by hand. Rule-based systems have humans writing the patterns - regexes, grammars, gazetteers. Statistical NLP learns the parameters from data but humans still engineer the features: HMMs, CRFs, n-gram models, TF-IDF. Neural NLP learns features and parameters together, end to end, from Word2Vec and LSTMs through BERT and GPT.\par\vspace{3pt}
I would still reach for rules in about five situations, and I do regularly. When precision matters far more than recall and the target has a known format - dates, phone numbers, email addresses - a regex is exactly right and a model is the worse tool. When the domain is genuinely closed, like URL parsing or input validation. When a regulator or an auditor has to read the logic, which is common in content moderation. When speed is the constraint, since rule engines are orders of magnitude faster than any neural model. And when there is no training data at all and you have to encode domain knowledge directly.\par\vspace{3pt}
What I would actually say about production systems is that the good ones are hybrids: rules on the front for normalisation and the obvious cases, a neural model in the middle, rules again on the back as guardrails and post-processing. Framing it as rules versus neural is the wrong frame - the question is which layer each belongs in.\par\vspace{3pt}
\flab{Loses the signal}
\begin{itemize}
\item ``Rule-based methods are obsolete'' with no qualification.
\item Cannot give one concrete case where a regex beats a model.
\item Vague about what feature engineering meant in the statistical era.
\end{itemize}
\lab{Then they ask}
\begin{itemize}
\item Give an example of a hybrid rule-plus-model system you would build.
\item How do you test and maintain a large rule set?
\item When does a rule system get too complex to keep?
\end{itemize}
\vspace{2pt}\noindent{\color{muted}\rule{\linewidth}{0.3pt}}
\needspace{6\baselineskip}
\qhead{L5}{Conceptual}{Why use a CRF layer on top of BERT for NER instead of just a softmax classifier?}
\lab{Say this}
Because a softmax head makes each token's decision independently, and NER labels are a sequence with hard constraints between neighbours. BERT gives you rich input-side context --- every token already sees the whole sentence --- but nothing in a per-token softmax stops the model from emitting an impossible tag sequence like B-PER followed by I-ORG. Input-side dependency and output-side dependency are different things, and that distinction is really the whole answer.\par\vspace{3pt}
A CRF adds a learned transition matrix over label pairs, scoring how compatible adjacent tags are, and then decodes the best-scoring whole sequence with Viterbi rather than taking an argmax at each position independently. It learns from data that I-PER can follow B-PER or I-PER but not B-ORG, and that an I- tag can't open an entity. You train it end to end with the encoder; the likelihood needs a partition function over all possible sequences, which the forward algorithm gives you in time linear in sequence length.\par\vspace{3pt}
Where it earns its keep is low-data regimes, many entity types, and richer tag schemes like BIOES where there are more ways to be inconsistent. I'd also say the honest thing: with a large fine-tuned encoder and plenty of data the gain is often a few tenths of a point, and a constrained Viterbi decode with hand-written transition rules captures most of it without the extra training complexity.\par\vspace{3pt}
\flab{Loses the signal}
\begin{itemize}
\item ``BERT already models dependencies'' without separating input-side from output-side
\item Not knowing what the transition matrix does or how the sequence is decoded
\item Claiming the CRF always helps --- with a strong encoder and abundant data it's often marginal
\end{itemize}
\lab{Then they ask}
\begin{itemize}
\item How is the CRF trained end to end with BERT?
\item What is the partition function and how is it computed?
\item What is label bias, and does a CRF suffer from it?
\end{itemize}
\vspace{2pt}\noindent{\color{muted}\rule{\linewidth}{0.3pt}}
\needspace{6\baselineskip}
\qhead{L5}{Mathematical}{Explain the Word2Vec Skip-gram model. What is the training objective? What is negative sampling and why is it needed?}
\lab{Say this}
Skip-gram learns embeddings by predicting context words from a center word. Every word gets two vectors --- an input vector used when it's the center, an output vector used when it's context --- and that detail matters, because people forget it and then can't explain what gets shipped at the end, which is the input matrix. The objective is to maximize the log probability of the true context words, summed over every center-context pair in the corpus, where that probability is a softmax over the whole vocabulary.\par\vspace{3pt}
The problem is that denominator. It sums over the entire vocabulary --- a hundred thousand words or more --- for every single training pair, so every gradient step touches every output vector. That's the entire reason negative sampling exists.\par\vspace{3pt}
Negative sampling replaces the multi-class softmax with binary classification. For each real center-context pair you draw k noise words, and you train a logistic classifier to call the observed pair real and the noise pairs fake. That takes the per-example cost from order-vocabulary down to order-k. The one non-obvious piece is the noise distribution: it isn't the raw unigram distribution, it's unigram raised to the three-quarters power. That flattens it, so you draw rare words as negatives more often than their frequency warrants and get more informative gradients than if you sampled ``the'' over and over.\par\vspace{3pt}
\lab{Numbers}
\begin{itemize}
\item Negative sampling: O(V) $\to$ O(k) per example
\item Noise distribution: unigram raised to the 3/4 power
\end{itemize}
\flab{Loses the signal}
\begin{itemize}
\item Not knowing Skip-gram has two embedding matrices, input and output.
\item ``Negative sampling just samples wrong words'' without the binary-classification reformulation.
\item Confusing Skip-gram with CBOW --- predict context from center, not the reverse.
\end{itemize}
\lab{Then they ask}
\begin{itemize}
\item What's the connection between Word2Vec and PMI?
\item How do you choose the number of negative samples?
\item What is subsampling of frequent words and why does it help?
\end{itemize}
\vspace{2pt}\noindent{\color{muted}\rule{\linewidth}{0.3pt}}
\needspace{6\baselineskip}
\qhead{L5}{Conceptual}{What are the differences between Word2Vec, GloVe and FastText? When would you choose each one?}
\lab{Say this}
All three give you one static vector per word. They differ in what they train on and whether they can handle a word they've never seen.\par\vspace{3pt}
Word2Vec trains a shallow network on local context windows - skip-gram predicts context from the center word, CBOW does the reverse. It streams, so it scales to enormous corpora, but words are atomic: an unseen word has no vector at all.\par\vspace{3pt}
GloVe goes the other way. It builds the global co-occurrence matrix first, then fits vectors so dot products approximate log co-occurrence counts, using corpus-wide statistics directly rather than hoping local windows aggregate to them. The cost is materializing that matrix, and it's equally helpless on out-of-vocabulary words.\par\vspace{3pt}
FastText is Word2Vec plus character n-grams: a word's vector is the sum of its subword n-gram vectors. That gives you OOV handling for free and captures morphology, which matters a great deal in inflected languages.\par\vspace{3pt}
So FastText for typos, user-generated text, morphologically rich languages, or heavy domain vocabulary; GloVe if I just want good pretrained English vectors off the shelf. But the bigger decision is whether static embeddings are right at all - a contextual model beats all three when you can afford it, and static vectors earn their place on latency, memory and edge deployment.\par\vspace{3pt}
\flab{Loses the signal}
\begin{itemize}
\item ``They're all the same, just use BERT'' - dodges the question and misses where static vectors still win.
\item Doesn't know GloVe uses global co-occurrence counts while Word2Vec uses local windows.
\item Doesn't know FastText gets OOV handling from character n-grams.
\end{itemize}
\lab{Then they ask}
\begin{itemize}
\item How does Levy and Goldberg's result connect Word2Vec to SVD on a PMI matrix?
\item When would you use static embeddings in a system today?
\item How is FastText's n-gram decomposition different from BPE?
\end{itemize}
\vspace{2pt}\noindent{\color{muted}\rule{\linewidth}{0.3pt}}
\needspace{6\baselineskip}
\qhead{L5}{Conceptual}{Why are contextual embeddings better than static embeddings? Give a concrete example of when static embeddings fail.}
\lab{Say this}
Because a static embedding has to give one vector to a word that means several different things, and that vector is a compromise that is wrong for every sense.\par\vspace{3pt}
The clean example is ``crane''. The crane lifted the steel beam - construction equipment. We saw a crane at the lake - a bird. She had to crane her neck - a verb meaning to stretch. Word2Vec or GloVe puts all three at a single point, somewhere in the average of machinery, wildlife and body movement, which is a location that describes none of them. A contextual model gives you three different vectors, and it does word-sense disambiguation for free, without anyone building a sense inventory.\par\vspace{3pt}
Polysemy is the headline but not the whole story. Contextual models also encode syntactic role - ``running'' as a noun and as a verb come out differently - and sentence-level effects like negation and quantification propagate into the token representations, which static vectors have no mechanism for at all.\par\vspace{3pt}
The empirical jump was large: BERT moved GLUE up by roughly seven points over the previous state of the art, which is not a normal increment on that benchmark.\par\vspace{3pt}
I would not say static embeddings are dead, though. They are extremely cheap, they make honest baselines, and they are still fine inside larger systems where you need a fast lookup rather than a forward pass.\par\vspace{3pt}
\lab{Numbers}
\begin{itemize}
\item BERT: \textasciitilde{}7 points absolute on GLUE over the previous state of the art
\end{itemize}
\flab{Loses the signal}
\begin{itemize}
\item ``Contextual is better because BERT is better'' with no mechanism.
\item No concrete polysemy example when asked for one.
\item Claiming static embeddings are useless today.
\end{itemize}
\lab{Then they ask}
\begin{itemize}
\item What do ELMo's different layers capture?
\item How does BERT build context differently from ELMo?
\item Where would you still use static embeddings today?
\end{itemize}
\vspace{2pt}\noindent{\color{muted}\rule{\linewidth}{0.3pt}}
\needspace{6\baselineskip}
\qhead{L5}{Conceptual}{Explain BERT's pretraining objectives. Why was NSP removed in
RoBERTa?}
\lab{Say this}
BERT trained on two objectives at once. Masked language modelling selects 15 percent of tokens and predicts them from both directions; of those selected positions, 80 percent are replaced with the mask token, 10 percent with a random token, and 10 percent left unchanged. That 80/10/10 split exists because the mask token never appears at fine-tuning time --- if the model only ever had to reconstruct positions flagged with a special token, it would learn to do nothing at unflagged positions, so the random and unchanged cases force it to keep a useful representation everywhere.\par\vspace{3pt}
Next sentence prediction was the second objective: given two segments, decide whether they're consecutive. RoBERTa removed it and got consistently better results. The explanation is that it's trivially easy --- the negatives came from a different document, so topic overlap alone solves it and the model never has to learn discourse coherence. You spend capacity and part of every batch on a task with almost no signal.\par\vspace{3pt}
Worth adding that RoBERTa's other changes probably mattered more than dropping NSP: dynamic masking rather than one mask pattern baked in at preprocessing, much larger batches, and 160 gigabytes of text against BERT's 16. The real headline of that paper is that BERT was simply undertrained.\par\vspace{3pt}
\lab{Numbers}
\begin{itemize}
\item MLM: 15\% of tokens selected, 80/10/10 mask/random/unchanged --- only 12\% actually masked
\item RoBERTa: 160GB of training text vs BERT's 16GB
\end{itemize}
\flab{Loses the signal}
\begin{itemize}
\item Saying 15 percent of tokens are replaced with [MASK] --- 15 percent are selected and only 12 percent are actually masked
\item Can't explain why the 80/10/10 split exists
\item Doesn't know RoBERTa dropped NSP, or can't say why it was a weak objective
\end{itemize}
\lab{Then they ask}
\begin{itemize}
\item What is ELECTRA and why is it more sample-efficient?
\item Why does a 15 percent masking ratio work --- what happens if you mask 50?
\item What replaced sentence-level objectives in later models?
\end{itemize}
\vspace{2pt}\noindent{\color{muted}\rule{\linewidth}{0.3pt}}
\needspace{6\baselineskip}
\qhead{L5}{Conceptual}{Explain scaling laws for LLMs. What does the Chinchilla paper tell us about compute-optimal training?}
\lab{Say this}
Scaling laws are the empirical finding that language model loss falls as a power law in model size, dataset size and compute - no plateau in the ranges anyone has measured, and predictable enough that you fit the curve at small scale and extrapolate a run you haven't done yet. The bridge from a budget to a model is that training compute is about six N D: roughly six FLOPs per parameter per token, forward and backward.\par\vspace{3pt}
Chinchilla showed that for a fixed compute budget people were allocating it wrong. The optimum scales parameters and tokens together, roughly twenty tokens per parameter. GPT-3, at 175 billion parameters on 300 billion tokens, was badly undertrained; Chinchilla at 70 billion on 1.4 trillion matched it with a fraction of the parameters.\par\vspace{3pt}
Then the part where people get stuck: almost nobody trains Chinchilla-optimal anymore, and that's deliberate. Chinchilla minimizes loss for a fixed training budget and ignores inference entirely. If you're going to serve a model billions of times you want the smallest model that clears your quality bar, so you overtrain hard - Llama 3 8B on 15 trillion tokens is roughly ninety times the Chinchilla-optimal ratio. You pay the training cost once and collect cheaper serving forever.\par\vspace{3pt}
\lab{Numbers}
\begin{itemize}
\item Training compute C \textasciitilde{} 6ND FLOPs
\item Chinchilla: \textasciitilde{}20 tokens per parameter
\item Llama 3 8B: 15T tokens, \textasciitilde{}94x the Chinchilla-optimal token ratio
\end{itemize}
\flab{Loses the signal}
\begin{itemize}
\item Quotes 20:1 as a rule to obey rather than a compute-optimal point that ignores inference cost.
\item Doesn't know C \textasciitilde{} 6ND, so can't turn a compute budget into a model size.
\item Confuses compute-optimal with inference-optimal.
\end{itemize}
\lab{Then they ask}
\begin{itemize}
\item Estimate the FLOPs to train a 70B model on 2T tokens.
\item What happens when you run out of high-quality data?
\item Do scaling laws look the same for fine-tuning?
\end{itemize}
\vspace{2pt}\noindent{\color{muted}\rule{\linewidth}{0.3pt}}
\needspace{6\baselineskip}
\qhead{L5}{Conceptual}{Why does chain-of-thought prompting improve reasoning? What are its limitations?}
\lab{Say this}
Chain-of-thought works because the generated tokens are computation, not just output. A single forward pass gives the model a fixed amount of compute per token; letting it write intermediate steps lets it spend more passes on a hard problem and carry state in the text itself.\par\vspace{3pt}
Three mechanisms, concretely. Decomposition: a multi-step problem becomes several steps, each of which is inside what the model can do in one shot. Working memory: the reasoning tokens act as an external scratchpad, so the model does not have to hold everything in its hidden state. And distributional shift - pretraining data is full of worked step-by-step explanations, textbooks and tutorials and StackOverflow answers, so prompting for reasoning activates a mode the model already learned rather than teaching it something new.\par\vspace{3pt}
The limitations matter as much. It costs tokens, so latency and spend go up on every request. The chain can be unfaithful - the model produces reasoning that sounds right and lands on the wrong answer, or bad reasoning and the right answer anyway - which means you cannot treat the chain as an explanation of what the model actually did. It does not help on tasks with no sequential structure, like sentiment or NER, and can hurt by adding noise. And it needs scale: small models produce incoherent chains and get worse, not better.\par\vspace{3pt}
\flab{Loses the signal}
\begin{itemize}
\item ``It makes the model think'' with no mechanism.
\item Not raising unfaithfulness - treating the chain as a true explanation of the model's computation.
\item Claiming CoT always helps, including on simple or non-sequential tasks.
\end{itemize}
\lab{Then they ask}
\begin{itemize}
\item What is self-consistency and why does it help?
\item How does zero-shot CoT work?
\item When does CoT actively hurt?
\end{itemize}
\vspace{2pt}\noindent{\color{muted}\rule{\linewidth}{0.3pt}}
\needspace{6\baselineskip}
\qhead{L6}{Conceptual}{Explain in-context learning. Why can models learn from examples in the prompt without gradient updates?}
\lab{Say this}
In-context learning is a model doing a task from demonstrations in the prompt with no weight updates, and the honest answer is that we have three partial explanations rather than one settled one.\par\vspace{3pt}
The cleanest is mechanistic: induction heads. Olsson and colleagues identified a two-head circuit --- one head that carries information about the previous token, and one that looks back for an earlier occurrence of the current token and copies whatever followed it. That gives you ``A was followed by B earlier, so predict B again,'' which is literally pattern completion, and the circuit forms during training.\par\vspace{3pt}
The second framing is Bayesian: pretraining leaves the model with a prior over tasks, and the demonstrations don't teach a new mapping so much as identify which task it already knows. The third is distributional --- the pretraining corpus is full of implicit input-output pairs, Q\&A threads, translation pairs, docstrings next to code --- so continuing such a pattern is exactly what next-token prediction rewarded.\par\vspace{3pt}
There's also work arguing that attention during in-context learning produces something resembling a gradient step on the demonstrations --- implicit fine-tuning inside the forward pass. I'd cite it but flag that it's derived under linear-attention simplifications and how far it carries to real softmax attention is contested.\par\vspace{3pt}
Practically: format consistency matters most, example selection and ordering matter, and even random labels partly work --- which tells you the demonstrations mostly specify the task and the input distribution rather than the label mapping.\par\vspace{3pt}
\flab{Loses the signal}
\begin{itemize}
\item ``It just learns from the examples'' with no mechanism
\item Saying the model memorizes or updates on the examples --- no parameters change
\item Not knowing induction heads or any mechanistic account, and not noting that ICL needs scale
\end{itemize}
\lab{Then they ask}
\begin{itemize}
\item Why do random labels still partly work?
\item How does ICL performance scale with model size?
\item How would you select and order demonstrations in practice?
\end{itemize}
\vspace{2pt}\noindent{\color{muted}\rule{\linewidth}{0.3pt}}
\needspace{6\baselineskip}
\qhead{L5}{Trade-off}{Compare few-shot prompting, fine-tuning, and RAG for adapting
an LLM to a new task or domain. When would you use each?}
\lab{Say this}
It comes down to three things: how much labeled data you have, how often the knowledge changes, and what your per-query cost budget is.\par\vspace{3pt}
Few-shot prompting when you have under about fifty examples, the task is well-specified, or it changes frequently. No training, no infrastructure, you iterate in minutes. You pay at inference --- long prompts on every call --- and it gets inconsistent on genuinely complex tasks while eating context you might want for something else.\par\vspace{3pt}
Fine-tuning, full or LoRA, when you have a thousand-plus labeled examples and the task is stable. You get consistency, and you often get cheaper inference, because a fine-tuned small model with a short prompt beats a large model carrying a huge few-shot prompt on cost per call. The price is training infrastructure, catastrophic-forgetting risk, and a retraining treadmill as your data drifts.\par\vspace{3pt}
RAG when the knowledge changes faster than you can retrain, when you need grounding or citations, or when the corpus is simply too large to put in weights. It's the right answer whenever being factually current matters. The catch is that retrieval quality caps everything downstream --- a bad retriever gives you a confidently wrong answer with a citation attached --- and you now own an index and its freshness pipeline.\par\vspace{3pt}
In practice they compose, and I'd say so rather than picking one. Fine-tune for format and behavior, RAG for facts, few-shot for the last mile. That's the standard production shape.\par\vspace{3pt}
\flab{Loses the signal}
\begin{itemize}
\item Recommending one approach without the tradeoffs.
\item Not mentioning that they compose --- fine-tune for behavior, retrieve for facts.
\item Ignoring per-query cost and latency, which is often what actually decides it.
\end{itemize}
\lab{Then they ask}
\begin{itemize}
\item When would you combine fine-tuning with RAG?
\item What do you do in a domain with zero labeled data?
\item How would you evaluate which strategy is winning?
\end{itemize}
\vspace{2pt}\noindent{\color{muted}\rule{\linewidth}{0.3pt}}
\needspace{6\baselineskip}
\qhead{L5}{Conceptual}{How does an LLM actually call a tool? Walk me through the full round trip from tool definition to final answer.}
\lab{Say this}
Five stages, and the thing to get right up front is that the model never executes anything.\par\vspace{3pt}
First, the tool definitions go into the context: each one a JSON-Schema block with a name, a description and typed parameters, serialized by the model's chat template. Second, the model emits a call instead of prose. It's been fine-tuned on tool-use trajectories to produce a structured block --- tool name plus JSON arguments --- and it halts with a distinguished stop reason, \texttt{tool\_use} rather than \texttt{end\_turn}. Serving stacks typically layer constrained decoding against the schema on top, so the call is guaranteed to parse.\par\vspace{3pt}
Third, the application executes: it receives the call, validates the arguments, runs the function. This is the step candidates get wrong --- the model only emits text describing what it wants run. Fourth, the observation goes back, appended as a tool-result message keyed to that call's ID --- the result, or the error explicitly flagged as an error. Fifth, the model reads it and either answers or issues more calls. An agent loop is just this cycle iterated.\par\vspace{3pt}
Two details that signal depth: parallel tool calls, several independent calls in one turn executed concurrently to save round trips; and tool-choice forcing --- auto, required, a named tool, or none. Forcing a named tool whose input schema is your desired output schema is how structured extraction is actually implemented. And constrained decoding only guarantees syntax; picking the right tool with the right arguments is a model-quality problem no grammar fixes.\par\vspace{3pt}
\flab{Loses the signal}
\begin{itemize}
\item Saying or implying the model executes the function --- it emits a request, the application runs it.
\item Treating tool use as pure prompting, with no awareness it's trained in via SFT on trajectories and special formats.
\item Believing constrained decoding makes tool use correct --- it guarantees syntax, not tool or argument choice.
\item Not knowing tool errors go back to the model as observations rather than being swallowed or retried silently.
\end{itemize}
\lab{Then they ask}
\begin{itemize}
\item How would you get reliable structured JSON with no tool involved at all?
\item What breaks if your serving stack renders tools in a different format than the model was trained on?
\item How do you cap cost when a model gets stuck re-calling a failing tool?
\end{itemize}
\vspace{2pt}\noindent{\color{muted}\rule{\linewidth}{0.3pt}}
\chapter{Alignment \& Safety}
\addtocontents{toc}{\protect\thispagestyle{fancy}}
\needspace{6\baselineskip}
\qhead{L6}{First Principles}{Walk me through the DPO loss function. How does it eliminate the need for a separate reward model?}
\lab{Say this}
DPO trains on preference pairs with a plain supervised loss, and it works because the reward model was never strictly necessary - you can solve for it in closed form.\par\vspace{3pt}
The argument: the RLHF objective is maximize reward minus a KL penalty against a reference policy, and that has a known optimal solution, the reference reweighted by the exponential of reward over beta with a normalizing partition function. Invert it and reward becomes beta times the log ratio of policy to reference, plus a term depending only on the prompt. Push that into the Bradley-Terry preference model and the prompt-only term is identical on both sides, so it cancels - partition function and all. What's left is a logistic loss on the difference of log ratios between the chosen and the rejected response, trained exactly like supervised fine-tuning.\par\vspace{3pt}
So the reward model isn't gone, it's implicit: the log-probability ratio to the reference is the reward, and beta controls how far you're allowed to drift from it.\par\vspace{3pt}
What you gain is one fewer training stage and no PPO to tune. What you give up is a fixed offline preference set with no online exploration, and DPO refines a decent reference model rather than teaching from scratch. Thin preference data or a need for online collection, and RLHF still earns its cost.\par\vspace{3pt}
\flab{Loses the signal}
\begin{itemize}
\item Says DPO has no reward model at all; it has an implicit one, the log ratio to the reference.
\item Can't get from the RLHF objective to the loss, so can't say why the partition function cancels.
\item Claims DPO always beats RLHF.
\end{itemize}
\lab{Then they ask}
\begin{itemize}
\item What is beta doing intuitively?
\item When would you still run full RLHF?
\item What is KTO, and how does it differ?
\end{itemize}
\vspace{2pt}\noindent{\color{muted}\rule{\linewidth}{0.3pt}}
\needspace{6\baselineskip}
\qhead{L6}{Trade-off}{PPO vs DPO vs RLHF---complete comparison for LLM alignment.}
\lab{Say this}
The distinction that matters is online versus offline. PPO-based RLHF is real RL --- it generates fresh responses during training, scores them with a reward model, updates the policy. DPO is offline supervised learning on a static preference set: no generation loop, no reward model. And RLHF names the pipeline, not the optimizer, so I'd gently correct the framing.\par\vspace{3pt}
DPO wins on simplicity and cost: two models in memory instead of four, roughly 1.5 times SFT compute against three to five for PPO, and no RL loop to babysit. What I'd push back on is the folklore that DPO is safe from reward hacking. It has a different hacking profile, not a lower one --- its implicit reward gets over-optimized the same ways, length and style bias are well documented in DPO-trained models, and because training is offline you have no live RM-score-versus-KL curve to watch it happen.\par\vspace{3pt}
PPO's real advantages are exploration and reward flexibility. It can discover good outputs that aren't in any static dataset, and it accepts any scalar reward, including verifiable ones like unit tests passing or a proof checking.\par\vspace{3pt}
So: DPO when you have good static preference data and limited compute --- that's the Llama 3 choice. PPO when you have a verifiable signal or need the policy to explore. Frontier labs use both, often DPO for initial alignment and PPO with a process reward model for reasoning.\par\vspace{3pt}
\lab{Numbers}
\begin{itemize}
\item DPO: 2 models resident vs 4 for PPO
\item Training cost: DPO \textasciitilde{}1.5$\times$ SFT, PPO \textasciitilde{}3--5$\times$ SFT
\end{itemize}
\flab{Loses the signal}
\begin{itemize}
\item ``DPO is strictly better because it's simpler'' --- it cannot explore beyond its static dataset.
\item Confusing RLHF the pipeline with PPO the optimizer.
\item Asserting DPO is immune to reward hacking rather than differently exposed to it.
\end{itemize}
\lab{Then they ask}
\begin{itemize}
\item What is DPO's implicit reward model?
\item How would you run iterative DPO?
\item Why did Llama 3 go with DPO?
\end{itemize}
\vspace{2pt}\noindent{\color{muted}\rule{\linewidth}{0.3pt}}
\needspace{6\baselineskip}
\qhead{L5}{Conceptual}{What is reward hacking? Give 3 concrete examples and how to mitigate each.}
\lab{Say this}
Reward hacking is Goodhart's law inside RLHF: the reward model is a proxy for human preference, and once you optimize hard against a proxy the policy finds where the proxy is wrong rather than where the answer is better.\par\vspace{3pt}
Three that show up constantly. Verbosity: annotators equate length with effort, the RM learns it, the policy produces bloated repetitive answers. Fix it by length-normalizing the reward and seeding the RM training set with pairs where the shorter response wins. Sycophancy: the RM rewards agreement, so the model tells the user they're right that the earth is flat. Fix it with adversarial pairs where truthful disagreement is preferred, plus a separate truthfulness model blended into the reward. Third, distribution-shift exploitation: the policy drifts into regions the RM never saw and it extrapolates absurdly high scores. Tighter KL, and retrain the RM on the current policy's outputs.\par\vspace{3pt}
Detection is the part candidates skip. Track RM score, KL from the reference, and periodic human win-rate together. If RM score is climbing while human win-rate is flat or falling, you're hacking. Cheap secondary tells: response length inflating, vocabulary diversity collapsing, formatting locking into a template.\par\vspace{3pt}
And I'd frame it as managed, not solved. Any imperfect proxy under enough optimization pressure gets exploited.\par\vspace{3pt}
\flab{Loses the signal}
\begin{itemize}
\item Treating reward hacking as a bug to be fixed rather than a consequence of optimizing an imperfect proxy.
\item Offering a single mitigation --- the interviewer is looking for layered defenses plus detection.
\item Claiming DPO is immune; it still overfits biases baked into the preference data.
\end{itemize}
\lab{Then they ask}
\begin{itemize}
\item Is DPO immune to reward hacking?
\item How do process reward models help?
\item What's the relationship between the KL budget and reward hacking?
\end{itemize}
\vspace{2pt}\noindent{\color{muted}\rule{\linewidth}{0.3pt}}
\needspace{6\baselineskip}
\qhead{L6}{Debugging}{Your RLHF-tuned model becomes sycophantic (agrees with everything). Diagnose and fix.}
\lab{Say this}
Sycophancy is a specific reward hack, so I'd diagnose it as one rather than reaching for the architecture. Three hypotheses, each with a cheap test.\par\vspace{3pt}
First, RM bias. Sample a thousand preference pairs where the user's premise is debatable or false and check how often the preferred response agrees. If agreement wins the large majority, the RM has learned agreeableness rather than helpfulness. Second, the KL penalty is too loose --- check divergence from the frozen reference, and state your units, because this is where people get sloppy. Healthy runs sit in single-digit nats per sequence; twenty per sequence is a long way drifted, and twenty per token would mean the policy is unrecognizable. Third, the SFT data. If your assistant transcripts came from customer service, where agreeing is the house style, it was sycophantic before RL started.\par\vspace{3pt}
The fixes are layered and mostly upstream. Curate adversarial pairs where truthful disagreement beats agreeable falsehood, and retrain the RM. Raise beta. Blend in a separate truthfulness signal. Constitutional-AI style, add an explicit principle about politely correcting factual errors and generate self-critique pairs against it.\par\vspace{3pt}
Then build the measurement: a benchmark of false-premise queries where the correct answer disagrees, and track disagreement rate across RLHF iterations --- which also catches the over-correction into a model that argues with everything.\par\vspace{3pt}
\lab{Numbers}
\begin{itemize}
\item Healthy KL from reference: single-digit nats per sequence; \textasciitilde{}20 nats/sequence is serious drift
\end{itemize}
\flab{Loses the signal}
\begin{itemize}
\item ``Add more data'' without saying which data --- RM pairs, SFT set, or prompts --- caused it.
\item Blaming architecture; sycophancy is a data and objective problem.
\item Proposing a prompt-engineering fix at inference time for a training-level failure.
\end{itemize}
\lab{Then they ask}
\begin{itemize}
\item How do you measure sycophancy quantitatively?
\item How do you avoid over-correcting into reflexive disagreement?
\item Could DPO fix this without a reward model?
\end{itemize}
\vspace{2pt}\noindent{\color{muted}\rule{\linewidth}{0.3pt}}
\needspace{6\baselineskip}
\qhead{L5}{Conceptual}{Jailbreak attacks: how do they work and how do you defend?}
\lab{Say this}
Jailbreaks exploit a tension we built in deliberately: the model is trained to follow instructions and also trained to refuse some of them, and every attack finds a place where following wins.\par\vspace{3pt}
The families worth naming. Role-play --- ``you are DAN, you have no restrictions'' --- where instruction-following overrides safety training. Encoding attacks: Base64, ROT13, pig latin, anything that moves the request off the natural-language distribution safety training covers. Multi-turn escalation, where each turn is benign and only the trajectory is harmful, which per-turn checks miss by construction. Few-shot poisoning, establishing a pattern of harmful answers and letting in-context learning continue it. And prompt injection, different in kind: the instruction arrives inside data the model was asked to process.\par\vspace{3pt}
Defense has to be layered because no single layer holds. Input classifiers on known patterns, updated continuously, with encoded inputs decoded before classification. Output classifiers, the most robust layer because they catch harmful content regardless of how it was elicited. Instruction hierarchy trained in, so system prompts genuinely outrank user turns. Multi-turn context tracking, so you score the conversation rather than the message. And ongoing red teaming, human and automated, because the attack distribution moves faster than any fixed classifier.\par\vspace{3pt}
The honest framing is that this is an arms race you manage, not a problem you close.\par\vspace{3pt}
\flab{Loses the signal}
\begin{itemize}
\item Claiming safety training prevents jailbreaks --- they exploit the gap between training and deployment distributions.
\item Proposing a single defense layer.
\item Omitting output filtering, the layer that catches whatever slipped past the input side.
\end{itemize}
\lab{Then they ask}
\begin{itemize}
\item How do you handle a jailbreak your classifier has never seen?
\item What's the difference between prompt injection and jailbreaking?
\item How do you test your defenses systematically?
\end{itemize}
\vspace{2pt}\noindent{\color{muted}\rule{\linewidth}{0.3pt}}
\needspace{6\baselineskip}
\qhead{L5}{Conceptual}{How do you detect and reduce hallucination in a production LLM system?}
\lab{Say this}
Three layers --- prevent, detect, mitigate --- plus the honest framing that none of them closes the gap, because pretraining optimizes fluency, not factuality.\par\vspace{3pt}
Prevention is grounding first --- RAG over a curated corpus --- then scope: a narrow system prompt, domain fine-tuning, low temperature on factual queries, and a model trained to abstain rather than fabricate. Detection splits by type. Intrinsic, where the answer contradicts the retrieved context, is NLI entailment of each response sentence against that context. Extrinsic is self-consistency sampling and citation verification, with token-level uncertainty as a cheap auxiliary signal. Mitigation is what you do with a flag: citations and disclaimers on low-stakes traffic, human review on high-stakes, feedback routed back into the eval set.\par\vspace{3pt}
What makes that a production answer rather than a menu is three things the menu doesn't give you. Pick the operating point against review capacity, not a benchmark --- the detector is a rare-positive classifier, so its false-positive volume, not its precision, decides whether anyone acts on the flags. Instrument the rate: sample production traffic, run the verifier offline where the latency budget doesn't bind, and alert on the trend, because a hallucination rate you aren't measuring is one you'll hear about from users. And where the output drives tools or actions, the control that matters is permissions and gating, not text detection --- a detector on the final answer runs after the tool call already fired.\par\vspace{3pt}
\flab{Loses the signal}
\begin{itemize}
\item ``RAG solves hallucination'' --- the model can still hallucinate from retrieved context, or ignore it.
\item One technique instead of layered prevent/detect/mitigate.
\item Claiming RLHF eliminates it; sycophancy is itself a hallucination mode.
\end{itemize}
\lab{Then they ask}
\begin{itemize}
\item Intrinsic versus extrinsic hallucination --- what's the difference?
\item What do you do when the retrieved context is itself wrong?
\item Where does the detector's false-positive budget come from?
\end{itemize}
\vspace{2pt}\noindent{\color{muted}\rule{\linewidth}{0.3pt}}
\needspace{6\baselineskip}
\qhead{L5}{Conceptual}{Walk me through the RLHF pipeline step by step. What are the three stages?}
\lab{Say this}
Three stages: supervised fine-tuning, then a reward model, then RL against that reward.\par\vspace{3pt}
SFT first. You fine-tune the pretrained base on high-quality demonstration data, instruction and response pairs, so the model learns the format and register of a helpful answer. That checkpoint does double duty - it's the starting policy, and it becomes the reference the KL penalty is measured against later.\par\vspace{3pt}
Stage two, the reward model. For each prompt you sample two or more responses and have annotators say which is better, then train under Bradley-Terry - a logistic loss on the difference in scores between preferred and dispreferred - so the model only ever learns relative quality, never an absolute scale. Architecturally it's the same base model with the language modeling head swapped for a scalar head.\par\vspace{3pt}
Stage three, PPO. Optimize the policy to maximize expected reward minus beta times the KL from the reference. The KL term is not optional: without it the policy walks off into whatever degenerate text the reward model happens to score highly, which is reward hacking.\par\vspace{3pt}
The practical detail interviewers like: stage three holds four models in memory at once - policy, frozen reference, reward model, and the value function critic. That footprint is a big part of why DPO got popular.\par\vspace{3pt}
\lab{Numbers}
\begin{itemize}
\item PPO stage holds 4 models in memory: policy, reference, reward model, critic
\end{itemize}
\flab{Loses the signal}
\begin{itemize}
\item Gets the order wrong, or jumps from SFT straight to RL with no reward model.
\item Doesn't mention the KL penalty, or can't say what it prevents.
\item Confuses the reward model with the value function; different models, different jobs.
\end{itemize}
\lab{Then they ask}
\begin{itemize}
\item How many models are resident during PPO?
\item What happens if the reward model is weak?
\item Why is the SFT model the reference rather than the pretrained base?
\end{itemize}
\vspace{2pt}\noindent{\color{muted}\rule{\linewidth}{0.3pt}}
\needspace{6\baselineskip}
\qhead{L5}{Conceptual}{What is reward hacking? How does the KL penalty help? When does it fail?}
\lab{Say this}
Reward hacking is the policy learning to exploit the reward model instead of actually getting better. The reward model is a proxy trained on finite preference data, so it has blind spots, and PPO will find them, because finding the argmax of the reward is exactly what it is optimising. The concrete forms you see are verbosity, where the RM prefers longer answers so the policy pads; sycophancy, where it agrees with whatever the user asserted; and impressive-looking formatting with nothing inside it.\par\vspace{3pt}
The KL penalty against the reference policy is the main defence, and the logic is distributional. The reward model was trained on completions from the SFT model or something close to it, so its scores are trustworthy near that distribution and increasingly unreliable as you move away. Penalising KL keeps the policy inside the region where the proxy still tracks the thing you actually care about.\par\vspace{3pt}
It fails in three ways. If the reward model's bias is systematic rather than a narrow quirk - it just likes long responses everywhere - the policy can hack it without leaving the KL ball at all, because the hack is available nearby. If beta is set too low you have bought nothing. And the penalty anchors you to the reference model, so anything wrong with the SFT model is something the KL term actively preserves.\par\vspace{3pt}
\flab{Loses the signal}
\begin{itemize}
\item Defining reward hacking loosely as ``the model gets things wrong''.
\item Not connecting the KL penalty to where the reward model's scores are actually reliable.
\item Claiming the KL penalty solves reward hacking.
\end{itemize}
\lab{Then they ask}
\begin{itemize}
\item What does the over-optimisation curve look like?
\item How would you detect reward hacking in production?
\item How do reward-model ensembles help?
\end{itemize}
\vspace{2pt}\noindent{\color{muted}\rule{\linewidth}{0.3pt}}
\needspace{6\baselineskip}
\qhead{L6}{Mathematical}{Explain DPO. How does it relate to RLHF, and what are its limitations?}
\lab{Say this}
DPO turns RLHF into a single supervised loss by noticing that the reward model never had to be a separate artifact. The KL-constrained objective RLHF optimizes has a closed-form optimum: the optimal policy is the reference policy reweighted by the exponentiated reward. Invert that and the reward is beta times the log ratio of policy to reference, plus a term that depends only on the prompt. Substitute it into the Bradley-Terry preference model and that prompt term cancels, because you're comparing two responses to the same prompt. What's left is a loss on the policy itself --- push up the log ratio on the chosen response, push it down on the rejected one, with a sigmoid so the gradient fades once you're comfortably winning the pair.\par\vspace{3pt}
The practical win is fewer moving parts: two models instead of four, no reward model, no rollouts, no PPO. It trains like ordinary supervised learning, so it's far more stable and much cheaper.\par\vspace{3pt}
The limitations are real. It's offline --- you optimize against a fixed preference set with no exploration, so it can't discover responses better than what's in the data. There's no explicit reward model left, which you'd want for monitoring, for best-of-n sampling, or for reranking at inference. And it's sensitive to the reference policy: DPO from a poorly supervised-fine-tuned reference inherits that weakness.\par\vspace{3pt}
\lab{If they hand you the marker}
$r(x,y)=\beta\log\frac{\pi_\theta(y\mid x)}{\pi_{\text{ref}}(y\mid x)}+C(x)$ --- write this single line and say: the reward is implicit in the policy, and when you put it into Bradley-Terry the $C(x)$ cancels in the difference between chosen and rejected. That cancellation is the whole trick.\par
\flab{Loses the signal}
\begin{itemize}
\item Can't say where the prompt-dependent term cancels
\item Presenting DPO as just another RLHF variant, with no link to the KL-constrained optimum
\item Claiming DPO strictly dominates RLHF; not mentioning that it's offline
\end{itemize}
\lab{Then they ask}
\begin{itemize}
\item What is the implicit reward in DPO?
\item When would you still run full RLHF?
\item What does beta control in practice?
\end{itemize}
\vspace{2pt}\noindent{\color{muted}\rule{\linewidth}{0.3pt}}
\chapter{Evaluation}
\addtocontents{toc}{\protect\thispagestyle{fancy}}
\needspace{6\baselineskip}
\qhead{L6}{Debugging}{Our fraud model has 0.92 AUC, but the risk team says the probabilities are useless---their auto-decline rule fires at 0.9 and almost nothing crosses it. What is wrong, and how do you fix it?}
\lab{Say this}
AUC and calibration are different things, and 0.92 AUC tells you nothing about whether 0.9 means ninety percent. AUC only measures ordering: it is invariant to any strictly increasing transform of the scores, so it cannot distinguish p from a stretched or squashed version of p. The risk team is consuming the score as a number; the metric they were shown does not evaluate it as one.\par\vspace{3pt}
I would diagnose in order. Plot a reliability diagram on held-out data and report log loss or Brier alongside AUC. Then check the base rate - at a third of a percent fraud, a well-calibrated model should rarely exceed 0.9, so the threshold may be what is wrong, not the model. Then check for negative downsampling, the most common production cause: if negatives were kept with probability w, the scores sit on a shifted base rate and must be mapped back. A true one percent rate downsampled at w equal to 0.1 presents as 9.2 percent - a factor of nine, invisible to AUC. Then the model family - boosting pushes scores outward, random forests compress them toward the middle - and then drift.\par\vspace{3pt}
Then two independent fixes. Recalibrate with a monotone map fit on a held-out set: Platt under about a thousand points, isotonic above ten thousand. Both are monotone, so AUC barely moves. And re-derive the threshold from costs rather than a round number - five dollars to decline a good transaction against a hundred-and-twenty-dollar chargeback gives 0.04, not 0.9.\par\vspace{3pt}
\lab{Numbers}
\begin{itemize}
\item Cost-optimal threshold p* = $c_{FP}$/($c_{FP}$ + $c_{FN}$): $5 vs $120 gives 0.04
\item 1\% base rate downsampled at w = 0.1 presents as 9.2\%
\item Platt below \textasciitilde{}1K calibration points, isotonic above \textasciitilde{}10K
\end{itemize}
\flab{Loses the signal}
\begin{itemize}
\item ``AUC is 0.92, the model is fine'' - the exact conflation the question is built to catch.
\item Proposing a full retrain when a two-parameter post-hoc map on held-out data is the first move.
\item Fitting the calibrator on training-set predictions - for a tree ensemble those are near-perfect and the map comes out near identity.
\end{itemize}
\lab{Then they ask}
\begin{itemize}
\item Does isotonic regression change the AUC?
\item How much calibration data do you need?
\item Calibrated overall but wrong on one merchant segment - now what?
\end{itemize}
\vspace{2pt}\noindent{\color{muted}\rule{\linewidth}{0.3pt}}
\needspace{6\baselineskip}
\qhead{L5}{Conceptual}{When is random k-fold cross-validation invalid, and what do you use instead?}
\lab{Say this}
Random k-fold assumes rows are exchangeable and that validation stands in the same relationship to training that production will. Three things break that.\par\vspace{3pt}
Time structure. If the data is time-ordered, or deployment means train on the past and score the future, shuffled folds train on the future to predict the past and share period-level structure - a promotion, an outage, a fraud campaign - across the split. Use rolling-origin or expanding-window splits, and if labels mature, like a thirty-day churn label, add purging: drop training rows whose label window overlaps validation, plus an embargo gap.\par\vspace{3pt}
Repeated entities. If one user, patient, merchant or query contributes many rows, random folds put the same entity on both sides and you're crediting the model for memorizing identity. Use GroupKFold, or split on a hash of the entity ID, and StratifiedGroupKFold when classes are imbalanced too. Near-duplicate rows are the same failure with a different key, so deduplicate before splitting.\par\vspace{3pt}
Small or imbalanced classification doesn't invalidate random folds, but unstratified folds can end up with zero positives - so stratify.\par\vspace{3pt}
The rule underneath all of it: the split must simulate the deployment gap. Write down what production will actually ask - new entity or known one, future period or the same - and build the split to match. That also tells you when grouping is wrong: if you'll only ever score users you've already seen, a user-disjoint split understates production performance and makes you reject a good model.\par\vspace{3pt}
\flab{Loses the signal}
\begin{itemize}
\item Recites ``use TimeSeriesSplit for time series'' with no notion of why, or applies grouping reflexively without asking what production scores.
\item Misses group leakage entirely - the failure that has invalidated published results.
\item Believes stratification fixes temporal or group problems.
\end{itemize}
\lab{Then they ask}
\begin{itemize}
\item Your labels mature over thirty days - exactly how does that change the split?
\item The model scores both existing and brand-new users. Design the evaluation.
\item How would you detect group leakage if you didn't know the group key?
\end{itemize}
\vspace{2pt}\noindent{\color{muted}\rule{\linewidth}{0.3pt}}
\needspace{6\baselineskip}
\qhead{L6}{Estimation}{Model B is 1.5 points better than model A on MMLU---do you ship it?}
\lab{Say this}
Not on that alone. First, put an error bar on it. Full MMLU is about fourteen thousand questions; at seventy percent accuracy that's four tenths of a point of standard error per score, so the difference between two independently measured scores carries about 0.55 and a 95 percent interval of plus or minus 1.1. So 1.5 points is barely outside noise --- and if the eval ran on a thousand-question subset, common for cost, the interval is three or four points and the gap means nothing.\par\vspace{3pt}
Second, do the paired analysis, because both models answered the same questions. Count the flips --- what B fixed versus what B broke --- and run McNemar or a paired bootstrap. Pairing typically halves the interval and turns ``marginal'' into an actual answer. I'd also cluster standard errors by subject, since MMLU has 57 of them, because a gain concentrated in two subjects is a completely different claim than a uniform lift.\par\vspace{3pt}
Third, protocol. Same prompt template, few-shot count, answer-extraction parser, sampling settings? A 1.5-point gap is smaller than typical prompt-format sensitivity, so any asymmetry voids the comparison. And contamination: if B trained on a newer crawl, check n-gram overlap and whether the gain survives on a post-cutoff held-out set.\par\vspace{3pt}
Even if it's real, one saturated benchmark shouldn't gate a ship. I'd ship if the paired comparison holds, it replicates on internal held-outs, and the regression suite shows no losses on safety, latency or cost.\par\vspace{3pt}
\lab{Numbers}
\begin{itemize}
\item Full MMLU: 14,042 questions across 57 subjects
\item At \textasciitilde{}70\% accuracy: SE $\approx$ 0.4 pts per score; 95\% CI on the gap $\approx$ $\pm$1.1 pts
\item On a 1--2K-question subset: 95\% CI widens to $\approx$ $\pm$3--4 pts
\end{itemize}
\flab{Loses the signal}
\begin{itemize}
\item ``Yes, 1.5 > 0'' --- treating a benchmark score as a noiseless measurement.
\item Not asking how many questions the eval used, or whether the comparison was paired.
\item Ignoring contamination when the two models have different training cutoffs.
\end{itemize}
\lab{Then they ask}
\begin{itemize}
\item How does your answer change on GPQA-Diamond, which has 198 questions?
\item How would you put a confidence interval on pass@k?
\item What if B's training data was collected after MMLU was published?
\end{itemize}
\vspace{2pt}\noindent{\color{muted}\rule{\linewidth}{0.3pt}}
\needspace{6\baselineskip}
\qhead{L5}{Conceptual}{Your model has 95\% accuracy but stakeholders are unhappy. What is going on?}
\lab{Say this}
Almost always one of two things: the class balance makes 95 percent meaningless, or accuracy is not the metric they care about. Start with balance, the fastest check. If 95 percent of examples are negative, a model that always predicts ``negative'' scores 95 percent and is worth nothing. So the first thing I pull is the confusion matrix, before tuning anything, then move to precision, recall and PR-AUC, which don't flatter a majority-class predictor.\par\vspace{3pt}
Second, objective mismatch. Stakeholders in fraud or safety are usually saying ``we're missing cases,'' which is a recall complaint against a model optimized for overall correctness. That's a conversation rather than a modelling problem: get them to state the business cost of a false positive against a false negative, and pick the operating point from that.\par\vspace{3pt}
Third, the test set may not represent production --- temporal drift, demographic skew, or leakage inflating the number. Fourth, calibration: the model can rank correctly while emitting probabilities that are wrong as numbers, which silently breaks anything downstream that bids or thresholds on them.\par\vspace{3pt}
And the one I'd always check: slicing. Ninety-five percent overall can hide forty percent on the segment stakeholders actually look at --- a geography, a device type, new users. Aggregate metrics are exactly where this kind of unhappiness hides, so I slice by segment, geography and time before I believe any headline number.\par\vspace{3pt}
\flab{Loses the signal}
\begin{itemize}
\item Jumping to hyperparameter tuning or more data without questioning the metric.
\item Not asking about class balance --- this is a filter question.
\item Telling the stakeholders that accuracy is fine and they're wrong.
\end{itemize}
\lab{Then they ask}
\begin{itemize}
\item How would you choose between precision and recall here?
\item What metric would you propose instead?
\item What do you do when two stakeholders want conflicting metrics?
\end{itemize}
\vspace{2pt}\noindent{\color{muted}\rule{\linewidth}{0.3pt}}
\needspace{6\baselineskip}
\qhead{L5}{Conceptual}{How do you evaluate a search ranking system?}
\lab{Say this}
Three layers, and the interesting part is the bridge between them.\par\vspace{3pt}
Offline, NDCG@10 against human-judged graded relevance on a zero-to-four scale, because graded judgments distinguish ``somewhat relevant'' from ``exactly what I wanted'' and binary labels can't. MRR alongside it for navigational queries, and Recall@100 to evaluate retrieval separately from the ranker. The non-negotiable part is slicing: head versus tail, navigational versus informational, by user segment. It is very easy to improve head queries while quietly destroying the long tail.\par\vspace{3pt}
Online, an A/B test on session success --- click-through on the top results plus low reformulation and abandonment rates, which capture ``did the user get what they came for'' better than raw clicks. Guardrails on p95 latency and result diversity, and a long-horizon check on 7- and 30-day retention, because clickbait reliably raises CTR while lowering retention.\par\vspace{3pt}
The bridge is where I'd spend the rest of the answer. Both offline judgments and click logs are biased toward whatever the current model already ranks highly, so offline NDCG systematically flatters the incumbent. Interleaving fixes much of that: mix both rankers' results into a single page and see which side gets the clicks. It's roughly one to two orders of magnitude more sensitive per unit of traffic than a conventional A/B, so it catches real differences you could never afford to run an A/B long enough to see.\par\vspace{3pt}
\lab{Numbers}
\begin{itemize}
\item Interleaving: \textasciitilde{}10--100$\times$ more sensitive than a conventional A/B per unit of traffic
\end{itemize}
\flab{Loses the signal}
\begin{itemize}
\item Offline metrics only, with no online evaluation plan.
\item Using accuracy instead of ranking metrics.
\item Not knowing what interleaving is, or why offline judgments favour the incumbent.
\end{itemize}
\lab{Then they ask}
\begin{itemize}
\item How do you handle position bias in click data?
\item Walk me through an interleaving experiment.
\item How long do you run the A/B, and how do you decide that up front?
\end{itemize}
\vspace{2pt}\noindent{\color{muted}\rule{\linewidth}{0.3pt}}
\needspace{6\baselineskip}
\qhead{L5}{Trade-off}{Precision is 0.95 but recall is 0.30. What do you do?}
\lab{Say this}
First question back: is that a problem? For content moderation, where a false positive means wrongly punishing a real user, 0.95 precision at 0.30 recall may be exactly the operating point you chose on purpose. For fraud or medical screening, missing seventy percent of positives is a crisis. So I want the business cost of a false positive against a false negative before I touch anything.\par\vspace{3pt}
Assuming recall matters, the cheapest move is the threshold --- you're probably at 0.5 for no reason at all. Plot the full precision-recall curve and pick the point that minimizes expected cost, not the one that looks balanced. That costs nothing and often gets most of the way.\par\vspace{3pt}
If the threshold trade is too steep, the model is the problem, and this specific pattern --- high precision, low recall --- usually means it only learned the easy, obvious positives. Hard negative mining, focal loss or class weighting, oversampling rare positives, and above all better features for the subtle cases.\par\vspace{3pt}
Structurally, a cascade often beats a single threshold: keep the high-precision model as stage one for automatic action, add a higher-recall second stage over what it misses, and route the disagreements to human review. You buy recall without paying full precision. And audit the labels. Low recall sometimes means positives were labelled negative in the first place, especially when the label comes from something having been caught. Sample a hundred false negatives by hand --- some fraction usually aren't errors at all.\par\vspace{3pt}
\flab{Loses the signal}
\begin{itemize}
\item ``Just lower the threshold'' without first asking which error is expensive.
\item Not entertaining that 0.95/0.30 might be the correct operating point.
\item Never questioning label quality as a cause of apparent low recall.
\end{itemize}
\lab{Then they ask}
\begin{itemize}
\item How do you choose the optimal threshold?
\item What if stakeholders demand both above 0.90?
\item How does this change for a real-time system with a latency budget?
\end{itemize}
\vspace{2pt}\noindent{\color{muted}\rule{\linewidth}{0.3pt}}
\needspace{6\baselineskip}
\qhead{L5}{Conceptual}{How do you know if a model improvement is statistically significant?}
\lab{Say this}
Two different questions depending on whether you mean offline or online, and people conflate them.\par\vspace{3pt}
Offline on a fixed test set: compute the metric difference and bootstrap it --- resample with replacement ten thousand times, recompute the difference each time, take the 2.5th and 97.5th percentiles. If that interval excludes zero you have significance at 0.05. For classification, McNemar's is better still because it's paired: it looks only where the two models disagree, discarding shared variance and buying power from the same test set.\par\vspace{3pt}
Online, the real work happens before launch. Power analysis to size the test --- for something like a half-percent relative lift on a three percent baseline CTR you're into tens of millions of sessions per arm, and knowing that up front stops you from running a test that could never have detected anything. Randomize by user rather than by request, run at least one or two full weeks so day-of-week effects average out, and test on per-user aggregates rather than raw events. If you're reading five metrics, correct for it --- Benjamini-Hochberg rather than Bonferroni once you have many.\par\vspace{3pt}
And significance isn't the decision. A p of 0.01 on a 0.01 percent lift is a real effect that isn't worth the deployment risk. Set a minimum detectable effect from what the business would actually act on, and size to that.\par\vspace{3pt}
\lab{Numbers}
\begin{itemize}
\item 0.5\% relative lift on a 3\% CTR baseline: \textasciitilde{}20M sessions per arm at 80\% power
\item Bootstrap: \textasciitilde{}10K resamples for a 95\% CI on the metric difference
\end{itemize}
\flab{Loses the signal}
\begin{itemize}
\item Not distinguishing offline significance from online significance.
\item No sample-size planning --- running an underpowered test and calling it flat.
\item Confusing statistical significance with practical significance.
\end{itemize}
\lab{Then they ask}
\begin{itemize}
\item Statistical versus practical significance --- where do you draw the line?
\item How do you handle novelty effects in the first week?
\item What do you do if the test is still inconclusive after four weeks?
\end{itemize}
\vspace{2pt}\noindent{\color{muted}\rule{\linewidth}{0.3pt}}
\needspace{6\baselineskip}
\qhead{L6}{Debugging}{Your model's AUC is 0.95 but calibration is poor (ECE = 0.15). Why does this matter and how do you fix it?}
\lab{Say this}
AUC and calibration measure different things and you can have one without the other. AUC is discrimination --- can the model rank positives above negatives --- and it's invariant to any monotonic transform of the scores. Calibration asks whether a 0.7 actually means seventy percent. ECE of 0.15 says that, averaged over confidence bins, predicted probability and observed frequency differ by about fifteen points.\par\vspace{3pt}
Whether that matters depends entirely on how the score gets consumed. If you only rank --- search results, a recommendation feed --- AUC of 0.95 is fine and calibration is cosmetic. The moment anything multiplies or thresholds the probability, it's broken. Ad auctions multiply predicted CTR by a bid, so miscalibration is systematic mispricing of every impression. A medical risk score at a fixed threshold either over-treats or misses high-risk patients. A cascade that routes to an expensive model when confidence is low will route wrong.\par\vspace{3pt}
The fix is post-hoc and cheap, which is the key insight --- you don't retrain. Temperature scaling: learn one scalar on a held-out calibration set and divide the logits by it. It's monotonic, so AUC is unchanged by construction, and it nearly always works. That's the default. Platt scaling for more flexibility, isotonic regression if you have plenty of calibration data. To prevent it at the source, label smoothing and mixup both cut the overconfidence that causes it.\par\vspace{3pt}
\lab{Numbers}
\begin{itemize}
\item ECE 0.15 $\approx$ a 15-percentage-point average gap between confidence and observed accuracy
\end{itemize}
\flab{Loses the signal}
\begin{itemize}
\item ``AUC is 0.95 so the model is fine'' --- conflating discrimination with calibration.
\item Not knowing what ECE measures or how the bins work.
\item Proposing a retrain when temperature scaling on a held-out set is a one-parameter fix.
\end{itemize}
\lab{Then they ask}
\begin{itemize}
\item How do you calibrate a multi-class model?
\item Does temperature scaling change the AUC, and why?
\item How do you validate calibration when the calibration set is small?
\end{itemize}
\vspace{2pt}\noindent{\color{muted}\rule{\linewidth}{0.3pt}}
\needspace{6\baselineskip}
\qhead{L5}{Mathematical}{What is perplexity? How does it relate to cross-entropy?
What are its limitations as an evaluation metric?}
\lab{Say this}
Perplexity is the exponential of the cross-entropy loss --- the same quantity in different units. If you measure cross-entropy in bits, perplexity is two to that; in nats, e to that. A model at 3 nats has perplexity about twenty.\par\vspace{3pt}
The intuition that makes it useful is effective branching factor: perplexity twenty means the model is on average about as uncertain as if it were choosing uniformly among twenty next words. Lower means it put more probability on what actually came next. A uniform model over a vocabulary of size V has perplexity exactly V --- that's the do-nothing baseline.\par\vspace{3pt}
The limitations matter more than the definition, and this is where the question usually goes. It's tokenizer-dependent, so two models with different vocabularies aren't comparable at all; if you have to compare across tokenizers you use bits per byte instead. It only measures local next-token prediction, so it says nothing about global coherence or whether the content is true. It's domain-specific --- the same model can look excellent on news and terrible on clinical notes. And crucially it doesn't predict downstream capability: two models within noise of each other on perplexity can be far apart on QA or summarization, which is why nobody reports perplexity as a capability claim anymore. It's a training-health metric and a language-modeling comparison within a fixed tokenizer, and that's about it.\par\vspace{3pt}
\lab{Numbers}
\begin{itemize}
\item Uniform model over a vocabulary of size V has perplexity exactly V
\end{itemize}
\flab{Loses the signal}
\begin{itemize}
\item Not being able to state the relationship to cross-entropy.
\item Comparing perplexity across models with different tokenizers.
\item Claiming perplexity measures fluency, coherence, or factuality.
\end{itemize}
\lab{Then they ask}
\begin{itemize}
\item What's the perplexity of a uniform model over vocabulary V?
\item How do you compute perplexity for a subword model?
\item Why doesn't a frontier model's perplexity tell you how good it is at reasoning?
\end{itemize}
\vspace{2pt}\noindent{\color{muted}\rule{\linewidth}{0.3pt}}
\needspace{6\baselineskip}
\qhead{L5}{Conceptual}{How would you evaluate a RAG system end-to-end? What metrics would you track?}
\lab{Say this}
You have to decompose it, because a single end-to-end score tells you nothing about what to fix. Three layers plus an end-to-end number.\par\vspace{3pt}
Retrieval first, and recall@k is the one that matters most --- if the right chunk isn't in the context, nothing downstream can save you, so it's a hard ceiling on the whole system. NDCG and MRR on top of that when ordering matters, since position in the context window affects how much the model uses a chunk. This needs a labeled set of query-to-relevant-chunk pairs, and building that honestly is the real work.\par\vspace{3pt}
Generation second: given the context that was actually retrieved, is the answer relevant, complete and coherent? That's human evaluation or LLM-as-judge. Third, and this is the RAG-specific one, faithfulness: does the answer contain only what the retrieved context supports? You measure it with NLI-style entailment checks against the retrieved chunks, LLM-as-judge faithfulness scoring, or citation verification where every claim has to map to a cited span. RAGAS computes versions of these automatically. Then end-to-end answer correctness against a gold set, plus production signals --- thumbs up and down, CSAT.\par\vspace{3pt}
The reason to keep the layers separate is triage. Wrong answer with the right chunk retrieved is a generation or prompting problem. Wrong answer with the chunk missing entirely is an embedding or hybrid-search problem. Chunk retrieved but ranked twentieth is a reranking problem. Different fixes, and a blended score can't tell you which one you have.\par\vspace{3pt}
\flab{Loses the signal}
\begin{itemize}
\item Scoring only the final answer, with no decomposition into retrieval, generation and faithfulness.
\item Never mentioning faithfulness --- the metric that's specific to RAG.
\item ``We'll just use human evaluation'' with no concrete plan or labeled set.
\end{itemize}
\lab{Then they ask}
\begin{itemize}
\item How would you build the labeled evaluation set?
\item What are the biases of LLM-as-judge?
\item How do you detect a quality regression in production?
\end{itemize}
\vspace{2pt}\noindent{\color{muted}\rule{\linewidth}{0.3pt}}
\needspace{6\baselineskip}
\qhead{L5}{Mathematical}{What is BLEU? Derive it step by step. What are its limitations?}
\lab{Say this}
BLEU scores a machine translation by n-gram precision against one or more references, with two corrections that stop it being trivially gameable.\par\vspace{3pt}
The first is clipping. For each n-gram in the candidate you count its occurrences but cap that count at the most times it appears in any single reference. Without it you could output ``the the the the'' and score perfect unigram precision. The second is the brevity penalty. Precision alone rewards being short --- a one-word output that happens to match is a hundred percent precise --- so if the candidate is shorter than the reference you multiply by an exponential penalty; longer or equal, no penalty.\par\vspace{3pt}
You compute modified precision for unigrams through 4-grams and combine them with a uniform-weight geometric mean. Geometric matters: if any order scores zero the whole thing is zero, which is why sentence-level BLEU on short sentences is so unstable and why you report corpus-level.\par\vspace{3pt}
For limitations --- it's precision-only, with the brevity penalty as a crude stand-in for recall. No semantic awareness, so a good paraphrase gets punished. Poor correlation with human judgment at the segment level. Biased against morphologically rich languages. And it's capped by the number and quality of your references.\par\vspace{3pt}
In practice I'd report COMET, which is neural, trained on human judgments and correlates far better --- and still report BLEU, because reviewers expect it.\par\vspace{3pt}
\flab{Loses the signal}
\begin{itemize}
\item ``It's n-gram overlap'' and nothing more --- no clipping, no brevity penalty.
\item Not knowing why the brevity penalty exists.
\item Saying BLEU measures recall, or being unable to name a single limitation.
\end{itemize}
\lab{Then they ask}
\begin{itemize}
\item What is SacreBLEU and why did it need to exist?
\item How does BLEU compare with ROUGE?
\item How would you evaluate translation into a language like Finnish?
\end{itemize}
\vspace{2pt}\noindent{\color{muted}\rule{\linewidth}{0.3pt}}
\needspace{6\baselineskip}
\qhead{L5}{Mathematical}{Derive NDCG from first principles, then compute NDCG@3 for a ranking with grades (3, 0, 2) using exponential gain, where the ideal available grades are (3, 2, 0).}
\lab{Say this}
Build it in three moves, each fixing the previous one's flaw. Cumulative gain maps each relevance grade to a gain and sums them, but it ignores order entirely --- a good document at rank ten counts the same as at rank one. So discount by position: divide the gain at rank i by log base two of i plus one, modelling attention that decays down the page. That's DCG, and it now rewards putting good documents early. But its scale depends on how many good documents that query happens to have, so you can't average it across queries. Normalize by the ideal DCG --- the same judged documents sorted best-first --- and you get NDCG between zero and one, comparable and averageable.\par\vspace{3pt}
For this ranking with exponential gain, two-to-the-relevance minus one turns grades three, zero, two into gains seven, zero, three. The discounts at ranks one to three are one, about 0.63, and 0.5. So DCG at 3 is seven plus zero plus one and a half, which is 8.5. The ideal ordering is three, two, zero, giving gains seven, three, zero, so IDCG is seven plus three times 0.63, about 8.89. NDCG at 3 is 8.5 over 8.89, roughly 0.96.\par\vspace{3pt}
I'd state the conventions explicitly, because both are choices: exponential gain versus linear, and the log-of-i-plus-one discount, which is what keeps rank one undiscounted.\par\vspace{3pt}
\lab{Numbers}
\begin{itemize}
\item Discounts at ranks 1--3: 1, 0.631, 0.5
\item DCG@3 = 8.5, IDCG@3 $\approx$ 8.89, NDCG@3 $\approx$ 0.96
\end{itemize}
\lab{If they hand you the marker}
$$\mathrm{DCG@}k=\sum_{i=1}^{k}\frac{2^{\mathrm{rel}_i}-1}{\log_2(i+1)},\qquad \mathrm{NDCG@}k=\frac{\mathrm{DCG@}k}{\mathrm{IDCG@}k}$$ Write the discount as $\log_2(i+1)$ and say aloud that the $+1$ is what stops rank 1 from dividing by zero --- then plug the numbers in on the board rather than in your head.\par
\flab{Loses the signal}
\begin{itemize}
\item Skipping normalization, or normalizing by the shown list re-sorted instead of the judged ideal
\item Writing the discount as log$_2$(i), which divides rank 1 by zero
\item Getting the arithmetic right but unable to justify the log discount or the exponential gain
\end{itemize}
\lab{Then they ask}
\begin{itemize}
\item What if the grade-2 document were unjudged instead?
\item Why log rather than 1/i?
\item What happens to averaging when a query's judged pool has one relevant document?
\end{itemize}
\vspace{2pt}\noindent{\color{muted}\rule{\linewidth}{0.3pt}}
\needspace{6\baselineskip}
\qhead{L6}{Debugging}{Your reranker improved offline NDCG@10 by 3\% on the judgment set, but the A/B test shows flat CTR. Walk me through your investigation.}
\lab{Say this}
Three suspects: the offline number, whether the gain reaches users, and whether CTR is even the right question. The highest-yield one is the middle one, so I'd start there.\par\vspace{3pt}
What's the trigger rate? On what fraction of queries do the two systems actually produce a visibly different top ten? A 3 percent NDCG gain on 10 percent of queries is a 0.3 percent aggregate effect, almost certainly below the experiment's detection floor --- so the answer isn't a longer test, it's triggered analysis restricted to affected traffic. I'd also check where in the list the change lives, since reordering below the fold doesn't move clicks, and whether downstream business rules, dedup, or diversity constraints are clobbering the reranker before anything renders.\par\vspace{3pt}
On the offline number: where did the labels come from? If they're click-derived, the ``gain'' may just be increased agreement with the old ranker's position bias. And check unjudged documents --- if the new model surfaces things the pool never covered and they score as irrelevant, the comparison is rigged against it.\par\vspace{3pt}
On the measurement: CTR is an attractiveness metric. A real relevance win can be CTR-neutral while moving long-click or reformulation rate.\par\vspace{3pt}
Then I'd run interleaving, which is far more sensitive than an A/B. If interleaving prefers the new ranker and the A/B stays flat, the gain is real but small --- and that's a decision about the metric hierarchy, not about CTR.\par\vspace{3pt}
\lab{Numbers}
\begin{itemize}
\item 3\% NDCG on a 10\% trigger rate is a \textasciitilde{}0.3\% aggregate effect --- usually under the detection floor
\end{itemize}
\flab{Loses the signal}
\begin{itemize}
\item Never asking what fraction of queries actually changed --- the single highest-yield question here.
\item Concluding ``offline metrics are unreliable'' --- or the reverse, ``ship it, offline is what matters.''
\item Not knowing that click-derived judgments inherit the old ranker's position bias.
\end{itemize}
\lab{Then they ask}
\begin{itemize}
\item Interleaving prefers the new ranker but the A/B is still flat after four weeks --- ship or not?
\item How would you have caught this before spending A/B traffic?
\item What changes if the online metric is revenue per search rather than CTR?
\end{itemize}
\vspace{2pt}\noindent{\color{muted}\rule{\linewidth}{0.3pt}}
\chapter{Experimentation}
\addtocontents{toc}{\protect\thispagestyle{fancy}}
\needspace{6\baselineskip}
\qhead{L5}{Estimation}{We want to detect a 1\% relative lift on a 2\% click-through rate. How many users do we need, and how long should we run?}
\lab{Say this}
Convert to an absolute effect first, because that's where people lose four orders of magnitude. A 1 percent relative lift on a 2 percent base rate is an absolute delta of 0.0002. Then n is about sixteen times p times one-minus-p, over delta squared, per group --- the sixteen is two times 2.8 squared, where 2.8 is 1.96 plus 0.84, so that's 80 percent power at alpha 0.05. Plugging in: sixteen times 0.02 times 0.98 over 4e-8 gives about 7.8 million per arm, so 15.7 million users total. At two million eligible users a day that's roughly eight days, and I'd round up to two full weeks so the test covers whole weekly cycles instead of a weekday-skewed population.\par\vspace{3pt}
Then the things that separate levels. I'd reframe it as MDE rather than sample size --- given the traffic we have, what's the smallest effect we could detect? If that's larger than any plausible effect, don't run the test; decide on other grounds or bundle it into a bigger change. CUPED on the pre-period click rate, with a correlation around 0.7, cuts the required n roughly in half. Only triggered users count, so if the feature fires for 30 percent of traffic the calendar time triples. And CTR is a ratio metric --- with user-level randomization, compute variance by the delta method or a bootstrap over users, not over impressions, or your standard error will be far too small.\par\vspace{3pt}
\lab{Numbers}
\begin{itemize}
\item n $\approx$ 16p(1$-$p)/$\delta$$^2$ per arm at 80\% power, $\alpha$ = 0.05
\item 1\% relative on a 2\% base $\to$ $\delta$ = 0.0002 $\to$ \textasciitilde{}7.8M per arm, 15.7M total
\item CUPED at $\rho$ $\approx$ 0.7 roughly halves required n
\end{itemize}
\flab{Loses the signal}
\begin{itemize}
\item Using the relative lift as if it were absolute --- understates the requirement by four orders of magnitude
\item Quoting a sample size with no alpha, power or variance assumption attached
\item Sizing off total site traffic instead of the triggered population
\end{itemize}
\lab{Then they ask}
\begin{itemize}
\item How does the answer change at 95 percent power?
\item What if the metric is revenue per user instead?
\item What do you do when the MDE exceeds any plausible effect?
\end{itemize}
\vspace{2pt}\noindent{\color{muted}\rule{\linewidth}{0.3pt}}
\needspace{6\baselineskip}
\qhead{L6}{Trade-off}{Your test came back flat, but you are confident the feature helps. What do you do?}
\lab{Say this}
I'd push back on ``flat.'' Flat isn't ``no effect,'' it's ``no effect larger than my MDE,'' so the first move is to compute the confidence interval and ask what it excludes. A CI of minus 0.5 to plus 0.7 percent tells you nothing about the 0.3 percent effect you were hoping for. A CI of minus 0.05 to plus 0.06 genuinely rules out anything worth shipping for. Those are different situations, and everything downstream depends on which one you're in.\par\vspace{3pt}
If it was underpowered, buy power rather than looking harder. CUPED or stratification on pre-period behavior --- about a two-times effective sample saving at pre-post correlation 0.7. Restrict to the triggered population so untouched users stop diluting the effect. Pick a metric closer to the mechanism: the click this feature actually changes, not sitewide revenue. Extend in whole weeks, or raise allocation.\par\vspace{3pt}
If power was adequate and the interval is tight, take the result seriously. Verify the feature actually shipped --- instrumentation, exposure rates, an SRM check. Then check pre-registered segments and the effect-by-day curve for novelty or dilution. A win in an unregistered segment is a hypothesis for the next test, not a finding.\par\vspace{3pt}
Then decide, because that's what's really being asked. If the change is cheap or reduces complexity, ship it and say out loud that you're shipping without measured benefit. If it adds risk or maintenance cost, don't. Either way write down the MDE, so the next team knows what this test could and couldn't detect.\par\vspace{3pt}
\lab{Numbers}
\begin{itemize}
\item CUPED at pre-post correlation 0.7: \textasciitilde{}2x effective sample-size saving
\end{itemize}
\flab{Loses the signal}
\begin{itemize}
\item Peeking until it turns significant, or slicing segments until one does.
\item Treating p greater than 0.05 as proof of no effect.
\item Shipping anyway while claiming the experiment supported it.
\end{itemize}
\lab{Then they ask}
\begin{itemize}
\item How would you decide whether to ship on flat?
\item What do you tell a PM who wants the win?
\item When is sequential testing the right answer here?
\end{itemize}
\vspace{2pt}\noindent{\color{muted}\rule{\linewidth}{0.3pt}}
\needspace{6\baselineskip}
\qhead{L6}{Debugging}{Your model's online metrics diverge from offline metrics after 2 weeks in production. What is happening?}
\lab{Say this}
The two-week timescale is the clue. An overnight break is a bug; a gradual divergence over two weeks is the system feeding back on itself, or the world moving.\par\vspace{3pt}
I'd check the feedback loop first, because it's the cause that specifically worsens with time. The model's rankings decide what users see, users click what they're shown, and those clicks become the next training set --- position bias compounds into the labels and the model looks better to itself every week, while the offline set, frozen at launch, doesn't move. The tell is offline metrics stable while online engagement flattens. The fix is randomized exploration on a traffic slice, and unbiased learning-to-rank with inverse propensity weighting.\par\vspace{3pt}
Next, concept drift: the relationship between features and label shifted --- trends, seasonality, a competitor launch --- so the offline set is simply stale. That's distinct from data drift, where the inputs shifted: new demographics, a product change bringing in users the model never saw. Segmenting performance by cohort separates them. The unglamorous candidate is stale features: batch features fresh at launch are two weeks old if a refresh pipeline degraded quietly, and the model leans on outdated signal a little harder every day.\par\vspace{3pt}
The structural fix is to stop trusting a frozen offline set --- rolling evaluation windows, and off-policy estimates that account for how the data was collected.\par\vspace{3pt}
\flab{Loses the signal}
\begin{itemize}
\item Assuming it's a bug rather than a systemic property of a deployed model.
\item Never considering feedback loops or position-bias contamination of the labels.
\item ``Just retrain'' with no root-cause diagnosis.
\end{itemize}
\lab{Then they ask}
\begin{itemize}
\item How do you distinguish concept drift from data drift?
\item How do you build an offline eval that actually predicts online performance?
\item What is counterfactual evaluation and when would you use it?
\end{itemize}
\vspace{2pt}\noindent{\color{muted}\rule{\linewidth}{0.3pt}}
\needspace{6\baselineskip}
\qhead{L5}{Trade-off}{Shadow deployment vs A/B test vs multi-armed bandit---when do you use each?}
\lab{Say this}
They sit on a spectrum from zero risk to full optimization, and in practice you use them in sequence rather than picking one.\par\vspace{3pt}
Shadow deployment runs the new model on live traffic without serving its outputs. It answers ``does this break'' --- latency, resources, error rates, wild predictions --- before any user is exposed. The limitation is fundamental: nobody ever acts on those predictions, so you compare outputs but never outcomes. Shadow mode cannot tell you the model is better.\par\vspace{3pt}
A/B testing is the gold standard for that. Random assignment, real business metrics, clean causal inference. The price is speed and waste: you commit to a sample size up front, wait a week or two so day-of-week effects average out, and half your traffic sits on the worse variant the entire time.\par\vspace{3pt}
Bandits --- Thompson sampling or UCB --- shift traffic toward whatever is winning as evidence accumulates. Use them with many variants, say ten ranking strategies, and when cumulative regret matters more than a clean p-value. The cost is that adaptive allocation makes significance testing awkward and isolating any individual change much harder. So: shadow to prove nothing breaks, A/B to prove it helps, bandits afterwards for ongoing optimization across many variants. For ranking I'd slot interleaving in before the A/B --- far more sensitive per unit of traffic.\par\vspace{3pt}
\flab{Loses the signal}
\begin{itemize}
\item Treating bandits as strictly better than A/B tests.
\item Not knowing that shadow mode produces no user-facing outcome data at all.
\item Going straight to A/B for a risky change, skipping shadow.
\end{itemize}
\lab{Then they ask}
\begin{itemize}
\item When would you use interleaving instead?
\item How do you handle explore-exploit in a bandit with non-stationary traffic?
\item What is a switchback experiment and when do you need one?
\end{itemize}
\vspace{2pt}\noindent{\color{muted}\rule{\linewidth}{0.3pt}}
\chapter{Retrieval \& Ranking}
\addtocontents{toc}{\protect\thispagestyle{fancy}}
\needspace{6\baselineskip}
\qhead{L6}{Trade-off}{Pointwise vs.\textbackslash\{\} pairwise vs.\textbackslash\{\} listwise ranking losses---when do you use each?}
\lab{Say this}
Pick by what your labels look like and by what the score has to do downstream.\par\vspace{3pt}
Pointwise --- cross-entropy on click, MSE on a relevance grade --- treats every item independently and you get the ranking by sorting. Its real virtue is calibration: the score is a probability, which you need if it feeds an ad auction or a threshold. Its weakness is that it knows nothing about position, so an error at rank one costs the same as one at rank a thousand.\par\vspace{3pt}
Pairwise --- BPR, RankNet, margin loss --- optimizes relative order within a query. That's the natural fit for implicit feedback, where all you really have is ``the user clicked this one and not that one.'' The cost is that every pair is weighted equally, so a swap between one and two costs the same as ninety-nine and a hundred, and the pair count grows quadratically.\par\vspace{3pt}
Listwise --- LambdaRank, ListNet --- optimizes the whole list. LambdaRank scales each pair's gradient by the NDCG change that swap would cause, so top-of-list swaps dominate, which is the closest you get to optimizing your actual metric. The price is graded relevance labels and a harder debugging story.\par\vspace{3pt}
So: calibrated scores for bidding, pointwise. Binary implicit feedback, pairwise. Editorial graded labels with NDCG as the target, listwise. In a two-stage system I'd run InfoNCE or pairwise in the retriever and listwise or pointwise in the ranker.\par\vspace{3pt}
\flab{Loses the signal}
\begin{itemize}
\item ``Listwise is always better because it's more sophisticated'' --- it needs graded labels and adds real complexity.
\item ``Pointwise can't rank'' --- it ranks fine, it just doesn't optimize the ranking metric.
\item Never mentioning calibration, which is the whole reason to keep a pointwise head.
\end{itemize}
\lab{Then they ask}
\begin{itemize}
\item How does LambdaRank approximate NDCG optimization?
\item How do you handle position bias in the click data?
\item Can you combine pointwise and pairwise objectives in one model?
\end{itemize}
\vspace{2pt}\noindent{\color{muted}\rule{\linewidth}{0.3pt}}
\needspace{6\baselineskip}
\qhead{L5}{Debugging}{Your two-tower model retrieves documents that share keywords with the query but aren't actually relevant. How do you fix this?}
\lab{Say this}
The model has learned lexical matching dressed up as embeddings, and the cause is almost always the negatives you trained on. With random or in-batch negatives, the only thing the model has to do is separate a relevant document from a completely unrelated one --- keyword overlap solves that, and once the loss is low the gradient from those easy negatives is essentially zero. It never sees an example that shares keywords and isn't relevant, so it never learns the distinction you actually care about.\par\vspace{3pt}
I'd confirm it cheaply first. Take a hundred diverse queries, pull the top 50, label the false positives, and check whether they share salient query terms. Then score the same queries with BM25 --- if BM25 and the two-tower model rank nearly the same documents, you've paid a lot of money for BM25 in embedding space.\par\vspace{3pt}
Then fix it in order of impact. Mine hard negatives from BM25's top 100, drop anything that's actually a positive, and train against those; mix them with random negatives, roughly 30 percent hard, because training on hard negatives alone is unstable and amplifies label noise. Second, add a cross-encoder reranker over the top 100 --- a dual encoder structurally cannot model word-level interaction between query and document, and this is the most reliable fix in production. And start from a retrieval-pretrained encoder like E5 or BGE rather than vanilla BERT.\par\vspace{3pt}
\flab{Loses the signal}
\begin{itemize}
\item ``Make the model bigger'' --- the training signal is the problem, not capacity
\item Never mentioning hard negatives
\item Stripping keywords out of the query, which destroys information rather than fixing the model
\end{itemize}
\lab{Then they ask}
\begin{itemize}
\item How often do you refresh hard negatives during training?
\item What breaks if you train on hard negatives only?
\item How would you measure lexical versus semantic matching directly?
\end{itemize}
\vspace{2pt}\noindent{\color{muted}\rule{\linewidth}{0.3pt}}
\needspace{6\baselineskip}
\qhead{L5}{Trade-off}{Siamese vs Triplet vs Two-Tower---decision framework for a new retrieval project.}
\lab{Say this}
Scale and symmetry decide it, and for a new retrieval project the answer is usually two-tower.\par\vspace{3pt}
Siamese means one shared encoder for both sides. You want it when the two inputs are the same type and the comparison is symmetric --- sentence to sentence, face to face --- and especially when data is limited, because weight sharing is strong regularization. It holds up to maybe a hundred thousand items, where pairwise comparison is still tractable. Face and signature verification are the classic cases.\par\vspace{3pt}
Triplet is the same idea with anchor, positive, negative and a margin. You reach for it when you care about relative ranking rather than absolute similarity and you have enough examples per class to mine informative triplets. That mining is the whole game --- easy triplets give zero gradient, hard ones destabilize training, semi-hard is the sane default. Image retrieval and person re-identification live here.\par\vspace{3pt}
Two-tower uses separate encoders for query and item, and that asymmetry is what lets you precompute and index every item embedding offline. That's the unlock above a million items with a real-time latency budget, and it handles genuinely different input types --- a text query against a product record. You train it with InfoNCE, in-batch negatives plus mined hard negatives.\par\vspace{3pt}
Honestly, most modern systems skip straight to two-tower with InfoNCE, which subsumes both other objectives when you use in-batch negatives --- then add a cross-encoder reranker wherever precision beats latency.\par\vspace{3pt}
\flab{Loses the signal}
\begin{itemize}
\item ``Triplet networks are outdated'' --- still the right tool for fine-grained similarity with lots of data per class.
\item ``Two-tower is always better'' --- with limited data and symmetric same-type inputs, a shared encoder is more data-efficient.
\item Not naming corpus scale as the primary decision factor.
\end{itemize}
\lab{Then they ask}
\begin{itemize}
\item When would you run both a Siamese and a two-tower model in the same system?
\item How does the loss choice differ across the three?
\item What's the minimum training data you'd want for each?
\end{itemize}
\vspace{2pt}\noindent{\color{muted}\rule{\linewidth}{0.3pt}}
\needspace{6\baselineskip}
\qhead{L5}{Conceptual}{Cosine similarity vs dot product vs L2 distance for embeddings---when each?}
\lab{Say this}
For L2-normalized embeddings all three give you the same ranking, so the choice is about computation, not semantics. Cosine is the dot product after normalization, and squared L2 between unit vectors is two times one minus cosine - a monotone transform. So normalize at index time and use dot product, because that's what maximum inner product search and every ANN library is built for.\par\vspace{3pt}
The choice only matters when embeddings aren't normalized, and then it's a question about magnitude. Dot product rewards large norms, and norms correlate with frequency and popularity, because those items got more gradient updates. In a recommender that's often a feature - dot product naturally boosts popular items. In semantic search it's usually a bug, and cosine strips it out.\par\vspace{3pt}
L2 I reach for when I need an actual metric with the triangle inequality, so k-means or clustering on the embeddings, or when the space was trained with a distance-based loss like triplet or contrastive with Euclidean margins. Match the retrieval metric to the training objective - that's the rule that actually matters.\par\vspace{3pt}
CLIP and most modern encoders already normalize, so the question is moot there. It comes straight back the moment you concatenate unnormalized side features like price or recency.\par\vspace{3pt}
\lab{Numbers}
\begin{itemize}
\item For unit vectors: ||u - v||\textasciicircum{}2 = 2(1 - cos(u,v)), so all three rank identically
\end{itemize}
\flab{Loses the signal}
\begin{itemize}
\item ``Cosine is always right for embeddings'' - only if magnitude carries nothing you want.
\item Doesn't know the three are rank-equivalent for normalized vectors; that's the most commonly tested fact here.
\item Picks a metric without asking what loss the embeddings were trained with.
\end{itemize}
\lab{Then they ask}
\begin{itemize}
\item What is MIPS, and why is it the standard retrieval primitive?
\item How does the metric choice interact with your ANN index?
\item You want cosine semantics but dot-product speed - what do you do?
\end{itemize}
\vspace{2pt}\noindent{\color{muted}\rule{\linewidth}{0.3pt}}
\needspace{6\baselineskip}
\qhead{L5}{Conceptual}{Explain BM25. Why is it still used in production search systems?}
\lab{Say this}
BM25 is a bag-of-words scoring function that fixes the two places raw TF-IDF is naive, and it's still the default lexical retriever in Lucene, Elasticsearch and Solr.\par\vspace{3pt}
First is term-frequency saturation. In TF-IDF a document mentioning your query term twenty times scores twenty times one that mentions it once, which is obviously wrong. BM25 runs the count through a saturating function, controlled by k1, so the fifth occurrence adds far less than the second. Second is length normalization: the denominator scales with document length relative to the corpus average, controlled by b, so a document twice the average needs roughly twice the term frequency to score the same. That stops long documents winning by containing everything.\par\vspace{3pt}
It survives because it's an inverted-index lookup, so milliseconds over hundreds of millions of documents; it needs no training data, just corpus statistics, so it works on day one in a new domain; and it's interpretable, so you can say exactly which query terms matched and what each contributed.\par\vspace{3pt}
But the real reason is complementarity. Dense retrieval fails on exact matches - part numbers, rare proper nouns, code identifiers - which is precisely what lexical matching is good at, so every serious production system is hybrid.\par\vspace{3pt}
\lab{Numbers}
\begin{itemize}
\item Typical parameters: k1 around 1.2-2.0, b around 0.75
\end{itemize}
\flab{Loses the signal}
\begin{itemize}
\item ``BM25 is just TF-IDF'' without being able to name what changed.
\item Doesn't know what k1 and b control.
\item Treats it as legacy - it is the lexical half of essentially every production hybrid system.
\end{itemize}
\lab{Then they ask}
\begin{itemize}
\item What are typical values for k1 and b?
\item What is BM25F for?
\item How would you fuse BM25 with a dense retriever?
\end{itemize}
\vspace{2pt}\noindent{\color{muted}\rule{\linewidth}{0.3pt}}
\needspace{6\baselineskip}
\qhead{L6}{Debugging}{Your new embedding model wins offline --- +4 points Recall@100 on your evaluation set --- but loses the online A/B test. Walk me through your investigation.}
\lab{Say this}
I'd treat this as two measurements that disagree and enumerate what differs between them, cheapest and most likely first --- and the highest-prior cause usually isn't the model, it's latency. If a much larger encoder replaced a small one, twenty to forty milliseconds of extra query latency will cost more engagement than four points of Recall@100 buys. So I check per-arm latency distributions and guardrail metrics, and confirm the test is bucketed and powered, before touching relevance at all.\par\vspace{3pt}
Next, serving parity. Was the entire corpus re-embedded? A partially backfilled index fails silently and catastrophically, because vectors from two different models aren't comparable --- the un-migrated slice just scores as noise. Then check preprocessing: many of these models expect a query prefix or instruction, and if the offline eval had it and serving doesn't, the advantage is gone.\par\vspace{3pt}
Third, the ANN interaction. Offline you almost certainly scored with exact search, while production serves from HNSW or IVF-PQ tuned on the old model's geometry. Measure ANN-versus-exact overlap at 100 in both arms and retune the index before judging the model.\par\vspace{3pt}
Fourth, downstream coupling --- a reranker or click model trained on the old candidate distribution can mis-score unfamiliar candidates, converting better recall into worse final rankings. Only then would I question the plus-four itself: stale eval sets, and click labels carrying position and selection bias toward whatever the old system surfaced.\par\vspace{3pt}
\flab{Loses the signal}
\begin{itemize}
\item ``The A/B is underpowered, ship it anyway'' --- or the reverse, ``offline metrics are meaningless''
\item Not knowing that two models' embeddings can't share one index
\item Never considering latency, or assuming production search is exact rather than ANN
\end{itemize}
\lab{Then they ask}
\begin{itemize}
\item How would you retune HNSW for a new embedding distribution?
\item How do you get unbiased relevance labels out of click logs?
\item Retrain the reranker before or during the embedder A/B?
\end{itemize}
\vspace{2pt}\noindent{\color{muted}\rule{\linewidth}{0.3pt}}
\needspace{6\baselineskip}
\qhead{L5}{Trade-off}{Compare BM25 vs.\textbackslash\{\} dense retrieval. When would you use each? When
would you combine them?}
\lab{Say this}
In production, almost always both --- they fail in complementary ways, which is the whole argument.\par\vspace{3pt}
BM25 is lexical: an inverted index scoring documents on query-term overlap with TF-IDF-style weighting. No training data, extremely fast, interpretable, and unbeatable on exact strings --- entity names, part numbers, error codes, legal citations. Its failure is vocabulary mismatch: search ``automobile'' and miss the document that says ``car.''\par\vspace{3pt}
Dense retrieval embeds queries and documents into a shared space with trained encoders and does nearest-neighbor search. It handles paraphrase and conceptual queries, which is precisely BM25's blind spot. Its failures are the mirror image: it can miss a rare exact token entirely, it needs training data and GPU infrastructure, and when it's wrong you can't tell why.\par\vspace{3pt}
So BM25 alone when you have no training data or the domain is precise-terminology-heavy --- medical codes, statutes, SKUs. Dense alone when queries and documents genuinely use different vocabulary and you have data to train on. But the honest production answer is hybrid, fused with reciprocal rank fusion, and it consistently beats either leg because the errors don't correlate.\par\vspace{3pt}
The thing I'd add is the cost of that: hybrid doubles your operational surface --- two indexes, two freshness pipelines, two things that page you. So I'd want per-slice evaluation showing where each leg actually earns its keep, rather than a global average that says hybrid is better and tells you nothing about which queries need it.\par\vspace{3pt}
\flab{Loses the signal}
\begin{itemize}
\item Calling BM25 obsolete or legacy-only.
\item Can't explain the vocabulary mismatch problem concretely.
\item Never mentioning hybrid --- or mentioning it without acknowledging the cost of running two indexes.
\end{itemize}
\lab{Then they ask}
\begin{itemize}
\item What is SPLADE and how does it bridge sparse and dense?
\item How would you choose the RRF constant?
\item If you had to drop one leg, which and why?
\end{itemize}
\vspace{2pt}\noindent{\color{muted}\rule{\linewidth}{0.3pt}}
\needspace{6\baselineskip}
\qhead{L5}{Trade-off}{When should you use RAG vs. fine-tuning vs. long context?}
\lab{Say this}
They solve different problems, so the question is which one you need, and usually the answer is more than one. The clean split: RAG is for knowledge, fine-tuning is for behavior, long context is for a document you have in hand right now.\par\vspace{3pt}
RAG when the knowledge changes, when you need citations, when the corpus is far too big for any context window, or when someone has to audit why the model said what it said - legal, medical, compliance. Its real advantage is that you update it by re-indexing, with no training run at all.\par\vspace{3pt}
Fine-tuning when you want to change how the model behaves rather than what it knows: tone, output format, instruction-following style, domain reasoning patterns, or distilling a big model into a small one. Knowledge baked into weights can't be updated cheaply and carries no attribution, which is why using fine-tuning as a knowledge store disappoints.\par\vspace{3pt}
Long context when the corpus fits, when the task genuinely is ``analyze this PDF,'' or when pipeline simplicity beats efficiency. The limits are that cost scales with context length on every call, retrieval inside a long context degrades in the middle, and the window is still finite.\par\vspace{3pt}
In production it's usually fine-tune for behavior plus RAG for knowledge.\par\vspace{3pt}
\flab{Loses the signal}
\begin{itemize}
\item Treats them as three competing options rather than complementary layers.
\item Can't articulate the behavior-versus-knowledge line, so recommends fine-tuning to teach facts.
\item Ignores update cadence, per-call cost, and attribution requirements.
\end{itemize}
\lab{Then they ask}
\begin{itemize}
\item What is RAFT, retrieval-augmented fine-tuning?
\item How do the trade-offs change at million-token context windows?
\item Would you fine-tune the retriever or the generator?
\end{itemize}
\vspace{2pt}\noindent{\color{muted}\rule{\linewidth}{0.3pt}}
\needspace{6\baselineskip}
\qhead{L5}{Conceptual}{Explain the bi-encoder versus cross-encoder tradeoff in retrieval. How would you use both in a production system?}
\lab{Say this}
Bi-encoders are cheap and shallow, cross-encoders are expensive and deep, and in production you use both in a funnel.\par\vspace{3pt}
A bi-encoder encodes query and document independently into vectors and scores with a dot product. Because the documents are encoded offline and sit in an ANN index, search is sublinear in corpus size, so you can run it over hundreds of millions of documents. What you give up is any interaction between query and document tokens - the model has to commit to a document representation before it has seen the query.\par\vspace{3pt}
A cross-encoder concatenates query and document and runs them through the transformer together, so every query token can attend to every document token. That is dramatically better relevance, and it is why it wins every quality comparison. It is also one forward pass per document, so full-corpus search is linear and hopeless past a few hundred or a thousand candidates.\par\vspace{3pt}
So the standard pattern is two stages: bi-encoder, or hybrid BM25 plus dense, retrieves the top hundred, and a cross-encoder reranks those down to a final five. Reranking a hundred documents costs a couple of hundred milliseconds, which most applications absorb.\par\vspace{3pt}
The point I would emphasise is the evaluation consequence. Because the expensive model only sees a few candidates, the first stage is judged on recall at its handoff depth, never on precision - if the right answer is not in the hundred, no reranker recovers it.\par\vspace{3pt}
\lab{Numbers}
\begin{itemize}
\item Reranking top-100 with a cross-encoder: \textasciitilde{}200-500 ms
\end{itemize}
\flab{Loses the signal}
\begin{itemize}
\item Cannot say why cross-encoders win - no mention of joint encoding or cross-attention.
\item Proposing a cross-encoder over the whole corpus.
\item Evaluating the first stage on precision rather than recall at the handoff depth.
\end{itemize}
\lab{Then they ask}
\begin{itemize}
\item Where does ColBERT sit between the two?
\item How do you decide how many candidates to rerank?
\item How does cross-encoder to bi-encoder distillation work?
\end{itemize}
\vspace{2pt}\noindent{\color{muted}\rule{\linewidth}{0.3pt}}
\needspace{6\baselineskip}
\qhead{L6}{Debugging}{Your RAG system answers factual lookup questions well but fails on questions that require combining facts across documents --- ``which customers on the legacy plan have an open P1 ticket?'' or ``did the Q3 pricing change contradict the renewal terms?'' Retrieval looks healthy: every returned chunk is on topic. Diagnose it.}
\lab{Say this}
This looks like a reasoning failure but it's an evidence-set failure, and that distinction is the diagnosis. The retriever optimizes chunk-level similarity, so every chunk it returns is individually on topic. Nothing scores whether the set is sufficient to answer. ``A relevant chunk was retrieved'' and ``the answer is in the context'' are different claims, and recall@k only checks the first.\par\vspace{3pt}
There's a structural reason: a composed question embeds to roughly the average of its parts, landing near documents that discuss both topics generically rather than the two chunks that each answer one half. And for a genuine second hop the bridging entity isn't in the query at all, so no single-shot retrieval at any k reaches it.\par\vspace{3pt}
Confirm it before fixing: annotate a slice of reasoning queries with the complete evidence set and measure all-hops recall --- did we get every required chunk --- next to ordinary recall@k. The gap is the diagnosis. Then hand the generator the gold chunks; if it answers, the defect is retrieval completeness, and if it doesn't, you have a real synthesis limit.\par\vspace{3pt}
Fixes cheapest first: query decomposition for independent sub-questions, iterative retrieval with a hop cap when hop two depends on hop one, and retrieval units that match the question --- customer records, not paragraphs mentioning one. But ``which customers on the legacy plan have an open P1'' is a join with a filter; that belongs in SQL, and the mature system routes set-valued queries there.\par\vspace{3pt}
\flab{Loses the signal}
\begin{itemize}
\item Reaching for a stronger model or prompt tweaks when the facts were never in context
\item Citing recall@k as proof retrieval is healthy without asking whether the set was sufficient
\item Never considering that set-valued or aggregate queries belong in a database or knowledge graph
\end{itemize}
\lab{Then they ask}
\begin{itemize}
\item How would you build a multi-hop eval set without hand-annotating every hop?
\item What termination condition and budget would you give an iterative retrieval loop?
\item When do you tell the team a query class needs SQL rather than a better retriever?
\end{itemize}
\vspace{2pt}\noindent{\color{muted}\rule{\linewidth}{0.3pt}}
\needspace{6\baselineskip}
\qhead{L6}{Trade-off}{Lexical, dense, or hybrid --- argue the retrieval split for a marketplace search engine.}
\lab{Say this}
Hybrid --- but I'd argue it from the query mix rather than assert it as a default. Marketplace traffic splits three ways: identifier and attribute queries like SKUs and model numbers; head category queries like ``running shoes''; and a long conceptual tail --- ``gift for a coffee snob'', vernacular, misspellings, other languages.\par\vspace{3pt}
Lexical owns the first bucket, and not by accident. Embeddings blur X100 and X200 by construction; the entire point of the geometry is that similar strings land near each other, which is exactly wrong for a one-character part number. BM25 also gives you freshness for free --- a new listing is searchable the moment it's indexed, with no encoder version to wait on --- plus structured filters on price and category in the same engine, and cheap, explainable head-query serving. Dense owns the tail, where vocabulary mismatch dominates and there's no click history to boost with --- plus cross-lingual matching and typo robustness.\par\vspace{3pt}
They fail differently and the error sets are complementary, so union the candidates and fuse --- reciprocal rank fusion, or better, both scores as features into the reranker. Never sum BM25 and cosine directly; the scales aren't comparable.\par\vspace{3pt}
Two staff-level additions. Make it a routing decision, not only a fusion decision --- detect identifier-like tokens and weight lexical up rather than running everything at full depth. And be honest about cost: hybrid doubles the operational surface, so a team without eval infrastructure should fix lexical and query understanding first: spell correction, synonyms, zero-result relaxation.\par\vspace{3pt}
\flab{Loses the signal}
\begin{itemize}
\item ``Dense replaces lexical'' --- fatal in a marketplace with SKUs, filters and churning inventory
\item A 50/50 hand-wave with no query-mix analysis and no routing story
\item Summing BM25 and cosine scores without addressing the scale mismatch
\end{itemize}
\lab{Then they ask}
\begin{itemize}
\item Your dense retriever tanks on part-number queries --- fix it without removing it.
\item One retrieval improvement this quarter: which one?
\item How does the answer change for a documentation-search product?
\end{itemize}
\vspace{2pt}\noindent{\color{muted}\rule{\linewidth}{0.3pt}}
\needspace{6\baselineskip}
\qhead{L5}{Conceptual}{How does HNSW work, and why is it fast?}
\lab{Say this}
HNSW is a proximity graph over your vectors with a small-world hierarchy on top, searched greedily with a bounded beam.\par\vspace{3pt}
Every vector is a node. Each draws a random top level from a geometric distribution, so layers thin out by a constant factor going up --- layer zero holds everything, the top layer a handful. Within a layer each node keeps at most M links, twice that at layer zero, and they're picked with a diversity heuristic rather than being the M literal nearest: you keep neighbors covering different directions, so you get real long-range edges instead of a tight local clique.\par\vspace{3pt}
Search enters at the top and greedily hops toward the query, dropping a layer whenever no neighbor improves. At layer zero you widen into a best-first beam of size efSearch and stop when nothing unexplored beats your worst kept result.\par\vspace{3pt}
Three things make it fast. The hierarchy solves the entry problem --- you land in the query's neighborhood in roughly log-many hops instead of walking there through layer zero. Each hop is only M distance computations against a constant-degree list. And the beam gives cheap backtracking, which matters because distance concentration in high dimensions creates local optima that greedy walks fall into. That's why efSearch is the recall-latency dial.\par\vspace{3pt}
The caveat I'd volunteer: log scaling is empirical, on data with low intrinsic dimension, not a theorem. And you pay in RAM --- vectors plus edges, all resident, with random access that fits disk badly.\par\vspace{3pt}
\lab{Numbers}
\begin{itemize}
\item Degree: at most M links per node per layer, 2M at layer 0
\end{itemize}
\flab{Loses the signal}
\begin{itemize}
\item Describing it as a tree or a binary search --- it's a graph, best-first, with no partitioning and no pruning certificate.
\item Stating O(log n) as a proven worst-case bound.
\item Not knowing what efSearch does --- it's the parameter every production tuning session revolves around.
\end{itemize}
\lab{Then they ask}
\begin{itemize}
\item What happens if efSearch is less than k?
\item Why is deletion hard?
\item When would you pick IVF-PQ over HNSW?
\end{itemize}
\vspace{2pt}\noindent{\color{muted}\rule{\linewidth}{0.3pt}}
\needspace{6\baselineskip}
\qhead{L5}{Conceptual}{Bi-encoder vs. cross-encoder: why not cross-encode everything, and what exactly does the cross-encoder buy?}
\lab{Say this}
Because it isn't a constant-factor difference, it's a different complexity class. A bi-encoder's score decomposes: document vectors are computed offline and indexed, so a query costs one encoder pass plus a sublinear index probe, whether the corpus is a million documents or a billion. A cross-encoder scores the pair jointly, so nothing precomputes --- scoring N documents means N full forward passes. At 100 million documents and roughly 5$\times$10$^1$$^0$ FLOPs per pair for BERT-base at 256 tokens, that's 5$\times$10$^1$$^8$ FLOPs for one query, thousands of GPU-seconds. No hardware fixes that, because the problem is that the score isn't decomposable.\par\vspace{3pt}
What the joint pass buys is term-level interaction. The bi-encoder has to compress query and document into single vectors before they ever meet, so matching is one dot product over pooled representations and fine-grained correspondence is gone. The cross-encoder attends across the pair token by token, so it can resolve exact identifiers, negation, and which passage answers which aspect of the query. That's why it reliably adds NDCG on top of any first stage.\par\vspace{3pt}
So it isn't either-or, it's a funnel: recall-oriented decomposable retrieval feeds a precision-oriented cross-encoder over the top hundred or so, each stage spending more per candidate on fewer candidates. ColBERT sits in between --- per-token document embeddings precomputed, MaxSim at query time --- recovering much of the interaction quality at a storage multiple. And the stages are coupled in training too: the cross-encoder is the bi-encoder's distillation teacher and the denoiser for its mined hard negatives.\par\vspace{3pt}
\lab{Numbers}
\begin{itemize}
\item BERT-base, 256 tokens: \textasciitilde{}5$\times$10$^1$$^0$ FLOPs per query-document pair
\item Cross-encoding 10$^8$ docs = \textasciitilde{}5$\times$10$^1$$^8$ FLOPs per query --- thousands of GPU-seconds
\item Affordable rerank window: 10$^2$--10$^3$ candidates
\end{itemize}
\flab{Loses the signal}
\begin{itemize}
\item Framing it as ``cross-encoders are slow'' without the decomposability argument --- that's why no hardware fixes it.
\item Treating it as an either-or choice rather than two stages of one funnel.
\item Proposing a cross-encoder as the first-stage retriever over a large corpus.
\end{itemize}
\lab{Then they ask}
\begin{itemize}
\item Where does ColBERT sit on this spectrum, and what does it cost?
\item How would you use the cross-encoder to make the bi-encoder better?
\item Your reranker budget is 20 ms --- what do you serve?
\end{itemize}
\vspace{2pt}\noindent{\color{muted}\rule{\linewidth}{0.3pt}}
\needspace{6\baselineskip}
\qhead{L6}{Debugging}{Your new reranker improved offline NDCG by 4\%, but online CTR dropped. Walk me through the diagnosis.}
\lab{Say this}
CTR dropping is different from CTR being flat - something actively got worse - and the cheapest check is also the most common cause: latency. A heavier reranker makes pages arrive later, and latency is itself a relevance metric; a hundred milliseconds measurably moves engagement. Worse, if the model blows its L2 time budget the system may silently shrink rerank depth from a hundred candidates to twenty, capping you at the recall ceiling online while the offline eval ran with no budget at all. So: p95 latency and timeout rates in the treatment arm, first.\par\vspace{3pt}
Second, presentation and what CTR actually measures. Clicks respond to titles, snippets and images, not relevance in the abstract, so a better ordering that surfaces uglier results loses clicks. And clicks can drop because you succeeded - a good snippet satisfies the user without one. Check conversion, satisfied-session rate and reformulation before calling it a regression.\par\vspace{3pt}
Third, L3 and score consumers. Business rules may be reordering your new top twenty, so check override rates. And anything calibrated on the old score distribution - thresholds, vertical blending, merge rules - is now misfiring on scores that mean something different.\par\vspace{3pt}
Fourth, the offline eval itself: judged-only NDCG sees a candidate pool production doesn't, click labels inherit the old ranker's position bias, and gains on tail segments can hide losses on head queries that dominate impressions.\par\vspace{3pt}
My order: latency dashboards, segment-sliced online deltas, L3 and score-consumer audit, then the label-bias audit.\par\vspace{3pt}
\lab{Numbers}
\begin{itemize}
\item \textasciitilde{}100 ms of added latency measurably reduces engagement
\end{itemize}
\flab{Loses the signal}
\begin{itemize}
\item Concludes ``offline metrics are useless'' or retrains on fresh clicks with no diagnosis.
\item Never mentions latency - the most common real cause and the cheapest to check.
\item Treats CTR as ground truth, with no good-abandonment or presentation reasoning.
\end{itemize}
\lab{Then they ask}
\begin{itemize}
\item Which single dashboard do you open first, and why?
\item How would you build an offline eval that would have predicted this?
\item CTR dropped but conversions rose - what does that tell you?
\end{itemize}
\vspace{2pt}\noindent{\color{muted}\rule{\linewidth}{0.3pt}}
\needspace{6\baselineskip}
\qhead{L5}{Conceptual}{You run BM25 and dense retrieval in parallel. How do you combine them - and why does half BM25 plus half cosine fail?}
\lab{Say this}
The weighted sum fails because the two scores are not on comparable scales, and the mismatch is query-dependent, which is what makes it unfixable by rescaling. BM25 is unbounded and its magnitude moves with query length and term rarity, so a query containing one very rare term inflates BM25 for every candidate and drowns the dense signal at any fixed alpha. Cosine is bounded and concentrates in a narrow band. Any alpha you tune on one query mix quietly breaks on another.\par\vspace{3pt}
Per-query min-max or z-scoring does not rescue it either. You are normalising by order statistics of a small, noisy candidate sample, assuming distribution shapes neither system has, and throwing away absolute-confidence information.\par\vspace{3pt}
What I would ship first is reciprocal rank fusion: sum one over k plus rank across the retrievers, with k around sixty. It is rank-based, so it is invariant to any monotone rescaling of either score - the incompatibility problem disappears by construction. The k constant damps top-rank dominance; at sixty, rank one and rank ten contribute comparably, so a document both retrievers agree on can outrank one system's top hit. No tuning, no training data, and it degrades gracefully if a retriever goes down.\par\vspace{3pt}
The end state is better than RRF, though. RRF throws away score magnitude and cannot learn that dense should dominate conceptual queries while lexical dominates identifier lookups. So you union the candidates and hand the ranker both raw scores, both ranks, and query features, and fusion becomes learned and query-conditional.\par\vspace{3pt}
\lab{Numbers}
\begin{itemize}
\item RRF: score = sum of 1/(k + rank), k \textasciitilde{} 60
\end{itemize}
\flab{Loses the signal}
\begin{itemize}
\item Proposing the weighted sum without flagging the scale problem - interviewers use this as a competence gate.
\item ``Just normalise per query'' - that is the second layer of the trap, not the fix.
\item Knowing the RRF formula but not what k does or why rank-based fusion is robust.
\end{itemize}
\lab{Then they ask}
\begin{itemize}
\item What does k = 0 do to RRF?
\item When would learned fusion beat RRF materially?
\item One retriever goes down - what does your fusion degrade to?
\end{itemize}
\vspace{2pt}\noindent{\color{muted}\rule{\linewidth}{0.3pt}}
\needspace{6\baselineskip}
\qhead{L6}{Debugging}{Your click-trained ranker keeps favoring the items that have historically sat at position 1---better new items never rise. Diagnose and fix.}
\lab{Say this}
This is the position-bias feedback loop, and it gets in through three doors at once --- labels, features, and the loop itself. You have to close all three, because fixing one leaves the others reintroducing the same bias.\par\vspace{3pt}
Labels first. Clicks are examination-gated --- position one gets several times the examination of position five --- so click-derived labels systematically overrate whatever the old ranker put on top. The model learns to predict what was ranked high, not what's relevant.\par\vspace{3pt}
Features second, and this is the door people miss. Historical CTR and engagement aggregates are accumulated exposure, not merit. Incumbents arrive with a fat prior; a new item arrives with an empty one the model reads as ``bad'' rather than ``unknown.'' Feature importance on those priors tells you immediately if that's what's happening.\par\vspace{3pt}
And the loop: this model's output decides tomorrow's exposure, which decides tomorrow's labels and features. Every retrain entrenches the incumbents further.\par\vspace{3pt}
Fixes in order of leverage. On labels, train counterfactually --- inverse propensity weighting on examination, propensities from adjacent-pair randomization or regression EM, clipped for variance; the cheap version is position as a training feature held constant at serving. On features, impression-normalize to rates, Bayesian-smooth toward a category prior so empty history reads as prior rather than bad. On the loop, dedicate exploration traffic, keep a small always-on randomized slice as your unbiased eval set, and put new-item exposure share into the launch criteria.\par\vspace{3pt}
\flab{Loses the signal}
\begin{itemize}
\item Proposing a hand-tuned freshness boost and stopping --- treats the symptom, leaves all three doors open.
\item No propensities and no randomization anywhere in the answer.
\item Debiasing labels but leaving the engagement features, which reintroduce the same bias through a second door.
\end{itemize}
\lab{Then they ask}
\begin{itemize}
\item How exactly would you estimate the propensities here?
\item Your IPS-trained model's offline NDCG dropped versus the biased baseline --- is that bad?
\item How much exploration traffic is enough?
\end{itemize}
\vspace{2pt}\noindent{\color{muted}\rule{\linewidth}{0.3pt}}
\chapter{Recommenders}
\addtocontents{toc}{\protect\thispagestyle{fancy}}
\needspace{6\baselineskip}
\qhead{L5}{Conceptual}{How do you handle the cold start problem for new users and new items?}
\lab{Say this}
Cold start is a data problem, not a model problem --- the job is to design the system so it collects useful signal fast, and new users and new items need different answers.\par\vspace{3pt}
For new users, you fall back on what you know without interactions: popularity, plus content and context features --- device, locale, referral source, whatever the signup flow gives you. Ask a couple of preference questions at onboarding if the product tolerates it. Then run a bandit --- Thompson sampling or epsilon-greedy --- so the first handful of interactions actually move the profile instead of re-serving the same popularity list.\par\vspace{3pt}
For new items, place the item in embedding space from its content: run title, description and image through encoders, or give it a content-derived semantic ID so it inherits signal from similar items that already have history. Then buy the interaction data deliberately --- a small forced-impression budget --- and log the propensities, so what you collect is usable for training rather than one more biased sample.\par\vspace{3pt}
The Staff-level addition is where in the funnel this happens. Cold items have to be handled at retrieval, not just reranking: if the candidate generator never surfaces the item, nothing downstream rescues it. And cold start ends at different scales --- a handful of interactions gives you a usable user profile, but you need hundreds of impressions before item statistics are stable.\par\vspace{3pt}
\flab{Loses the signal}
\begin{itemize}
\item ``Just use collaborative filtering'' --- there is no interaction data, that's the premise.
\item No exploration story, so the system never collects the signal it's missing.
\item Fixing it only in the ranker --- a cold item retrieval never surfaces can't be rescued downstream.
\end{itemize}
\lab{Then they ask}
\begin{itemize}
\item How do you balance exploration against user experience?
\item How would you evaluate whether the cold-start strategy is working?
\item How much forced-impression budget is enough?
\end{itemize}
\vspace{2pt}\noindent{\color{muted}\rule{\linewidth}{0.3pt}}
\needspace{6\baselineskip}
\qhead{L5}{Conceptual}{How do you handle the cold start problem for new users and new items?}
\lab{Say this}
Cold start is a data problem, not a modeling problem: the system's job is to collect useful signal fast, and the model is downstream of that.\par\vspace{3pt}
Split it into users and items, because they need opposite things. For a new user you have no interactions but you're not blind - device, locale, referral source, sometimes a social login. Serve popularity and content-based recommendations conditioned on whatever you have, ask for a few explicit preferences at onboarding if the product tolerates it, and run a contextual bandit, Thompson sampling or epsilon-greedy, so the first few interactions actively buy information instead of exploiting a weak prior.\par\vspace{3pt}
For a new item it's the opposite: rich content, no behavior. Use content features - title, description, images, category, price - to place the item near similar known items in embedding space so it inherits their collaborative signal until it has its own, which is where CLIP-style cross-modal embeddings help. Then explicitly budget exploration: force it in front of a small slice of traffic to earn interaction data, and apply inverse propensity weighting when you train on that data so you don't bake the exploration policy into the model.\par\vspace{3pt}
The failure I'd call out is treating exploration purely as a cost. If nothing is shown, nothing is learned, and the tail of the catalog stays permanently cold.\par\vspace{3pt}
\flab{Loses the signal}
\begin{itemize}
\item ``Just use collaborative filtering'' - there is nothing to collaborate on yet.
\item No exploration story, so new items never accumulate the data that would rank them.
\item Doesn't separate user cold start from item cold start; they need opposite fixes.
\end{itemize}
\lab{Then they ask}
\begin{itemize}
\item How do you balance exploration against exploitation here?
\item When does an item stop being cold?
\item How would you evaluate whether the cold-start path is working?
\end{itemize}
\vspace{2pt}\noindent{\color{muted}\rule{\linewidth}{0.3pt}}
\needspace{6\baselineskip}
\qhead{L6}{Debugging}{Your recommendation model has great offline metrics but poor online performance. What went wrong?}
\lab{Say this}
The most likely answer is that your offline evaluation is scoring the model on data the old policy generated, so a model that reproduces the old policy looks excellent and adds nothing.\par\vspace{3pt}
That is selection bias: the logs only contain items that were shown, and you never observe the counterfactual. If you evaluate on that same biased distribution, a high metric means the model learned to predict the incumbent's behaviour. Position bias is the same problem one level down - top slots get clicked regardless of relevance, so a model that has quietly learned position effects shows strong AUC with no real ranking improvement.\par\vspace{3pt}
The other three I would check are leakage, shift and metric mismatch. Feature leakage is something in the offline pipeline computed over the full time range, with item popularity the usual culprit, that is not available at serving time. Temporal shift is validating on a distribution that no longer exists - December training data failing in January. And proxy mismatch is AUC or NDCG simply not moving the business metric, often because the model is accurate on easy cases and miscalibrated exactly on the high-value tail.\par\vspace{3pt}
The fixes follow the diagnosis: temporal splits and never random ones on time-dependent data, inverse propensity weighting or explicit position-debiased training, monitoring calibration and not only ranking metrics, and treating every offline win as a hypothesis that an A/B test confirms rather than as a result.\par\vspace{3pt}
\flab{Loses the signal}
\begin{itemize}
\item ``The offline metrics must be wrong'' - the gap is expected and has causes worth naming.
\item Never raising selection bias or position bias.
\item Using a random train/test split on logged interaction data.
\end{itemize}
\lab{Then they ask}
\begin{itemize}
\item How do you design a good A/B test for a recommender?
\item How would you detect feature leakage?
\item What is counterfactual evaluation and when is it worth it?
\end{itemize}
\vspace{2pt}\noindent{\color{muted}\rule{\linewidth}{0.3pt}}
\needspace{6\baselineskip}
\qhead{L6}{Estimation}{Estimate the memory required for DLRM embedding tables with 100M users, 10M items, and 1000 categorical features.}
\lab{Say this}
Around 70 gigabytes of embedding tables in FP32, and the interesting part is which term dominates. Users are the biggest single table --- 100 million times 128 dimensions times 4 bytes is 51 gigabytes. Items at 10 million are only 5. So the user table alone is most of the user-plus-item total.\par\vspace{3pt}
For the thousand categorical features I'd bucket by cardinality rather than pretend they're uniform, and state the assumption out loud. Say fifty high-cardinality features like city or campaign ID at about a million entries each, a couple of hundred medium ones around ten thousand, and the remaining several hundred at a hundred or so, all at dimension 64. The high-cardinality group comes out at 13 gigabytes, medium at half a gigabyte, and the low-cardinality group is noise --- under 20 megabytes. Total lands near 70.\par\vspace{3pt}
The MLP is under a hundred million parameters, well under a gigabyte, which is the structural point about DLRM: it's a memory system with a small compute head attached.\par\vspace{3pt}
Training is where it bites. Adam keeps two moments per parameter, so 12 bytes rather than 4 --- call it 210 gigabytes. That doesn't fit on a single accelerator, so you shard the embedding tables model-parallel across GPUs and run the MLP data-parallel. At Meta scale the ID space is in the billions and there are thousands of features, which is how you get to trillion-parameter recommenders.\par\vspace{3pt}
\lab{Numbers}
\begin{itemize}
\item Users 100M $\times$ 128 $\times$ 4B = 51.2 GB; items = 5.1 GB
\item Total tables $\approx$ 70 GB; with Adam state $\approx$ 210 GB
\item Adam training: 12 bytes per parameter
\end{itemize}
\flab{Loses the signal}
\begin{itemize}
\item Forgetting optimizer state --- Adam triples parameter memory
\item Assuming all 1000 features share one vocabulary size
\item Not noticing that the user table dominates and the MLP is negligible
\end{itemize}
\lab{Then they ask}
\begin{itemize}
\item How would you shard these tables across GPUs?
\item What breaks if you double the embedding dimension?
\item How much would mixed-dimension embeddings save?
\end{itemize}
\vspace{2pt}\noindent{\color{muted}\rule{\linewidth}{0.3pt}}
\needspace{6\baselineskip}
\qhead{L6}{Debugging}{Your CTR model's offline AUC improved by 0.5\% but online revenue dropped 3\%. Diagnose.}
\lab{Say this}
Calibration, first hypothesis, before anything else. A 0.5 percent AUC lift is real, so the model genuinely ranks better --- but AUC is invariant to any monotone transform of the scores, and the auction doesn't consume ranks, it consumes probabilities. If the new model over-predicts pCTR by, say, one and a half times, you're charging advertisers above fair value, their daily budgets burn out earlier in the day, and you lose fill in the afternoon. Better ranking, less revenue. The diagnosis is a reliability plot --- predicted versus actual CTR by decile --- and if it's off the diagonal, recalibrate and rerun.\par\vspace{3pt}
Second, ask what the model did to the auction rather than to the ranking. Revenue is pCTR times bid, and AUC knows nothing about bids. If the new model shifted impressions toward higher-true-CTR but lower-bid advertisers, you lose money with a genuinely better model. Compare average winning bid and auction density before and after.\par\vspace{3pt}
Then the boring ones that are common anyway: training-serving skew --- log the features actually computed at serving time and diff them against offline features for the same impressions --- and offline-online distribution shift, where your test slice just isn't current traffic.\par\vspace{3pt}
The general lesson I'd state: for any model feeding an auction or a threshold, calibration is a launch gate, not a nice-to-have. AUC alone should never be enough to clear a model to ship.\par\vspace{3pt}
\flab{Loses the signal}
\begin{itemize}
\item ``0.5\% AUC is probably noise'' --- at this scale that's a meaningful, reproducible lift.
\item Not putting calibration first when the downstream consumer is an auction.
\item Blaming the A/B infrastructure before investigating model-level causes.
\end{itemize}
\lab{Then they ask}
\begin{itemize}
\item How would you fix the calibration?
\item What guardrail metrics would you add to the experiment?
\item How do you stop this happening on the next launch?
\end{itemize}
\vspace{2pt}\noindent{\color{muted}\rule{\linewidth}{0.3pt}}
\needspace{6\baselineskip}
\qhead{L5}{Trade-off}{Two-Tower vs. interaction-based models for candidate generation---when would you use each?}
\lab{Say this}
Two-tower for candidate generation, interaction-based for ranking, and the reason is the computational budget at each stage rather than which model is more accurate.\par\vspace{3pt}
Candidate generation takes ten million items down to a thousand, and the only thing that works at that width is a model whose item side can be precomputed. Two-tower encodes user and item independently and scores by dot product, so item embeddings go into an ANN index offline and serving is one user-tower pass plus a sub-millisecond lookup, independent of catalog size. What you pay is expressiveness --- the dot product is the only interaction, so there are no cross-features at all.\par\vspace{3pt}
Ranking takes a thousand down to a hundred, so you can afford a forward pass per candidate, and that's where interaction-based models earn their keep. DLRM, DCN, DIN --- cross networks, attention over behavior history, MLPs over the joint feature set. They capture things a two-tower structurally cannot, like a user who buys cheap items in one category and premium in another. And by definition they can't be precomputed, because the score depends on the pair.\par\vspace{3pt}
The hybrid worth raising is pushing lightweight interaction into the two-tower without losing precomputation --- query-item attention inside the item tower at training time, or distilling the ranker into the retrieval model. Retrieval quality goes up and the serving shape doesn't change.\par\vspace{3pt}
\lab{Numbers}
\begin{itemize}
\item Candidate generation \textasciitilde{}10M $\to$ 1K; ranking \textasciitilde{}1K $\to$ 100
\end{itemize}
\flab{Loses the signal}
\begin{itemize}
\item ``Always use interaction-based models, they're more accurate'' --- ignores the computational constraint at retrieval.
\item Not recognizing that the pipeline stage's budget is what drives the choice.
\item Conflating retrieval and ranking, which have fundamentally different constraints.
\end{itemize}
\lab{Then they ask}
\begin{itemize}
\item Can you make a two-tower model more expressive without losing precomputation?
\item How do you generate hard negatives for two-tower training?
\item Which ANN algorithm would you use and why?
\end{itemize}
\vspace{2pt}\noindent{\color{muted}\rule{\linewidth}{0.3pt}}
\needspace{6\baselineskip}
\qhead{L5}{Conceptual}{Why do recommendation models need calibration? What happens if pCTR is systematically over-confident?}
\lab{Say this}
Because the predicted probability isn't only used for ordering - it's consumed as a number in downstream decisions, and those break when it's wrong even if the ranking is perfect.\par\vspace{3pt}
Three places it's used as an absolute value. Auction pricing: ad rank is pCTR times bid, and what the advertiser is charged is a function of pCTR, so miscalibration is directly mispricing. Budget pacing: the system estimates expected spend per impression from pCTR to spread a daily budget evenly. And absolute thresholds - suppress anything below some pCTR - where a shift changes which items pass at all.\par\vspace{3pt}
If pCTR is systematically two times too high, here's the chain. Advertisers overpay for clicks that don't arrive, so measured ROI drops. Budgets burn out faster, so campaigns go dark by mid-afternoon, leaving inventory unfilled with fewer bidders competing - revenue falls at exactly the time of day you needed it. Then the slow damage: advertisers who consistently overpay lower their bids or leave.\par\vspace{3pt}
And here's what makes it insidious. If the overconfidence is uniform, the ranking is completely unchanged. The same ads win in the same order, AUC looks great, every ranking metric on your dashboard looks great. All the damage is in pricing and pacing, which is exactly why calibration has to be monitored separately from AUC.\par\vspace{3pt}
\lab{Numbers}
\begin{itemize}
\item Ad rank = pCTR x bid, so a 2x-overconfident pCTR leaves ranking untouched and doubles implied value
\end{itemize}
\flab{Loses the signal}
\begin{itemize}
\item ``Calibration doesn't matter if the ranking is correct'' - true for ordering, false for pricing and pacing.
\item Can't separate ranking quality from calibration quality, or assumes AUC covers both.
\item Treats calibration as a nice-to-have rather than a revenue-critical property.
\end{itemize}
\lab{Then they ask}
\begin{itemize}
\item How would you measure calibration?
\item Which recalibration method would you use, and why?
\item How often do you need to recalibrate?
\end{itemize}
\vspace{2pt}\noindent{\color{muted}\rule{\linewidth}{0.3pt}}
\chapter{Classical ML}
\addtocontents{toc}{\protect\thispagestyle{fancy}}
\needspace{6\baselineskip}
\qhead{L5}{Mathematical}{Derive the normal equations for linear regression. When would you not use the closed form?}
\lab{Say this}
Minimize the squared residual norm, expand, take the gradient, set it to zero: X-transpose X beta-hat equals X-transpose y, so beta-hat is the inverse of X-transpose X times X-transpose y, whenever X has full column rank. The Hessian is two X-transpose X, positive semidefinite, so that stationary point is a global minimum.\par\vspace{3pt}
The thing to say while writing it is the geometry: the condition is X-transpose times the residual equals zero, so the residual is orthogonal to the column space. OLS is the orthogonal projection of y onto the span of your features.\par\vspace{3pt}
When wouldn't I use it? Four separate reasons. Cost - it's cubic in the number of features, and at ten thousand features the Gram matrix alone is 800 megabytes. Singularity - more features than rows, or any exactly collinear feature like the dummy variable trap, and there's no inverse. Conditioning - the condition number of X-transpose X is the square of the condition number of X, so forming the Gram matrix costs you half your significant digits. And the setting - streaming or distributed data, or a non-quadratic loss, has no closed form at all.\par\vspace{3pt}
Then the implementation point that separates candidates: never call inverse. Cholesky when it's well conditioned, QR on X directly when it isn't, SVD when X may be rank deficient - which is what lstsq does, and why it returns a minimum-norm answer where the textbook formula would divide by zero.\par\vspace{3pt}
\lab{Numbers}
\begin{itemize}
\item Cost O(n d\textasciicircum{}2 + d\textasciicircum{}3); at d = 10,000 the Gram matrix alone is \textasciitilde{}800 MB
\item kappa(X\textasciicircum{}T X) = kappa(X)\textasciicircum{}2 - forming the Gram matrix squares the condition number
\end{itemize}
\lab{If they hand you the marker}
$X^{\top}X\hat{\beta} = X^{\top}y \;\Rightarrow\; \hat{\beta} = (X^{\top}X)^{-1}X^{\top}y$, then directly underneath it $X^{\top}(y - X\hat{\beta}) = 0$. Write the second line and say it out loud: the residual is orthogonal to every column of X, so OLS is a projection. That one line is what turns a memorized formula into a derivation, and it sets up the whole conditioning discussion.\par
\flab{Loses the signal}
\begin{itemize}
\item Offers inv(X.T @ X) @ X.T @ y as the recommended implementation.
\item Names only ``it's slow'' and misses singularity and conditioning.
\item Claims OLS needs normally distributed features or targets; normality is about the errors, and only for inference.
\end{itemize}
\lab{Then they ask}
\begin{itemize}
\item Why does kappa(X\textasciicircum{}T X) = kappa(X)\textasciicircum{}2 matter in practice?
\item What does the SVD give you when X is rank-deficient?
\item How does leverage let you compute LOOCV from a single fit?
\end{itemize}
\vspace{2pt}\noindent{\color{muted}\rule{\linewidth}{0.3pt}}
\needspace{6\baselineskip}
\qhead{L5}{First Principles}{Derive logistic regression's loss from maximum likelihood, give its gradient, and explain why we do not just use squared error.}
\lab{Say this}
The loss is not a design choice - it is the Bernoulli negative log-likelihood, and everything else follows.\par\vspace{3pt}
You assume each label is Bernoulli with parameter p and model the log-odds as linear. Inverting that link gives p equals sigmoid of x-transpose-beta, so the sigmoid is the inverse link, not an arbitrary squashing choice. Write the likelihood, take the negative log, and what falls out literally is binary cross-entropy - log loss and Bernoulli NLL are the same object.\par\vspace{3pt}
The gradient is the clean part. The loss derivative with respect to p carries p times one-minus-p in the denominator, and the sigmoid's derivative is p times one-minus-p - they cancel exactly, leaving prediction minus label. So the gradient is X-transpose times sigmoid-of-X-beta minus y, structurally identical to least squares, which is the canonical-link result in GLMs, not a coincidence. The Hessian is positive semi-definite, so the problem is convex; there is no closed form, and you fit with IRLS or L-BFGS.\par\vspace{3pt}
Why not squared error? Three separate reasons. It is the Gaussian likelihood and the label is Bernoulli. Composed with the sigmoid it is not convex in beta. And it saturates: at a confidently wrong point the squared-error gradient carries an extra sigmoid-prime and comes out roughly twenty thousand times smaller than log loss's, so the example that most deserves an update contributes almost nothing.\par\vspace{3pt}
One nuance: squared error on probabilities is the Brier score, which is proper, so the objection is about likelihood and geometry, not calibration.\par\vspace{3pt}
\lab{Numbers}
\begin{itemize}
\item At z = -10 with y = 1: squared-error gradient \textasciitilde{}4.5e-5 vs log loss \textasciitilde{} -1, about 22,000x smaller
\end{itemize}
\lab{If they hand you the marker}
$$\nabla_\beta \mathcal{L} = X^{\top}\!\left(\sigma(X\beta) - y\right), \qquad H = X^{\top} W X, \; W = \mathrm{diag}\!\left(p_i(1-p_i)\right)$$ Say while writing: the chain rule gives (p - y)/(p(1-p)) times p(1-p) - cross them out and only p minus y survives. Then W is non-negative, so the Hessian is PSD and the problem is convex.\par
\flab{Loses the signal}
\begin{itemize}
\item ``Logistic regression has a closed-form solution.''
\item Starting from ``we apply a sigmoid'' and being unable to say where the loss came from.
\item ``MSE gives uncalibrated probabilities'' - the Brier score is proper, so this is confidently wrong.
\end{itemize}
\lab{Then they ask}
\begin{itemize}
\item Show the Hessian and prove convexity.
\item What is IRLS actually doing?
\item How does this change for K classes, and what is the identifiability issue?
\end{itemize}
\vspace{2pt}\noindent{\color{muted}\rule{\linewidth}{0.3pt}}
\needspace{6\baselineskip}
\qhead{L6}{Trade-off}{When would you deploy logistic regression instead of gradient boosting in 2026?}
\lab{Say this}
Refuse the implied ``never'' and give the conditions, because there are four real ones. First, regulated or contested decisions --- credit, insurance, hiring, clinical. There you need a model whose behaviour can be stated as a monotone effect per feature and audited by people who aren't ML engineers. A coefficient table with odds ratios is an artifact a regulator accepts, and adverse-action reasons fall straight out of it. Scorecards --- binned features into logistic regression --- exist for that reason, not because nobody tried boosting.\par\vspace{3pt}
Second, extreme-scale sparse online serving, the ad-CTR heritage: hundreds of millions of hashed features, billions of daily examples, continuous retraining. Logistic regression with FTRL-Proximal trains online in one pass, the L1 term keeps the served model sparse and small, and it scores in microseconds. GBDTs don't update online and don't love ultra-high-cardinality sparse inputs.\par\vspace{3pt}
Third, small data with a genuinely near-linear signal --- a few thousand rows, where a regularized linear model often matches a tuned GBDT and is far more stable across refits. Fourth, calibration: log loss is the Bernoulli likelihood, so you get usable probabilities without a second stage, which matters when the score is multiplied by a value to make a bidding or expected-loss decision.\par\vspace{3pt}
Then be honest about the default: on mid-scale tabular data with real interactions, boosting wins, usually by a lot. The hybrids are splines and feature crosses to give the linear model curvature, or a GBDT for feature discovery with a linear model in serving.\par\vspace{3pt}
\flab{Loses the signal}
\begin{itemize}
\item ``Never, boosting is strictly better'' --- the trap the question is built around
\item ``Logistic regression is more interpretable'' as the entire answer, with no account of when that's a requirement
\item Not knowing the online sparse regime, where linear models still run at the largest scale in industry
\end{itemize}
\lab{Then they ask}
\begin{itemize}
\item What is FTRL-Proximal doing that plain SGD isn't?
\item How would you give a linear model nonlinearity without losing interpretability?
\item How do you show a regulator what the model does?
\end{itemize}
\vspace{2pt}\noindent{\color{muted}\rule{\linewidth}{0.3pt}}
\needspace{6\baselineskip}
\qhead{L5}{Mathematical}{Derive the optimal leaf weight and the split-gain formula in XGBoost.}
\lab{Say this}
XGBoost is Newton boosting, and everything falls out of one move: Taylor-expand the loss to second order around the current predictions.\par\vspace{3pt}
At round m you add a tree to frozen predictions, regularized by gamma per leaf plus L2 on leaf weights. Expand to second order in the tree's output, call g and h the loss's first and second derivatives at the current prediction, and the objective now involves only g, h and the tree.\par\vspace{3pt}
Now group by leaf. Every point in a leaf shares its weight, so with G the sum of gradients in that leaf and H the sum of hessians, the objective becomes independent one-dimensional quadratics, one per leaf, each minimized at minus G over H plus lambda. Substitute back and the best objective for a given tree structure is minus a half times the sum of G squared over H plus lambda, plus gamma times the leaf count.\par\vspace{3pt}
The split gain is just the change in that: a half of left plus right minus parent, each term G squared over H plus lambda, minus gamma. Gamma is the toll for adding a leaf, which makes it principled pruning rather than a heuristic.\par\vspace{3pt}
Two sanity checks: under squared loss the hessian is one and the leaf weight is a ridge-shrunk mean residual; under log loss g is p minus y and h is p times one minus p. And the split score needs nothing but G and H sums, which is why it distributes.\par\vspace{3pt}
\lab{If they hand you the marker}
$w_j^* = -\frac{G_j}{H_j+\lambda}$, then $\text{Gain} = \frac{1}{2}\left[\frac{G_L^2}{H_L+\lambda} + \frac{G_R^2}{H_R+\lambda} - \frac{(G_L+G_R)^2}{H_L+H_R+\lambda}\right] - \gamma$ --- write the leaf weight first, then say the gain is just parent objective minus children, and gamma is the price of the extra leaf.\par
\flab{Loses the signal}
\begin{itemize}
\item Only doing the first-order expansion --- there's no closed form without the hessian, and ``XGBoost is gradient boosting but faster'' misses the Newton objective entirely.
\item Reciting the formula without lambda in the denominator or gamma in the gain --- for the library whose selling point is regularization.
\item Unable to say what g and h are for log loss.
\end{itemize}
\lab{Then they ask}
\begin{itemize}
\item What changes with L1 on the leaf weights?
\item Why is \texttt{min\_child\_weight} different for regression versus classification?
\item Where does the histogram method approximate this derivation?
\end{itemize}
\vspace{2pt}\noindent{\color{muted}\rule{\linewidth}{0.3pt}}
\needspace{6\baselineskip}
\qhead{L5}{Conceptual}{Random forest vs. gradient boosting---how do they differ, and when do you reach for each?}
\lab{Say this}
They're opposite answers to the same problem --- a single deep tree has low bias and terrible variance.\par\vspace{3pt}
Random forest attacks the variance. Grow many deep, low-bias trees on bootstrap samples with feature subsampling at each split, then average. Averaging correlated estimators leaves the correlation-weighted part of the variance plus a term that shrinks with the number of trees, so more trees never hurts --- you cannot overfit in the number of trees, you only run out of compute. It trains embarrassingly in parallel and gives you out-of-bag validation for free.\par\vspace{3pt}
Boosting attacks the bias. Add shallow, high-bias trees sequentially, each fit to the negative gradient of the loss at the current predictions. So it's inherently sequential, and it does overfit in the number of rounds, which is why shrinkage and early stopping aren't optional.\par\vspace{3pt}
That also explains why the base learners differ, which is a good thing to have ready: averaging needs low-bias learners, so deep trees; sequential correction needs weak learners, so shallow ones.\par\vspace{3pt}
In practice, random forest is my first-hour baseline --- nearly tuning-free, robust on small or noisy data, free validation from OOB. A tuned LightGBM, XGBoost or CatBoost is what ships, and the gap is usually a few points in boosting's favor. I'd also reach for boosting specifically when I need a custom loss --- quantile, Poisson, ranking --- or monotonic constraints.\par\vspace{3pt}
\flab{Loses the signal}
\begin{itemize}
\item Inverting the mechanism --- ``RF reduces bias'' or ``boosting reduces variance.''
\item ``More trees overfits a random forest'' --- confusing the number of trees with boosting rounds.
\item Can't say why RF trees are deep while GBDT trees are shallow.
\end{itemize}
\lab{Then they ask}
\begin{itemize}
\item Derive the variance of an average of correlated estimators.
\item Why does feature subsampling help RF but column subsampling is optional for GBDT?
\item Which is more robust to label noise, and why?
\end{itemize}
\vspace{2pt}\noindent{\color{muted}\rule{\linewidth}{0.3pt}}
\needspace{6\baselineskip}
\qhead{L5}{First Principles}{Explain gradient boosting to someone who already understands gradient descent.}
\lab{Say this}
Gradient boosting is gradient descent where the space you're descending in is functions, not parameters. In ordinary gradient descent you have a parameter vector, compute the gradient of the loss with respect to it, and step against it. In boosting, the object you're updating is the vector of the model's predictions at the training points: you compute the gradient of the loss with respect to each prediction and step against that.\par\vspace{3pt}
The wrinkle is that the step is only defined at the n training points. You can't apply it directly - it says nothing about a new input and wouldn't generalize - so you project it onto a function class: fit a small tree to those negative gradients by least squares. The tree is your descent direction. Then update F to F plus eta times the tree, where eta is exactly a step size.\par\vspace{3pt}
Which is why ``boosting fits the residuals'' is a special case, not a definition. Under squared loss the negative gradient happens to be y minus F of x, so it looks like residual fitting. Under log loss the tree fits label minus predicted probability. Under a ranking objective, lambda gradients.\par\vspace{3pt}
Then round it out: rounds are iterations, the learning rate is the step size, and early stopping on a validation set is your convergence check - because unlike gradient descent on a convex problem, more rounds eventually overfit by construction.\par\vspace{3pt}
\flab{Loses the signal}
\begin{itemize}
\item Gives residual-fitting as the definition instead of the squared-loss special case.
\item No answer for why you fit a tree to the gradient instead of taking the gradient step directly.
\item Confuses the shrinkage learning rate with some learning rate inside the tree; trees are fit greedily, there is no SGD in there.
\end{itemize}
\lab{Then they ask}
\begin{itemize}
\item What are the pseudo-residuals under absolute loss, and what does that imply?
\item Why does a lower learning rate generalize better?
\item How does AdaBoost fit into this picture?
\end{itemize}
\vspace{2pt}\noindent{\color{muted}\rule{\linewidth}{0.3pt}}
\needspace{6\baselineskip}
\qhead{L6}{Trade-off}{XGBoost, LightGBM or CatBoost: 600K rows, 45 features of which 14 are categorical including a 40K-cardinality \texttt{merchant\_id}. Choose and defend.}
\lab{Say this}
CatBoost, and the 40,000-cardinality merchant ID is what decides it.\par\vspace{3pt}
All three share the same second-order boosted-tree objective, so accuracy differences are small and dataset-specific. What actually differs is categorical handling, speed at scale, overfitting posture and ecosystem - and with fourteen categoricals, one of them at forty thousand levels, categorical handling dominates. One-hot is out immediately; you would be adding forty thousand near-useless indicator columns. Naive target encoding is a leakage trap, because each row's own label leaks into its own feature.\par\vspace{3pt}
CatBoost's ordered target statistics solve exactly that: the encoding for each row uses only rows that precede it in a random permutation, so it is leak-free by construction, and its feature combinations pick up categorical interactions. At 600,000 rows you are well inside its comfortable speed envelope.\par\vspace{3pt}
The defensible alternative is LightGBM. Its native categorical splits sort levels by sum-gradient over sum-hessian, which handles the mid-cardinality features well, plus a properly out-of-fold target or frequency encoding for merchant ID. That iterates faster, but you own the encoding hygiene yourself. XGBoost is the fallback if your platform - Spark pipelines, existing serving - already speaks it; nothing in the data demands it.\par\vspace{3pt}
And I would say the tie-breaker out loud: one afternoon, all three with early stopping on an identical temporal validation split, comparing metric, training time and inference latency. Expect differences of tenths of a point and pick on operational grounds.\par\vspace{3pt}
\flab{Loses the signal}
\begin{itemize}
\item One-hot encoding the 40K-cardinality feature without hesitating.
\item Choosing on ``X is more accurate'' folklore and never mentioning categorical handling.
\item Proposing in-fold target encoding, or not knowing why CatBoost's ordering avoids the leak.
\end{itemize}
\lab{Then they ask}
\begin{itemize}
\item What exactly goes wrong if you target-encode with the full training set?
\item New merchants appear daily - what happens to each library's encoding?
\item Now it is 500 million rows. Does your answer change?
\end{itemize}
\vspace{2pt}\noindent{\color{muted}\rule{\linewidth}{0.3pt}}
\needspace{6\baselineskip}
\qhead{L6}{Conceptual}{Why does a gradient-boosted tree beat your MLP on this tabular dataset---and when would it not?}
\lab{Say this}
It's an inductive-bias match, not a capacity gap, and leading with the mechanism rather than a benchmark is the point. Three findings from the controlled comparisons Grinsztajn and colleagues ran in 2022. Tabular targets are irregular --- sharp thresholds and plateaus, exactly what piecewise-constant splits fit natively, while neural network training is biased toward smooth functions. Tables carry many uninformative columns, which a tree ignores by simply never splitting on them, while an MLP has to learn to suppress them. And MLPs are rotation-invariant while tabular data emphatically is not: each feature axis means something specific, so the tree's axis-aligned ``limitation'' is precisely the right prior. The clean demonstration is that randomly rotating the features barely affects an MLP and wrecks a tree, which tells you the axes carry the signal.\par\vspace{3pt}
Then the pragmatics that matter in practice: native handling of NaNs and mixed types, invariance to monotone transforms, minutes of CPU training, and graceful degradation when hyperparameters are wrong.\par\vspace{3pt}
When would it not hold? Heavy learned embeddings for high-cardinality IDs --- that's the recommender regime, and it's neural for good reason --- multimodal inputs, online or continual updating, and very large data with dense interactions. And on genuinely small tables, in-context learners like TabPFN v2 now beat tuned GBDTs outright. So boosting is the mid-scale tabular default, not a law of nature.\par\vspace{3pt}
\flab{Loses the signal}
\begin{itemize}
\item ``Neural nets are more powerful, it must be a tuning issue'' --- capacity isn't the axis, prior match is
\item No mechanism, just ``trees win on tabular'' as folklore
\item Unaware that the small-data regime has shifted, or that recommenders are neural for a reason
\end{itemize}
\lab{Then they ask}
\begin{itemize}
\item Design the experiment that tests whether uninformative features are the MLP's problem.
\item What is TabPFN actually doing at inference time?
\item Your tabular data has one free-text column --- now what?
\end{itemize}
\vspace{2pt}\noindent{\color{muted}\rule{\linewidth}{0.3pt}}
\needspace{6\baselineskip}
\qhead{L5}{Mathematical}{Derive PCA from first principles. Whiteboard it.}
\lab{Say this}
Set it up before writing anything: X is n by d, rows are examples, and it's centered. S is the sample covariance, X-transpose-X over n minus one. The goal is the unit direction of maximum projected variance.\par\vspace{3pt}
Because the data are centered, projections onto a direction w have mean zero, so their variance is just w-transpose S w. You need the unit-norm constraint --- without it you scale w up and the objective is unbounded, and it's worth saying that out loud. Form the Lagrangian, differentiate, and you get S w equals lambda w, so the stationary points are exactly the eigenvectors of the covariance. Then left-multiply by w-transpose: w-transpose S w equals lambda. The multiplier is the explained variance itself, which is the line interviewers are waiting for. For the second component you add orthogonality to the first, the extra multiplier drops out by symmetry, and induction gives you the rest.\par\vspace{3pt}
Then three things that turn a correct derivation into a strong one. Minimizing reconstruction error is the same problem --- by Pythagoras the error is the norm of x minus the squared projection, and that first term doesn't depend on w. In practice you don't eigendecompose the covariance, you SVD X, because forming X-transpose-X squares the condition number. And on preconditions: center always, standardize when units differ, and remember PCA is unsupervised --- nothing guarantees the top components are the discriminative ones.\par\vspace{3pt}
\lab{If they hand you the marker}
$\max_{\|w\|=1} w^\top S w \;\Rightarrow\; Sw = \lambda w \;\Rightarrow\; w^\top S w = \lambda$ --- write the constrained objective, then the eigenvalue equation, then left-multiply by $w^\top$ and say aloud that the multiplier *is* the variance.\par
\flab{Loses the signal}
\begin{itemize}
\item Omitting the norm constraint, or writing the Lagrangian and never saying that lambda is the explained variance.
\item Forgetting to center --- without it the first component points at the mean vector.
\item ``I'd eigendecompose X-transpose-X'' as the implementation plan, with no mention of SVD or conditioning.
\end{itemize}
\lab{Then they ask}
\begin{itemize}
\item Show me that maximizing variance and minimizing reconstruction error are the same problem.
\item What is whitening and when does it hurt?
\item Ten million rows, fifty components --- what do you actually run?
\end{itemize}
\vspace{2pt}\noindent{\color{muted}\rule{\linewidth}{0.3pt}}
\needspace{6\baselineskip}
\qhead{L6}{Mathematical}{Derive the EM updates for a Gaussian mixture model, and explain why the likelihood increases monotonically.}
\lab{Say this}
You need EM because the log-likelihood has a log sitting outside a sum over components, so nothing separates and there's no closed form. Introduce a latent assignment variable per point and the complete-data log-likelihood puts the log inside the sum, so every parameter separates.\par\vspace{3pt}
The E step is just Bayes. Given current parameters, compute the responsibility of component k for point i: the mixing weight times that component's density, normalized over components. Because the complete-data log-likelihood is linear in that indicator, taking its expectation is literally substituting the responsibility.\par\vspace{3pt}
The M step is then the ordinary maximum-likelihood update with hard counts replaced by soft responsibilities. The mean is the responsibility-weighted average of the points, the covariance the weighted scatter about it, and the mixing weight the total responsibility mass over n, which needs a Lagrange multiplier because the weights sum to one. One sentence to remember: soft counts in place of hard counts.\par\vspace{3pt}
Monotonicity comes from the ELBO. For any distribution over the latents, the log-likelihood equals the ELBO plus a KL to the true posterior. The E step sets that distribution to the posterior, driving the KL to zero, so the bound touches the likelihood; the M step then raises the bound. Bound touched, bound raised - the likelihood cannot go down. Monotone ascent to a stationary point, local not global. Practical consequence: a decreasing log-likelihood in your training log is a bug.\par\vspace{3pt}
\lab{If they hand you the marker}
$\gamma_{ik} = \dfrac{\pi_k\,\mathcal{N}(x_i \mid \mu_k, \Sigma_k)}{\sum_j \pi_j\,\mathcal{N}(x_i \mid \mu_j, \Sigma_j)}$ and, if they push on why it ascends, $\log p(X \mid \theta) = \mathcal{F}(q,\theta) + \mathrm{KL}\!\left(q \,\|\, p(Z \mid X, \theta)\right)$. Write the responsibility first and say every M-step formula is the plain MLE with gamma in place of a hard count; write the second line only when asked, and point at the KL going to zero as the moment the bound touches the likelihood.\par
\flab{Loses the signal}
\begin{itemize}
\item ``EM alternates two steps until convergence'' with no ELBO or Jensen argument for why the likelihood rises.
\item Can't say why direct MLE fails - the log of a sum.
\item Claims EM finds the global maximum, or doesn't know the unregularized GMM likelihood is unbounded above.
\end{itemize}
\lab{Then they ask}
\begin{itemize}
\item What happens if a component collapses onto a single point?
\item How do you choose the number of components?
\item Where does this become the VAE?
\end{itemize}
\vspace{2pt}\noindent{\color{muted}\rule{\linewidth}{0.3pt}}
\needspace{6\baselineskip}
\qhead{L5}{Trade-off}{How do you choose k? Your teammate says the elbow plot shows 8 clusters.}
\lab{Say this}
I would push back gently: there usually is no true k. Inertia falls monotonically and hits zero at k equals n, so the objective cannot select k - every method is a heuristic trading fit against parsimony or against a null model, and on real data they disagree by a factor of two.\par\vspace{3pt}
On the elbow specifically: fine as a first look, terrible as a final answer. The curve is usually smooth, so the kink you see depends on the plot's aspect ratio; two people read the same chart and pick different k. I would ask him to overlay other methods before defending eight.\par\vspace{3pt}
I would run a portfolio. Silhouette, which has a genuine optimum in k and gives per-point diagnostics, though it favours compact convex clusters and is quadratic in distances, so subsample. The gap statistic, which compares log inertia against a uniform reference and is the only common method that can return k equals one - no cluster structure at all - which is its real value. BIC on a GMM if a probabilistic model is acceptable. And stability: recluster bootstrap or split-half resamples and measure label agreement with adjusted Rand. Partitions that do not reproduce are not real, and that is the most persuasive evidence.\par\vspace{3pt}
Then the actual move: decide by the downstream use. If these are marketing segments a team has to staff and write copy for, k is bounded by headcount - eight might be right because it is operable, not because inertia bent.\par\vspace{3pt}
\flab{Loses the signal}
\begin{itemize}
\item Naming the elbow and stopping, with no acknowledgement that it is subjective.
\item Trying to choose k by minimising inertia - it is monotone, which is the tell.
\item Never asking what the clusters are actually for.
\end{itemize}
\lab{Then they ask}
\begin{itemize}
\item What does the gap statistic compare against?
\item Silhouette says 2, BIC says 11 - now what?
\item How would you convince a PM these clusters are real?
\end{itemize}
\vspace{2pt}\noindent{\color{muted}\rule{\linewidth}{0.3pt}}
\needspace{6\baselineskip}
\qhead{L6}{Debugging}{Your t-SNE plot of user embeddings shows five clean, well-separated clusters. What can you conclude?}
\lab{Say this}
Almost nothing yet, and saying that immediately is the whole test. What a clean picture licenses is a hypothesis: there may be local neighbourhood structure in the embedding space. It does not tell me there are five clusters.\par\vspace{3pt}
Three specific reasons. t-SNE optimizes a local-neighbourhood objective with a heavy-tailed kernel in the low-dimensional space, which actively creates gaps --- it exists to solve the crowding problem. Run it on i.i.d. Gaussian noise at low perplexity and you'll get crisp islands. Cluster sizes on the plot are meaningless too, because the per-point bandwidth is tuned to hit a fixed perplexity, expanding sparse regions and contracting dense ones. And distances between clusters carry essentially no information: the objective is asymmetric, so pushing dissimilar points far apart is nearly free, which makes reading distances off the plot the classic error.\par\vspace{3pt}
What I'd do instead. Re-run at perplexity 5, 30 and 100, with PCA initialization and a few seeds, and keep only what survives all of them. Then test the hypothesis in the original space: take the visual groups as tentative labels and train a simple classifier --- if held-out accuracy is near chance, the groups don't exist. Cluster properly in the original space and check agreement by adjusted Rand. And colour by the obvious confounds --- signup cohort, country, device, tenant, model version --- because five beautiful islands usually turn out to be five values of a categorical feature.\par\vspace{3pt}
One non-negotiable: never fit a downstream model on the two-dimensional coordinates.\par\vspace{3pt}
\flab{Loses the signal}
\begin{itemize}
\item ``We have five user segments'' as a conclusion
\item Reading inter-cluster distances or cluster sizes as meaningful
\item Using t-SNE coordinates as features, or having no plan to validate in the original space
\end{itemize}
\lab{Then they ask}
\begin{itemize}
\item Would UMAP fix any of this?
\item How would you actually cluster these embeddings?
\item What would convince you the five groups are real?
\end{itemize}
\vspace{2pt}\noindent{\color{muted}\rule{\linewidth}{0.3pt}}
\needspace{6\baselineskip}
\qhead{L5}{Conceptual}{When does k-means fail, and what do you use instead?}
\lab{Say this}
The clean way to answer this is from the model rather than a list. Minimizing within-cluster sum of squares is maximum likelihood for a mixture of isotropic, equal-variance, equal-weight Gaussians under hard assignment. Every failure is one of those assumptions breaking, so you can derive the list on the spot instead of memorizing it.\par\vspace{3pt}
Non-convex shapes --- rings, moons --- break because centroid assignment gives linear Voronoi boundaries, and you can't carve a ring with straight lines; that's DBSCAN or spectral clustering. Elliptical or unequally spread clusters break isotropy and equal variance; that's a GMM with full covariance. Very unequal cluster sizes break equal weight --- the objective would rather split a big mass and swallow a small one --- so free the mixing weights. Outliers break it because centroids are means and means aren't robust; k-medoids, or DBSCAN if you want them labeled as noise. High dimension breaks it through distance concentration; PCA first, or spherical k-means on normalized vectors.\par\vspace{3pt}
Two more that aren't about the model. Unscaled features: whichever column has the widest range owns the objective --- dollars spent versus visits per month is a millionfold mismatch in squared distance, so standardize, always. And bad initialization, the most common practical failure of all; k-means++ with restarts.\par\vspace{3pt}
That said, k-means is still the right first move on most tabular problems: fast, the baseline everyone compares against, and its failures are visible if you plot a PCA projection colored by cluster.\par\vspace{3pt}
\flab{Loses the signal}
\begin{itemize}
\item An unstructured list of failures with no unifying reason, and no alternative named for each.
\item Not mentioning feature scaling --- the single most common real-world cause of nonsense clusters.
\item Suggesting ``more clusters'' as the fix for non-convex shapes; that patches a ring with blobs and destroys interpretability.
\end{itemize}
\lab{Then they ask}
\begin{itemize}
\item Why exactly are k-means boundaries linear?
\item Would scaling fix the rings problem?
\item How do you pick k, and how do you know the partition is real?
\end{itemize}
\vspace{2pt}\noindent{\color{muted}\rule{\linewidth}{0.3pt}}
\needspace{6\baselineskip}
\qhead{L5}{Debugging}{Training RMSE 0.31, validation RMSE 0.94, and the two learning curves are still far apart at 200K rows. Diagnose it and name the knobs, on the model of your choice.}
\lab{Say this}
That's the variance signature, and the most useful thing on that plot isn't the gap - it's that the curves still haven't met at 200 thousand rows, which tells you more data will help. That's the decision-relevant fact.\par\vspace{3pt}
Before I touch a hyperparameter, though, I'd audit the split, because a gap that large has two non-model explanations that no amount of regularization fixes. Leakage or a broken split: a temporal problem split randomly, group leakage where the same user or merchant lands on both sides, or a feature derived from the target. Compare a random split against a grouped or temporal one - if the gap moves, that's your answer. And distribution shift: if validation is a later period, you may be looking at drift rather than variance, and regularization won't touch drift.\par\vspace{3pt}
Once that's clear, the knobs depend on the model. GBDT: early stopping on rounds first, because boosting overfits in the number of rounds by construction, then a lower learning rate with more rounds, shallower trees, higher \texttt{min\_child\_weight}, row and column subsampling, more L2. Random forest: this pattern is unusual, since averaging already kills variance, so suspect leakage before tuning, then raise \texttt{min\_samples\_leaf} and lower \texttt{max\_features}. Ridge or lasso: raise lambda. kNN: raise k, whose variance is exactly sigma squared over k - the cleanest knob in classical ML.\par\vspace{3pt}
And across all of them the same three levers: more rows, fewer features, more averaging.\par\vspace{3pt}
\lab{Numbers}
\begin{itemize}
\item kNN variance is sigma\textasciicircum{}2 / k, so k is the cleanest variance knob available
\end{itemize}
\flab{Loses the signal}
\begin{itemize}
\item Reads a train-validation gap this size as high bias and proposes adding capacity.
\item Jumps straight to hyperparameters without questioning the split.
\item Recites the bias-variance decomposition but can't say the expectation is over the draw of the training set.
\end{itemize}
\lab{Then they ask}
\begin{itemize}
\item How would you estimate the irreducible noise?
\item The curves have already met - now what?
\item Does any of this change for a 7B-parameter network?
\end{itemize}
\vspace{2pt}\noindent{\color{muted}\rule{\linewidth}{0.3pt}}
\chapter{Applied Craft}
\addtocontents{toc}{\protect\thispagestyle{fancy}}
\needspace{6\baselineskip}
\qhead{L6}{Debugging}{Your churn model has 0.99 AUC offline and 0.62 in production. Walk me through it.}
\lab{Say this}
The magnitude tells you which branch to chase. Two families explain an offline-online gap: leakage, meaning the offline number was never real, and skew or drift, meaning it was real and then something changed. A 0.99 churn AUC is not plausible under any circumstances --- churn models live at 0.75 to 0.85 --- so this is leakage until proven otherwise, and I wouldn't spend the first hour on drift.\par\vspace{3pt}
Then the playbook, cheapest first. Look at feature importance as an alarm, then train a model on the top feature alone; if one column reproduces most of the 0.99, that's the whole story. Ablate it and see where the metric lands --- if it falls to 0.81, that's your honest baseline.\par\vspace{3pt}
The step I'd actually insist on is a point-in-time audit: one row per feature giving the value's timestamp, the prediction timestamp, and whether it's retrievable at serving. For churn the classic offenders are an account-closed date populated only by the event you're predicting, total charges that stop accumulating at departure, a support-ticket count that includes the cancellation ticket, and any current-value dimension attribute joined onto historical rows --- the warehouse stores the present, so joining today's plan tier onto an eleven-month-old event fabricates knowledge.\par\vspace{3pt}
Then split integrity: temporal or shuffled, users repeated across folds, duplicate rows. The confirming experiment is rebuilding with a strict as-of cutoff and backtesting on the following period. If that comes back at 0.64 and matches production, the diagnosis is closed.\par\vspace{3pt}
\lab{Numbers}
\begin{itemize}
\item A realistic churn AUC is \textasciitilde{}0.75--0.85; 0.99 is a leakage signal, not a win
\end{itemize}
\flab{Loses the signal}
\begin{itemize}
\item ``Retrain on fresh data'' --- a fix before a diagnosis, and useless if the training set is contaminated
\item Treating 0.99 as plausible; no domain prior for what a churn AUC looks like
\item Naming leakage but having no procedure --- no ablation, no timestamp audit, no backtest
\end{itemize}
\lab{Then they ask}
\begin{itemize}
\item Which churn features would you audit first, and why?
\item The audit is clean and the gap is still 20 points --- now what?
\item How would you have caught this before launch?
\end{itemize}
\vspace{2pt}\noindent{\color{muted}\rule{\linewidth}{0.3pt}}
\needspace{6\baselineskip}
\qhead{L6}{Mathematical}{You want to target-encode a 50K-cardinality \texttt{merchant\_id}. Give me the exact construction, and tell me where it leaks if you get it wrong.}
\lab{Say this}
Start from the leak, because it motivates the construction. The naive version replaces each merchant with the mean target over its training rows --- including the row you're encoding. For a merchant with one transaction, the encoded value is that row's label. Trees find it instantly, training metrics look superb, and it collapses in validation.\par\vspace{3pt}
Two ingredients fix it. Shrinkage: encode with the empirical-Bayes posterior mean --- n times the merchant's own rate plus m times the global rate, over n plus m --- so a merchant with three thousand transactions is trusted and one with two is pulled to the global rate. And out-of-fold: split training into K inner folds, encoding each fold's rows from statistics on the others, so no row's own label reaches its own feature.\par\vspace{3pt}
Then the rules that make it survive a real pipeline. The encoder you ship is fit on the whole training set and applied to validation, test and serving --- those rows contributed no labels, so that isn't leakage. The whole procedure lives inside the outer CV fold; encoding before you split is the worst contamination there is, because the transform carries labels. The inner split must match the outer one's structure --- group-aware if that's group-aware, time-aware if temporal --- or you've moved the leak, not removed it. And unseen merchants fall back to the prior by construction, which matters at 50K cardinality because new ones appear daily.\par\vspace{3pt}
If you'd rather not own that apparatus, CatBoost does ordered target statistics natively.\par\vspace{3pt}
\lab{Numbers}
\begin{itemize}
\item Smoothing prior weight m: typically 5-100
\end{itemize}
\flab{Loses the signal}
\begin{itemize}
\item Computing the mean over the full training set and calling it done --- the definitional failure of this question.
\item Smoothing but not going out-of-fold, or out-of-fold but not smoothing; the singleton-merchant problem needs both.
\item Fitting the encoder before the train/test split --- ``I encode, then I cross-validate.''
\end{itemize}
\lab{Then they ask}
\begin{itemize}
\item What exactly does m trade off, and how would you choose it?
\item Your outer split is by month --- what does the inner split look like?
\item Same feature for a logistic regression: what changes?
\end{itemize}
\vspace{2pt}\noindent{\color{muted}\rule{\linewidth}{0.3pt}}
\needspace{6\baselineskip}
\qhead{L5}{Debugging}{Here is a pipeline. \texttt{StandardScaler} on the full dataframe, then median imputation, then \texttt{SelectKBest(k=100)} on all rows, then SMOTE to balance, then a 5-fold CV that reports 0.94 AUC. Find the bugs and rank them by damage.}
\lab{Say this}
Every stateful step is fit on data that includes the evaluation rows, so all four are train-test contamination --- but they aren't equally bad, and ranking them is the point.\par\vspace{3pt}
Worst is SelectKBest on all rows. Feature selection sees the labels of the rows it will later be scored on, so the held-out folds helped choose the features that predict them. It can manufacture a strong metric from pure noise: in the ESL demonstration, screening five thousand random predictors to the hundred most correlated on fifty samples cross-validates to three percent error where the truth is fifty. With k of a hundred, that's most of the 0.94.\par\vspace{3pt}
Second, SMOTE before the split: synthetic minority points interpolate training rows and land in the validation folds, so the model is graded on near-copies of what it memorized. Third, imputation: a global median is mild, but an iterative or kNN imputer fits a model per column across the whole dataset. Mildest is the scaler: a mean and standard deviation over a few extra percent of rows barely moves. Still wrong, still fix it, but it doesn't explain the number.\par\vspace{3pt}
The fix is structural: put every step in a pipeline object and cross-validate the pipeline, so each fold refits all of them on its own training portion, resampling training rows only. Then rerun and expect a lower, honest number. I'd also ask whether the split should be temporal or grouped, and whether SMOTE is needed at all versus class weights and a tuned threshold.\par\vspace{3pt}
\lab{Numbers}
\begin{itemize}
\item ESL \S{}7.10.2: 50 samples, 5,000 random predictors, screen top 100 $\to$ 3\% reported error vs 50\% truth
\end{itemize}
\flab{Loses the signal}
\begin{itemize}
\item ``Looks fine'' --- the single most damaging answer available in an applied loop.
\item Spotting the scaler (the famous example) but missing feature selection (the damaging one), or not noticing that SMOTE must never touch validation rows.
\item Fixing it by recomputing statistics per fold by hand instead of using a pipeline object.
\end{itemize}
\lab{Then they ask}
\begin{itemize}
\item How much of the 0.94 do you expect to survive?
\item Which of these is still wrong even with a single train/test split?
\item Where would you put early stopping in the fixed pipeline?
\end{itemize}
\vspace{2pt}\noindent{\color{muted}\rule{\linewidth}{0.3pt}}
\needspace{6\baselineskip}
\qhead{L6}{Trade-off}{Fraud is 0.1\% of transactions. Walk me through building the classifier.}
\lab{Say this}
I would deliberately not start with the imbalance. Two questions first. What is the absolute positive count? A tenth of a percent of 200 million transactions is 200,000 frauds, which is a large dataset with a skewed ratio and a perfectly normal supervised problem. A tenth of a percent of 20,000 rows is twenty frauds, which is a fundamentally different problem where I would be arguing for rules plus anomaly detection. Second, what decision does the score drive - auto-decline, step-up authentication, or a review queue with fixed daily capacity? That determines the metric.\par\vspace{3pt}
Then, in order. Metric: not accuracy, not raw ROC-AUC. PR-AUC reported against its 0.001 baseline, plus the operational metric - precision in the top k alerts per day, or recall at a fixed false-positive rate. The arithmetic that justifies that: at eighty percent TPR and one percent FPR, a million transactions gives 800 true alerts and about 10,000 false ones, so precision is seven percent while the ROC point looks respectable.\par\vspace{3pt}
Validation: strictly temporal, because fraud is adversarial, and grouped by card or account. Model: a GBDT on transaction, velocity and entity-history features, with every aggregate anchored point-in-time, because velocity features are the richest leakage surface here. Threshold: from the cost matrix - five dollars of review against a two-hundred-dollar chargeback is about 2.4 percent - or from review capacity. Recalibrate if you weighted anything. And only then resampling, mostly for training-loop economics rather than accuracy.\par\vspace{3pt}
\lab{Numbers}
\begin{itemize}
\item 0.1\% of 200M = 200K positives; 0.1\% of 20K = 20 positives - different problems
\item TPR 0.8 / FPR 1\% on 1M transactions: 800 true vs \textasciitilde{}10K false alerts, precision \textasciitilde{}7\%
\item Threshold from costs: $5 review vs $200 chargeback gives tau* \textasciitilde{} 2.4\%
\end{itemize}
\flab{Loses the signal}
\begin{itemize}
\item Opening with SMOTE, or any resampling, before naming a metric.
\item Reporting accuracy, or ROC-AUC alone, at this prevalence.
\item A random train/test split on adversarial time-series data, or leaving the threshold at 0.5.
\end{itemize}
\lab{Then they ask}
\begin{itemize}
\item The chargeback label takes 60 days - how do you evaluate last month?
\item Review capacity is halved. What changes?
\item How do you know the model is not just re-learning the rule engine?
\end{itemize}
\vspace{2pt}\noindent{\color{muted}\rule{\linewidth}{0.3pt}}
\needspace{6\baselineskip}
\qhead{L5}{Debugging}{How would you debug poor embedding quality?}
\lab{Say this}
I would go cheapest diagnostic first and only retrain once I know what is broken.\par\vspace{3pt}
First, look at them: a UMAP projection, not t-SNE, which does not preserve global structure, coloured by whatever labels I have. If there is no visible structure at all, the embeddings are fundamentally broken and I stop there. Then I make that quantitative with average intra-class cosine similarity against inter-class. What I care about is the gap, not the absolute numbers, because raw language-model features are anisotropic and everything looks similar, which people misdiagnose as collapse. Both high with no gap is collapse. Intra-class low means the model never learned category structure, and I would go look at the training data.\par\vspace{3pt}
Then a nearest-neighbour spot check on fifty to a hundred diverse queries: pull the top ten and actually read them. That is where you catch systematic failures, like neighbours that always share surface keywords but not meaning.\par\vspace{3pt}
If it is still not obvious, I would run an SVD on a sample of the embedding matrix - if the top few percent of dimensions carry more than ninety percent of the variance, that is dimensional collapse. And I would test on a proxy task, k-NN classification on frozen embeddings, because if the proxy is fine and the real task is not, the problem is downstream. Last, training diagnostics: is the loss still moving, are gradient norms sane, and for contrastive setups, the temperature.\par\vspace{3pt}
\lab{Numbers}
\begin{itemize}
\item Contrastive temperature: unstable below \textasciitilde{}0.01, too easy above \textasciitilde{}0.5
\end{itemize}
\flab{Loses the signal}
\begin{itemize}
\item Jumping to ``retrain on more data'' before diagnosing anything.
\item Relying on t-SNE alone, or on ``the plot looks fine'' with no numbers.
\item Calling uniformly high cosine similarity collapse without checking the intra-versus-inter gap.
\end{itemize}
\lab{Then they ask}
\begin{itemize}
\item What does collapse look like in a UMAP?
\item How do you tell bad embeddings from a bad evaluation metric?
\item How would you monitor embedding quality automatically in production?
\end{itemize}
\vspace{2pt}\noindent{\color{muted}\rule{\linewidth}{0.3pt}}
\needspace{6\baselineskip}
\qhead{L6}{Debugging}{You fine-tuned a model and it performs well on your test set but poorly in production. What happened?}
\lab{Say this}
Nine times out of ten the test set simply isn't production, and I'd check that before anything else.\par\vspace{3pt}
Distribution shift is hypothesis one. Curated test data is clean; production has typos, slang, truncated inputs, formats nobody anticipated, and users actively trying to break it. The check is cheap --- pull a stratified sample of real production traffic, label a few hundred, and score the model on that. If the gap reproduces there, you've found it and you can stop.\par\vspace{3pt}
Second, catastrophic forgetting. Aggressive fine-tuning overfits the model to a narrow slice and costs it the general capability it needs for everything outside that slice. Compare base against fine-tuned on a broad benchmark, not just your task. Fixes are lower learning rate, fewer epochs, LoRA rather than full fine-tuning, or mixing general-purpose data back into the fine-tuning mixture.\par\vspace{3pt}
Third, leakage. If your test set overlaps the training data --- near-duplicates, or just the same time window --- the offline number was fiction. Split temporally and run near-duplicate detection across the two sets.\par\vspace{3pt}
Fourth, metric mismatch. Accuracy or perplexity on clean inputs may just not correlate with what production cares about --- task completion, escalation rate, user satisfaction. That one is the most common and the least often caught, because everyone's dashboard keeps saying the model is fine.\par\vspace{3pt}
The durable fix for all four is the same: build the eval set out of production logs, and keep refreshing it.\par\vspace{3pt}
\flab{Loses the signal}
\begin{itemize}
\item Jumping to ``retrain on more data'' before diagnosing which failure it is.
\item Not treating distribution shift as the first hypothesis.
\item Blaming model capacity with no evidence.
\end{itemize}
\lab{Then they ask}
\begin{itemize}
\item How do you prevent catastrophic forgetting?
\item How do you build a test set that actually represents production?
\item What monitoring would have caught this earlier?
\end{itemize}
\vspace{2pt}\noindent{\color{muted}\rule{\linewidth}{0.3pt}}
\needspace{6\baselineskip}
\qhead{L6}{Debugging}{Your model's CTR dropped 5\% overnight. Walk through your diagnosis.}
\lab{Say this}
I triage in order of prior probability and speed of check, and ``retrain the model'' is last, not first. Overnight is the key word --- models don't drift five percent overnight, systems do. First, was there a deploy? New model, config change, feature-store migration, new hosts, load balancer. If a deploy lines up with the drop, roll back first and diagnose after. Then serving health: error rates, latency, timeouts. A latency spike that trips fallbacks to a non-personalized default looks exactly like a CTR drop.\par\vspace{3pt}
Second, the data pipeline. Check feature freshness --- a batch job that failed quietly leaves the model scoring on stale features. Check for nulls from an upstream schema change, and PSI on the top features to see whether the input distributions actually moved.\par\vspace{3pt}
Third, external causes. A frontend change that moved the module. A traffic-mix shift from a marketing push or a viral event, where the population changed rather than the model. Calendar effects.\par\vspace{3pt}
Only then the model, and only if it retrained overnight: compare offline metrics against the previous version and look at the prediction distribution. A shifted output distribution with no corresponding input shift points at a training bug. Running alongside all of that, confirm the metric is real. A logging change can manufacture a five percent drop that never happened, and that's an embarrassing thing to find on day three.\par\vspace{3pt}
\flab{Loses the signal}
\begin{itemize}
\item Starting with ``retrain the model'' before checking system health and data pipelines.
\item Never asking whether a deployment happened.
\item Not separating a measurement artifact from a real drop.
\end{itemize}
\lab{Then they ask}
\begin{itemize}
\item How do you tell a real drop from a measurement artifact?
\item What if the drop is confined to one user segment?
\item How would you build alerting that catches this in 30 minutes instead of overnight?
\end{itemize}
\vspace{2pt}\noindent{\color{muted}\rule{\linewidth}{0.3pt}}
\needspace{6\baselineskip}
\qhead{L5}{Conceptual}{How do you handle training-serving skew?}
\lab{Say this}
Training-serving skew is when the model sees different feature distributions at serving time than it did in training, and it's the single most common cause of ``works offline, fails in production.'' I handle it as prevention, detection and mitigation --- and prevention is where the leverage is.\par\vspace{3pt}
Prevention means removing the opportunity for two implementations to disagree. A feature store with shared definitions, so each feature is written once and used by both paths. The log-and-train pattern: log the exact feature vector you served and train on those logs rather than recomputing from raw sources --- the strongest guarantee available, because the training inputs are literally the serving inputs. And package preprocessing into the model artifact, so normalization and tokenization can't drift away from the model that expects them.\par\vspace{3pt}
Detection: alert on PSI above about 0.1 comparing serving to training feature distributions; monitor the prediction distribution against offline evaluation, which catches skew you never thought to instrument per-feature; and put an integration test in CI that pushes fixed inputs through both pipelines and requires matching outputs.\par\vspace{3pt}
Mitigation depends on cause. If a pipeline is late, fall back to last-known-good values rather than nulls the model has never seen. If the shift is real and permanent, retrain on data that reflects it. And keep this separate from concept drift: skew is an engineering mismatch, drift is the world moving.\par\vspace{3pt}
\lab{Numbers}
\begin{itemize}
\item Alert threshold: PSI > 0.1 comparing serving to training feature distributions
\end{itemize}
\flab{Loses the signal}
\begin{itemize}
\item ``Be careful with your features'' with no concrete mechanism.
\item Not knowing the log-and-train pattern.
\item Confusing training-serving skew with concept drift.
\end{itemize}
\lab{Then they ask}
\begin{itemize}
\item How do you detect skew in embedding features?
\item What if the skew is intentional, from an upstream product change?
\item How does log-and-train handle delayed labels?
\end{itemize}
\vspace{2pt}\noindent{\color{muted}\rule{\linewidth}{0.3pt}}
\end{document}